% =====================================================================
%  Harald Høffding, FILOSOFISKE PROBLEMER.  Indbydelsesskrift til
%  Kjøbenhavns Universitets Aarsfest til Erindring om Kirkens Reformation,
%  November 1902.  Kjøbenhavn: Universitetsbogtrykkeriet (J. H. Schultz),
%  1902.  90 pp.
%  TRANSCRIPTION (Danish) of the 1902 print, complete.
%
%  SOURCE: the Royal Library's (KB) scan, ~/bibliotek/Høffding, Harald/
%  1902-filosofiske-problemer.pdf (SCANS.tsv).  PAGE MAP (pagemap.py):
%  printed page = PDF page - 13; p. 1 = PDF 14 ... p. 90 = PDF 103.  Front
%  matter (unnumbered): the Indbydelsesskrift leaf PDF 4 (the same leaf
%  again at PDF 6), title PDF 8, preface PDF 10, contents PDF 12; these
%  carry "% --- front matter ---" comments and no \opage.  Not
%  transcribed: KB's cover sheet (PDF 1), blank and library leaves, the
%  bound-in university list after p. 90 (PDF 104-123).  Every page turn
%  is marked "% --- p. N ---" + \opage{N} at the first word of the printed
%  page; a word divided over a page turn ends in %.
%
%  CONVENTIONS: italic -> \textit; small capitals -> \textsc (normal case
%  inside); printed dashes -> ---, number ranges -> --; quotation marks
%  „…“ as printed; subscripts as printed (pp. 50-51).  Høffding's notes
%  are endnotes („Noter“, pp. 85-90, in small type), each headed by its
%  label and the page it belongs to; in the text the anchors are a
%  superscript figure and a parenthesis („fremgaaet“²⁸)“).  The labels run
%  1-12, 12b, 13-31, 31b, 31c, 32-60, as printed; in the Noter the 12b note
%  is labelled „12)“ (marked PRINTED AS IS).  Not transcribed: folios,
%  sheet signatures („12“, p. 89), library stamps.
%
%  METHOD (the e-book method, 2026-10-02; TRANSCRIPTION-PLAYBOOK §00 point
%  0).  No one typed the text.  The base text -- words, punctuation,
%  paragraphs, italics, small caps, headings, note texts -- is the SAGA
%  e-book (Filosofiske_problemer.epub, in this folder); it was aligned
%  letter by letter against the scan's own text layer by
%  ebook-method/filosofiske-problemer/build.py (tracked, repo root), which
%  also takes every page turn from the layer.
%  Of the letter and figure disagreements, the script settled by rule
%  those that are the layer's known faults (glyph confusions, letterspaced
%  small capitals, dropped accents, specks, the note labels); every other
%  one was decided at the image from crop sheets (sheets.py, same folder):
%  52 in the text, 62 in the notes, 22 punctuation disagreements in the
%  text and 30 in the notes (among them every comma closing an italic
%  title and every note's last mark), and 4 paragraph disagreements
%  (decisions.tsv, patches.tsv).  Spot pages 9 and 87 read in full
%  against the image afterwards; p. 87 led to the rule that a layer glyph
%  fault settles a disagreement only where the e-book's word is attested
%  elsewhere in the repository.  The print corrects the
%  e-book at: p. 2 „dér“, „énhver“ (accents); p. 14 Aarsagsforholdet;
%  p. 17 Flechsigs; p. 22 Sammenhøren; p. 23 saadant; p. 41 Genstandes;
%  p. 42 Molekuler; p. 46 a stray point removed; pp. 50-51 the
%  subscripts; p. 54 fører; p. 64 hele, Verdensprincipet; p. 65 føres;
%  p. 68 Resultat; p. 70 har; p. 71 Existens; p. 72 Ret, Principet;
%  p. 77 Friedrich, and „som Selvhævdelsen“, which the e-book omits; in
%  the Noter: Der, Roberto, apparet, p. 36, er, Bevidsthedsfænomenerne,
%  Psykologiske, experience, nur, nicht, Wikner, her, Grundlag,
%  sufficiunt, Wissenschaftslehre, „p. 32 f.“, the point (not comma)
%  after nine italic titles, the Greek of note 60, the note labels.
%  PRINTED AS IS: p. 47 „1dentitetsforhold“ (the figure 1 for I);
%  „(S 3)“, „(S 8)“, „(S 46)“ in the Noter; O with
%  subscript s for O-zero, pp. 50 and 51; the second „12)“.  One doubtful
%  point (p. 89, after „Hegel“) is marked % UNSURE.
%
%  RESIDUAL RISKS: (1) a misprint that the e-book and the layer both
%  normalise is invisible to the method; (2) italics and small capitals
%  are the e-book's, which the layer cannot witness; (3) punctuation is
%  witnessed only where the layer reads it, and quotation marks and
%  dashes come from the e-book alone; (4) letterspacing, if any, is not
%  recorded.
% =====================================================================
\documentclass[12pt,a4paper,openany]{book}
\usepackage[utf8]{inputenc}
\usepackage[T1]{fontenc}
\usepackage{textalpha}
\usepackage{libertinus}
\usepackage{libertinust1math}
\usepackage{graphicx}    % for \reflectbox (the old Danish ɔ: id-est mark)
\DeclareUnicodeCharacter{0254}{\reflectbox{c}}  % ɔ (U+0254) = reversed c
\usepackage[danish]{babel}
\usepackage[protrusion=true,expansion=true]{microtype}
\usepackage{setspace}
\usepackage[margin=1.2in]{geometry}
\usepackage{fancyhdr}
\setlength{\headheight}{28pt}
\addtolength{\topmargin}{-13.5pt}
\usepackage{hyperref}
\hypersetup{pdftitle={Filosofiske Problemer}, pdfauthor={Harald Høffding}, hidelinks}
\pagestyle{fancy}
\fancyhf{}
\fancyhead[LE,RO]{\thepage}
\fancyhead[RE]{\textit{Harald Høffding: Filosofiske Problemer}}
\fancyhead[LO]{\textit{1902}}
\setstretch{1.3}
\title{\textbf{Filosofiske Problemer}}
\author{Harald Høffding}
\date{Kjøbenhavn: Universitetsbogtrykkeriet (J. H. Schultz), 1902}
\usepackage{marginnote}
\newcommand{\opage}[1]{\marginnote{\footnotesize\textit{[#1]}}}
% Contents lines (front matter).
\newcommand{\tocline}[3]{\noindent\makebox[2.2em][l]{#1}#2\dotfill\ #3\par}
\newcommand{\tocsub}[2]{\noindent\hangindent=4.4em\hspace*{2.2em}#1\dotfill\ #2\par}
\newcommand{\tocplain}[1]{\noindent\hspace*{2.2em}#1\par}

\begin{document}
\maketitle
\thispagestyle{empty}
\clearpage

% --- front matter: Indbydelsesskrift (PDF 4; the same leaf again at PDF 6) ---
\begin{center}
INDBYDELSESSKRIFT\\[2pt]
{\scriptsize TIL}\\[2pt]
{\large KJØBENHAVNS UNIVERSITETS}\\
{\large AARSFEST}\\[2pt]
{\scriptsize TIL ERINDRING OM}\\[2pt]
{\large KIRKENS REFORMATION}\\[2pt]
NOVEMBER 1902\\[2em]
% (the University's seal)
HARALD HØFFDING:\\
FILOSOFISKE PROBLEMER\\[1em]
\rule{1cm}{0.4pt}\\[1em]
{\small KJØBENHAVN}\\
{\scriptsize TRYKT I UNIVERSITETSBOGTRYKKERIET (J. H. SCHULTZ)}\\
{\scriptsize 1902}
\end{center}
\clearpage

% --- front matter: title (PDF 8) ---
\begin{center}
{\Huge FILOSOFISKE}\\[10pt]
{\Huge PROBLEMER}\\[1.5em]
{\scriptsize AF}\\[1.5em]
{\large HARALD HØFFDING}\\[3em]
\rule{1.5cm}{0.4pt}\\[3em]
KJØBENHAVN\\
{\scriptsize TRYKT I UNIVERSITETSBOGTRYKKERIET (J. H. SCHULTZ)}\\
{\scriptsize 1902}
\end{center}
\clearpage

% --- front matter: preface, no heading, in italic (PDF 10) ---
\begin{itshape}
Det første Udkast til denne Afhandling laa til Grund for en Række
Forelæsninger, jeg i Februar Maaned d. A. holdt ved Upsala Universitet efter
Indbydelse fra dettes humanistiske Sektion. Efterat jeg nu har taget de Emner,
jeg behandlede i disse Forelæsninger, op til fornyet Drøftelse, vender min Tanke
sig med hjertelig Hilsen og Tak til den minderige Højskole ved Fyrisaa, som for
en Tid vilde regne mig blandt sine Lærere, --- og til de ældre og yngre Venner,
jeg der vandt.

\bigskip
16. August 1902.

\begin{flushright}HARALD HØFFDING.\end{flushright}
\end{itshape}
\clearpage

% --- front matter: contents (PDF 12) ---
\begin{center}{\large INDHOLD}\end{center}
\begin{flushright}{\scriptsize Side}\end{flushright}
\tocline{}{INDLEDNING}{1}
\tocline{I.}{BEVIDSTHEDSPROBLEMET}{6}
\tocsub{1. Personlighedsbegrebet og den analytiske Psykologi}{6}
\tocsub{2. Diskontinuiteter paa det psykiske Omraade}{12}
\tocsub{3. Psykologi og Fysiologi}{22}
\tocsub{4. Villie og Energi}{25}
\tocline{II.}{ERKENDELSESPROBLEMET}{28}
\tocsub{1. Erkendelsens Arter}{28}
\tocsub{2. Erkendelsens Principer}{31}
\tocsub{3. Kvalitet og Kvantitet}{39}
\tocsub{4. Elementært og idealt Aarsagsbegreb}{43}
\tocsub{5. Subjekt og Objekt}{48}
\tocline{III.}{TILVÆRELSESPROBLEMET}{53}
\tocsub{1. Problem og Metode}{53}
\tocsub{2. Metafysik som Kunst}{57}
\tocsub{3. Første Urfænomen (Enhed eller Mangfoldighed)}{59}
\tocsub{4. Andet Urfænomen (Aand eller Materie)}{62}
\tocsub{5. Tredie Urfænomen (Bestaaen eller Udvikling)}{65}
\tocline{IV.}{VURDERINGSPROBLEMET}{70}
\tocsub{1. Indledning}{70}
\tocplain{\hspace*{1em}\textit{A. Det etiske Problem.}}
\tocplain{\hspace*{3em}a. Det etiske Arbejde.}
\tocsub{2. Kontinuitetsprincipet i Etiken}{72}
\tocplain{\hspace*{3em}b. Den etiske Vurderings Rationalitet.}
\tocsub{3. Grundværdiernes Kamp}{75}
\tocsub{4. De individuelle Betingelsers Forskellighed}{77}
\tocplain{\hspace*{1em}\textit{B. Det religiøse Problem.}}
\tocsub{5. Kontinuitetsprincipet i Religionsfilosofien}{78}
\tocsub{6. Religionens psykologiske Sted}{80}
\tocsub{7. Historiske Former for Religion}{82}

\medskip
\noindent Noter S. 85.
\begin{center}\rule{2cm}{0.4pt}\end{center}
\clearpage

% ===== TEXT (pp. 1--84) =====
% --- p. 1 ---
\opage{1}

\begin{center}\textbf{\large Indledning.}\end{center}

Filosofiske Ideer kunne, som Filosofiens Historie viser, have en dobbelt
Betydning. De kunne være Forsøg paa Opstilling, Behandling eller Løsning af
visse Problemer, og de kunne staa som Symptomer paa visse Tendenser i det
menneskelige Aandsliv. Der vil være en stadig Vexelvirkning mellem disse to
Sider ved Filosofien, idet det kan staa som Symptom paa en ejendommelig
aandelig Bevægelse, at et Problem opstilles og behandles paa en bestemt Maade,
og idet paa den anden Side Problemernes Brod kan udløse aandelige Bevægelser,
der ellers ikke vilde blive til. Der ytrer sig i denne Vexelvirkning en indre
Sammenhæng mellem personligt Liv og videnskabelig Forsken, som mere eller
mindre gør sig gældende paa alle videnskabelige Omraader, men som efter Sagens
Natur ganske særligt vil træde frem, naar vi nærme os Videnskabernes
Grænseomraader. Den træder mere frem i Aandsvidenskaben end i Naturvidenskaben,
og den træder allermest frem i Filosofien, hvis Problemer væsentligt ere
Grænseproblemer.

Det er ganske vist en ikke sjelden Antagelse, at der er en Modstrid mellem
Personlighed og videnskabelig Forsken, saa at man helst bør afføre sig sin
Personlighed, naar man vil tænke videnskabeligt, og at paa den anden Side
videnskabelige Metoder og Resultater ikke betyde Noget for det personlige Liv i
dettes frie Udvikling. Jeg har selv i min Ungdom hyldet en saadan Antagelse
under den overvældende Indflydelse af S. Kierkegaards
% --- p. 2 ---
\opage{2}Ideer. Naar jeg mere og mere er kommen bort fra denne Antagelse, tør
jeg sige, at det er, fordi jeg har lært baade Videnskaben og det personlige Liv
bedre at kende.

Trangen til videnskabelig Forsken er en speciel Form for Trangen til
Overensstemmelse med sig selv under alle mangfoldige og skiftende Erfaringer.
Den ytrer sig baade i den formale Videnskab og i den reale, --- baade som Trang
til at forme Tankerækker, i hvilke det ene Led med indre Nødvendighed udvikler
sig af det andet, og som Trang til at forbinde foreliggende Erfaringer
saaledes, at der opnaas saa omfattende og saa fast en Kontinuitet som muligt.
Personlighed bestaar først og fremmest i en indre Enhed og Sammenhæng mellem
alle Forestillinger, Følelser og Bestræbelser. Det er da intet Offer,
Personligheden bringer, naar den sætter Livet ind paa at forske. Den
frembringer i Videnskaben --- som i Kunsten --- et objektivt Genbillede af sig
selv. Kun derved forstaas den inderlige Karakter, Erkendelsesglæden kan antage,
den amor intellectualis, som Forskeren føler ved at arbejde sig frem til et
Synspunkt, hvorfra det Enkelte kan ses som Led i en stor Sammenhæng. Den
hellige Ild, som stedse paany sætter Tanken i Bevægelse, trods Alt hvad der kan
dæmpe og hæmme, forstaas kun, naar man har dette for Øje. Og Personligheden paa
sin Side trænger til at lutres ved at bøje sig for Tankens objektive
Sammenhæng, den faste Tingenes Orden, indenfor hvilken ethvert enkelt Væsen har
sin bestemte Plads. Friheden vindes gennem Sandheden. Det hører til
Personlighedens Adel, at den ved egen Kraft kan finde de store Love, som
betinge dens Bestaaen, bestemme de Ønsker, dér opstaa hos den, og bestemme
Betingelserne for, at disse Ønsker kunne virkeliggøres. Der gives en Videnskab
om selve det personlige Liv, saavel som om alt Andet; denne Videnskab er maaske
ufuldkomnere end anden Videnskab, men den forudsættes ved énhver alvorlig
Drøftelse af det personlige Livs Former og Krav. Selv \textsc{Charles
Renouvier,} den nulevende Filosof, der lægger størst Vægt paa Modsætningen
mellem Tankenødvendighed og Personlighed, kræver dog „et \textit{rationelt}
Personlighedsbegreb“\textsuperscript{1}) og kan altsaa ikke afvise enhver
Analogi mellem dem.

% --- p. 3 ---
\opage{3}En saadan Analogi vil træde klart frem, naar man forsøger at formulere
de Hovedproblemer, hvormed den filosofiske Forsken beskæftiger sig. Disse
Problemer stilles paa én Gang fra Personlighedens og fra Videnskabens Side.

Ad rent historisk Vej har jeg i mit Værk \textit{Den nyere Filosofis Historie}
søgt at vise, at der gives fire saadanne Hovedproblemer, nemlig I) Problemet om
Bevidsthedslivets Natur (det psykologiske Problem), II) Problemet om
Erkendelsens Gyldighed (det logiske Problem), III) Problemet om Tilværelsens
Natur (det kosmologiske Problem) og IV) Vurderingsproblemet (det
etisk-religiøse Problem)\textsuperscript{2}). Det vil nu i denne Afhandling
være min Opgave at vise den indre Sammenhæng mellem disse Problemer. Det er et
og samme Grundproblem, der i dem fremtræder i forskellige Former og Anvendelser.

Det er meget forskellige Motiver, der kunne føre til filosofisk Forsken. ---
Meget ofte er det en etisk-religiøs, altsaa praktisk personlig Interesse, der
fører til Filosofi. Man begynder da med det fjerde Problem. Men det viser sig
snart, at dets Behandling kræver Indsigt i Bevidsthedslivets Natur, i
Erkendelsens Betingelser og i Beskaffenheden af den Tilværelse, i hvilken den
vurderende Personlighed existerer. --- Hvis det er en Interesse for
Iagttagelse, der ligger til Grund, vil man begynde med det psykologiske
Problem, saaledes som det er sket i den empiriske Filosofi i dens forskellige
Former. --- Hvis det derimod er Trang til at skelne mellem hvad der kan vides
og hvad der ikke kan vides, vil man begynde med Erkendelsesproblemet; denne Vej
gaar den kritiske Filosofi. --- Hvis man har fuld Tillid til Tankens Evne og
søger en afsluttet Verdensanskuelse, vil man --- som de dogmatiske og
spekulative Retninger --- begynde med Tilværelsesproblemet. --- Det vil
naturligvis ikke være ligegyldigt for Problembehandlingens Art og Maade,
hvilket Problem man begynder med; det ene Problem vil let blive stillet i
Skygge for det andet. For en sammenlignende Problemlære, som jeg her søger at
give et Bidrag til, vil det gælde om at lade hvert enkelt Problem komme til
dets fulde Ret. Jeg gaar da saaledes frem, at jeg --- i Tilslutning til den
ovenfor antydede Analogi mellem Personlighed og Videnskab --- først fremdrager Pro%
% --- p. 4 ---
\opage{4}blemet om det personlige Livs Natur, altsaa det psykologiske Problem,
og derefter gaar over til at undersøge den videnskabelige Erkendelses Natur,
altsaa Erkendelsesproblemet. Naar Tilværelsesproblemet stilles som det tredie i
Rækken, saa er Overgangen til det naturlig baade fra Personligheden, der staar
som en enkelt Del af den hele Tilværelse, og fra Videnskaben, hvis Opgave det
er at føre til en Verdensanskuelse. De tre hidtil nævnte Problemer vilde
stilles, selv om Mennesket var et rent intellektuelt, blot erkendende Væsen.
Det fjerde Problem stilles paa Grund af det Forhold, Mennesket som følende og
villende Væsen staar i til Tilværelsen. Saaledes hænge Tænkningens
Hovedproblemer sammen med Menneskets teoretiske og praktiske Interesser.

Men vigtigere end Spørgsmaalet om den Orden, i hvilken Problemerne fremstilles,
er Spørgsmaalet, om de kunne føres tilbage til et gennemgaaende Grundproblem.
Muligheden herfor finder jeg i den Betydning, Forholdet mellem Kontinuitet og
Diskontinuitet har indenfor ethvert af de enkelte Problemer. Dette Forhold
udtrykker baade Personlighedens og Videnskabens dybeste Interesse. Paa begge
Omraader stræbes der, som allerede bemærket, efter Enhed og Sammenhæng, og det
Diskontinuerlige staar forsaavidt som en Hindring, der skal overvindes. Men paa
den anden Side er det Diskontinuiteten (Forskellighed i Tid, Grad og Sted,
Kvalitetsforskel, Individualitetsforskel), der overalt, paa Videnskabens som
paa Livets Omraade, bringer nyt Indhold, udløser de bundne Kræfter og stiller
de store Opgaver. Intet af de to Elementer staar da paa Forhaand som ene
berettiget. Det vil have sin utvivlsomme Interesse at forfølge Forholdet mellem
dem under de fire forskellige Synspunkter, vore fire Hovedproblemer angive.

I det 19. Aarhundredes Filosofi fremtræder paa karakteristisk Maade
Kontinuitetsproblemets Betydning, idet de forskellige Retninger skiftevis tale
Kontinuitetens Sag. I Aarhundredets første Halvdel hævdede den filosofiske
Idealisme paa sin Vis Kontinuiteten i Tilværelsen og saa ned paa
Erfaringsvidenskaben paa Grund af dennes fragmentariske Karakter, medens
Positivismen (saaledes som Comte og Stuart Mill hævdede den) ind%
% --- p. 5 ---
\opage{5}skærpede Diskontinuiteten mellem de forskellige Grupper af Fænomener.
Ved Aarhundredets Slutning er det Realismen, der ved Udviklingshypotesens Hjælp
vil hævde Kontinuiteten, medens den idealistiske Retning er tilbøjelig til at
lægge Vægt paa den uundgaaelige Diskontinuitet i vor Viden\textsuperscript{3}).
Saaledes skifte Retningerne Plads i den store Holmgang, i hvilken der skal
kæmpes om Sandheden. Naar en betydningsfuld Tanke ikke længere kan gennemføres
ud fra et Synspunkt, da søger Forskningen uvilkaarligt et nyt Synspunkt, og
Tanken vinder i Klarhed ved disse skiftende Synspunkter. De forskellige
Retninger afløse hinanden, ligesom Løberne ved den antike Fakkelfest, men
Faklen er den samme. Og dersom ingen af de hidtil fremtraadte Retninger skulde
have ladet Grundtanken i alle de filosofiske Problemer vederfares fuld Ret, er
der saa meget des mere Grund til at arbejde paa at forme den med større Klarhed.

\begin{center}\rule{2cm}{0.4pt}\end{center}

% --- p. 6 ---
\opage{6}

\begin{center}I.\\[6pt]
\textbf{\Large Bevidsthedsproblemet.}\end{center}

1. Naar vi begynde Drøftelsen af de filosofiske Problemer med at undersøge
selve Personlighedsbegrebet, begynde vi med Psykologien, og det gælder da først
om at bestemme Psykologiens Forhold til Filosofien. Vi kunne her vistnok se
bort fra den Opfattelse, \textsc{Herbart} og \textsc{Lotze} endnu hævdede, at
Psykologien skulde være afhængig af den Opfattelse, man kommer til angaaende
Tilværelsesproblemet, altsaa afhængig af „Metafysiken“ eller Kosmologien. Der
gør sig fra alle Sider en Bestræbelse gældende for at hævde Psykologiens
Selvstændighed. Men selv hvor denne Bestræbelse raader, har der endnu været
yderst forskellige Anskuelser om Psykologiens Stilling.

Der er dem, der mene, at Psykologien er et med Filosofien i det Hele, og at
Erkendelsesteori, Metafysik, Æstetik og Etik ere dens forskellige Dele. Et
saadant Standpunkt indtoges i tidligere Tid af \textsc{Fries} og
\textsc{Beneke} og indtages i vore Dage af \textsc{Th.
Lipps}.\textsuperscript{4}) Dette Standpunkt er uholdbart, dels fordi det
personlige Liv kun er et af de Emner, der kunne være Genstand for Filosoferen,
dels fordi selve Psykologien, saavel som al anden Videnskab, forudsætter
Erkendelsens almindelige Former og Principer. Bevidsthedsproblemet viser
saaledes nødvendigvis ud over sig selv til Tilværelsesproblemet og
Erkendelsesproblemet, og trods al sin Selvstændighed er Psykologien kun en
Disciplin blandt flere. Psykologien kan da kun være en Del af Filosofien, ikke
udgøre hele Filosofien. Det er berettiget at begynde med
% --- p. 7 ---
\opage{7}Psykologien, da denne beskriver det Sted og de Forudsætninger, ud fra
hvilke vi søge at orientere os i Tilværelsen.

Men hører Psykologien nu i det Hele til Filosofien? Er den ikke en speciel
Videnskab, der i og for sig ikke har Mere med Filosofien at gøre end
Naturvidenskaben eller Historien? De specielle Veje, ad hvilke man i den nyeste
Tid søger psykologisk Kundskab, --- den experimentale, den fysiologiske og den
historiske Metode, --- synes jo dog at vidne om, at Psykologien nu er i Færd
med at blive en Specialvidenskab og maa udsondres fra Filosofien! --- Hertil
er, saavidt jeg ser, at svare, at Psykologien trods alle sine specielle Metoder
stedse forudsætter Evnen til Selviagttagelse. Det er den subjektive,
beskrivende og analyserende Psykologi, der stiller de Opgaver, der kunne
behandles i det Enkelte ved de specielle Metoders Hjælp. Der er her et
Grundlag, som er fælles for Psykologien som Specialvidenskab og for Filosofien.
Filosofien maa bygge paa et almindeligt Begreb om Bevidsthedslivets Natur, naar
den skal kunne behandle Erkendelses-, Tilværelses- og Vurderingsproblemet. Det
er Personlighedens Stilling som erkendende og vurderende og som Del af den hele
Tilværelse, disse Problemer angaa, og der forudsættes da et
Personlighedsbegreb, som ingen af hine specielle Metoder formaar at give. De
undersøge enkelte Ytringer af Bevidsthedslivet, men deres Undersøgelser maa
sammenarbejdes, for at et Personlighedsbegreb kan dannes.

Men her opstaar da netop Grundproblemet paa det psykologiske Omraade: Er det
muligt ad Erfaringens Vej at naa til et Personlighedsbegreb? Danner vort
Bevidsthedsliv en Helhed, et Kontinuum, en lille Verden for sig, eller skulde
det blot være et Aggregat, en Sum af Elementer og Fragmenter? Dette Spørgsmaal
er det, den filosofiske Psykologi (eller, om man hellere vil, Psykologiens
Filosofi) opkaster. Og i sin Behandling af det støtter den sig dels til
Selviagttagelsen og den deskriptive Psykologi, dels til den experimentale, den
fysiologiske og den historiske Metodes Resultater.

Erfaringen viser os større og mindre Afbrydelser af vort Bevidsthedsliv; der er
ubevidste Mellemrum mellem de bevidste Tilstande. Og naar vi nærmere undersøge
den enkelte Bevidst%
% --- p. 8 ---
\opage{8}hedstilstand, kunne vi indenfor den skelne forskellige Elementer, der
kunne genfindes i andre Tilstande, saa at den enkelte Tilstand, ikke mindre end
Bevidstheden i det Hele, kunde synes at være en Resultant af en Mangfoldighed
af Elementer eller Fragmenter, der staa som det egentligt Reale. Bevidstheden
vilde da være et Aggregat eller en Sum af enkelte Glimt. Et psykisk Kontinuum,
saaledes som det forudsættes i Personlighedsbegrebet, vilde der ikke gives.

Hertil maa nu først bemærkes, at det aldrig vil kunne afgøres med fuld
Sikkerhed, om Analysen er naaet ned til de sidste Elementer. Den Mulighed staar
stedse aaben, at de Elementer, vi stanse ved, og af hvilke vi ere tilbøjelige
til at tænke os Bevidstheden sammensat, selv ere sammensatte, og at der, hvis
vi kunde naa ned til endnu simplere Elementer, Elementer af anden eller tredie
Orden, vilde fremtræde en Kontinuitet, der ikke kunde paavises, saalænge vi
stansede ved Elementerne af første Orden. Vore Sansefornemmelser, som man har
været tilbøjelig til at betragte som en Art psykiske Atomer, bære dog, som
nyere Undersøgelser ad forskellige Veje have vist, Spor af
Sammensætningsprocesser af endnu mere elementær Art end dem, Iagttagelsen mere
direkte viser os.\textsuperscript{5})

Dette er dog kun en foreløbig Betragtning, der ikke fører til noget afgjort
Resultat, eftersom jo Diskontinuitetens Forsvarer meget vel kan bruge den til
at vise, at al den Kontinuitet, der kan paavises, kun er foreløbig. Hvad der
derimod er af afgørende Betydning, er, at de saakaldte psykiske Elementer
stedse ere bestemte ved den Sammenhæng, i hvilken de forekomme, og at det kun
er ved Abstraktion, man udenfor denne Sammenhæng tillægger dem Egenskaber, de
kun have i selve Sammenhængen. Det er paa denne Grundtanke, min Fremstilling af
Psykologien er bygget. Alt Sjæleliv virker paa genial Maade, frembringer
umiddelbart og uvilkaarligt sammenhængende Fænomener og Processer, som Analysen
bag efter paa mere eller mindre kunstig Maade søger at opløse i „Elementer“.
Hvad Sansefornemmelsen angaar, kunne vi her beraabe os paa Fechners Lov og paa
Kontrastloven, ifølge hvilke hver enkelt Fornemmelses Styrkegrad og Kvalitet
ere bestemte ved hele
% --- p. 9 ---
\opage{9}den Sammenhæng, i hvilken den bliver til. Sammenhængen kan altsaa ikke
opfattes som Produkt af de psykiske Elementer, da disse netop kun blive til som
disse bestemte Fornemmelser i Kraft af Sammenhængen. Forekom de i en anden
Sammenhæng, vilde de ikke være de samme Fornemmelser. Analogt er Forholdet,
hvad Forestillingerne angaar. Forestillingernes Association finder sin sidste
Forklaring i det Unaturlige i de enkelte Forestillingers Isolering; der gør sig
stedse --- og i højere Grad, jo mere livskraftig Bevidstheden er --- en Tendens
gældende til Supplering eller Udvidelse, hvorved den enkelte Forestilling
indgaar Forbindelse med andre Forestillinger efter bestemte Love. Den saakaldte
„Association-Psychology“ (til hvilken jeg selv undertiden med Urette er bleven
regnet) opfattede de enkelte Forestillinger som selvstændige Atomer, der bleve
bragte i Forbindelse med hverandre paa rent ydre, mekanisk Vis. Men det er
netop omvendt Helhedssammenhængen, der i Associationsprocessen øver sin Magt
over de enkelte Forestillinger, der tilmed aldrig forekomme i aldeles isoleret
Tilstand, saa at deres Forbindelse skulde kunne staa som et mekanisk Produkt.
Analysen forudsætter ogsaa her stedse Syntesen. Under den egentlige Tænkning
viser dette Forhold sig tydeligt, naar vi sammenligne Domsdannelse og
Begrebsdannelse. Disse forudsætte indbyrdes hinanden, idet vore Domme, der jo
ere Begrebsforbindelser, kun kunne være fuldkomne, naar de Begreber, der
forbindes i dem, ere fuldkomne, medens omvendt Begrebsdannelsen forudsætter en
Række af Domme, ved hvilke forskellige Elementer bestemmes i deres indbyrdes
Forhold. Det viser sig atter her, at Helheden og Delene gensidigt bestemme
hverandre. Der gives ingen aldeles isolerede Begreber, som først bagefter
skulde kunne forbindes til Domme.

Der ytrer sig her en Antinomi, der hænger nøje sammen med Bevidsthedens Væsen
og er ejendommelig for Personlighedsbegrebet. Bevidsthed og Personlighed kunne
ikke forklares som Produkter af forud givne Elementer, ligesaa lidt som det
organiske Liv kan forklares som Produkt af uorganiske Elementer. Men paa den
anden Side lægger Bevidsthedens og Personlighedens Væsen sig, ligesom det
organiske Livs Væsen,
% --- p. 10 ---
\opage{10}for Dagen i en stadig Sammenfatten af \textit{givne} Elementer,
Elementer, de ikke selv fra først af frembringe. Det er denne Antinomi, der gør
Livets og Personlighedens Opstaaen til saa store Gaader.\textsuperscript{6})

I det Foregaaende er der kun taget Hensyn til den mere formale Sammenhæng
indenfor Bevidsthedslivet. Men i enhver Bevidsthed er der et Maal, mod hvilket
der stræbes, en ledende Interesse, der tager eller stræber at tage Alt i sin
Tjeneste. Denne ledende Interesse --- Grundværdien kunde man kalde den --- kan
skifte i de forskellige Livsaldre; men der vil stedse paany være Tendens til
Udviklingen af en saadan, og den sætter i højere eller lavere Grad sit Præg paa
alle Bevidsthedens Elementer, giver dem deres Retning. Den konstituerer den
egentlige Sjæl (naar vi ved „Aand“ nærmest tænke paa den formale Sammenhæng).

Det er de her fremdragne Forhold, der gøre, at Personlighedsbegrebet stedse maa
staa som den centrale Tanke i Psykologien. Naar Analysen og de specielle
Metoder foretage deres Udsnit og søge at isolere enkelte Elementer eller
Øjeblikke, saa er det berettiget, ligesom det er berettiget at forsøge at
bestemme et irrationelt Tal ved at føje Decimal til Decimal, saa langt man kan.
Men det irrationelle Forhold mellem Helheden og Elementerne vedbliver at bestaa.

Fra idealistisk Side har man ofte været tilbøjelig til at antage
Personlighedsbegrebet for afsluttet og operere med det i kosmologiske
Spekulationer. Det overses herved, hvad den positivistiske Retning særligt
betoner, at vi stadigt arbejde paa at danne Personlighedsbegrebet, at stave os
frem til en psykologisk Helhedsopfattelse, ligesom Biologien endnu stadigt
staver paa en Definition af Livet. Og ligesom Biologien trods Anerkendelsen af
den levende Organismes Ejendommelighed ikke kender anden Metode end den, der
gennem Iagttagelse, Experiment og Analyse søger at forstaa de sammensatte
Processer ved Hjælp af de simplere, saaledes kan Psykologien, hvor meget den
end fastholder Bevidsthedslivets syntetiske Karakter, ikke anvende andre
Metoder end dem, al Videnskab anvender, Iagttagelse, Experiment og Analyse.
Personlighedsbegrebet staar
% --- p. 11 ---
\opage{11}som Idealet, henimod hvilket der styres, det stadige Problem, hvis
Belysning alle specielle Metoder tjene.\textsuperscript{7})

Det Irrationelle er her som overalt ikke blot Skranke, men ogsaa Opgave, en
Opgave, der stedse stilles paany. Den deskriptive Psykologi, som særligt fører
til at lægge Vægt paa den Helhedssammenhæng, i og med hvilken de sjælelige
Fænomener fremtræde, vil her --- især hvad de højere eller mere udviklede
Bevidsthedsfænomener angaar --- stedse bevare sin Selvstændighed og Betydning i
Forhold til den experimentale Psykologi; denne faar tilsidst stedse sine
Opgaver fra hin. Og den deskriptive Psykologi tjener tillige som Korrektiv
overfor den experimentale Psykologi, der ifølge sin Metode let kommer til at
isolere enkelte Elementer for meget, til at overse det rent Uvilkaarlige i
Bevidsthedslivet og til at lægge for stor Vægt paa Tilstandenes ydre
Symptomer.\textsuperscript{8}) Paa den anden Side vil den deskriptive Psykologi
kunne drage Fordel af den experimentale Psykologi, idet det, denne formaar at
paavise om de mere elementære psykologiske Processer, ad Analogiens Vej kan
kaste Lys over de højere Processers Natur og give Beskrivelsen større Fylde og
Nøjagtighed.

Den beskrivende Psykologi nærmer sig til Kunsten og bliver tilsidst en Kunst;
den experimentale Psykologi nærmer sig til Naturvidenskaben. Men der er
alligevel ingen principiel Forskel mellem dem. Der er ikke, som i den nyeste
Tid især \textsc{Münsterberg}\textsuperscript{9}) har villet hævde, et Svælg
imellem dem. Hin skal være „Værdividenskab“, denne Naturvidenskab; hin skal
operere med „Frihedsbegrebet“, denne med Aarsagsbegrebet. Ved Frihed forstaar
Münsterberg Evne til at handle efter Formaal. Men skulde da den psykiske
Proces, hvorved et Menneske kommer til at stille sig et Formaal og arbejde paa
at opnaa det, ikke være Genstand for videnskabelig Psykologi? Og kan den
forstaas uden ved at undersøges i Sammenhæng med de simplere Processer, der
foregaa, naar Lyst og Smerte, Glæde og Sorg opstaa, eller naar Forestillinger
forbindes og skilles? Her er en Mængde Erfaringer at gøre og at samle; det
gælder at forfølge en Udvikling Skridt for Skridt, at konstatere de forskellige
Bevidsthedsgrader og de forskelligartede Impulser og Motiver, der
% --- p. 12 ---
\opage{12}medvirke. Der vil fremdeles være at skelne mellem de forskellige
individuelle Typer, der kunne fremtræde med Hensyn til Villieslivets Art og
Retning, og en sammenlignende Individualpsykologi vil kunne være et nødvendigt
Hjælpemiddel til at forstaa Udviklingen indenfor den Enkeltes Bevidsthed. Det
er da en aldeles unaturlig Grænse, der sættes, naar man fastslaar en principiel
Modsætning mellem „personlige“ og „psykofysiske“ Kategorier --- og da især,
naar denne Modsætning lægges til Grund for en Modsætning mellem „Livets
Sandhed“ og „Videnskabens Sandhed“.\textsuperscript{10}) Det er jo dog
Videnskabens Opgave at forstaa Livet, selv om den fuldkomne Forstaaelse stedse
er et Ideal, og det er Livet, der yder Videnskaben det Erfaringsmateriale,
hvorpaa den bygger. Her er et stadigt Vexelvirkningsforhold, men ingen absolut
Modsætning. Og der er mange Mellemled og Mellemformer mellem Videnskabens
Sondren af Tilværelsens enkelte Elementer og Livets store Sammenspil af alle
Elementer, som tilsidst kun Kunsten kan fremstille i dets hele Fylde.

2. Imidlertid bliver der stedse store Vanskeligheder tilbage for den
videnskabelige Behandling af Bevidsthedsfænomenerne. Til en virkelig
Forstaaelse hører, at der paavises en Kontinuitet i de Processer, man
undersøger. Det vilde være Psykologiens Ideal ikke blot, om vi kunde vinde en
saa fuldkommen Beskrivelse af Bevidsthedstilstandene, at hver enkelt af disse
kom til at staa som ejendommeligt Led i den hele psykiske Proces, men ogsaa om
vi kunde reducere Forskellighederne mellem de skiftende Tilstande til saa
simple Former, at de følgende optraadte som Fortsættelser eller som Omsætninger
af de foregaaende. Men Udsigten til at naa dette Ideal spærres ved den
Diskontinuitet, Erfaringen synes at godtgøre. Der er jo ubevidste Mellemrum
mellem vore bevidste Tilstande --- i Afmagt, i den drømmeløse Søvn (forsaavidt
der gives en saadan), og der er kvalitative Forskelligheder mellem de
forskellige Bevidsthedstilstande og Bevidsthedselementer, saa at hver enkelt
Tilstand og hver enkelt Element synes at blive til af Intet, naar vi holde os
til den strengt psykologiske Betragtning. Imellem de forskellige individuelle
Bevidstheder er der i hvert Tilfælde en stor
% --- p. 13 ---
\opage{13}og iøjnefaldende Diskontinuitet. Den ene Bevidsthed kan endnu mindre
udledes af den anden, end den ene Bevidsthedstilstand kan udledes af den anden.

Disse psykiske Diskontinuitetsforhold træde særligt skarpt frem, naar de
stilles i Modsætning til den Kontinuitet og de Ækvivalensforhold, de materielle
Fænomener frembyde, og som man meget tidligt blev opmærksom paa. Hos de græske
Naturfilosofer var Materiens Bestaaen under alle Forandringer en fælles
Antagelse, og denne Antagelse har faaet sin Bekræftelse ved den nyere Tids
Kemi. I skarp Modsætning hertil stod da Afstanden mellem Sjælene indbyrdes.
Naar \textsc{Descartes} anvendte Substansbegrebet baade paa det psykiske og paa
det fysiske Omraade, var det med den store Forskel, at han antog mange
Sjælesubstanser, men kun én eneste materiel Substans.\textsuperscript{11})
Herved udtaltes skarpt og bestemt den psykiske Diskontinuitet i Modsætning til
den fysiske Kontinuitet. Det er kun paa halvt eller helt mystisk Vis, at vi i
Filosofiens Historie finder Sjælelivets Kontinuitet hævdet saaledes, at de
enkelte Bevidstheder tænkes i samme Forhold til en universel Sjælesubstans, som
de enkelte Legemer staa i til den universelle Materiesubstans. Saaledes hos
\textsc{Averroes, Spinoza} og \textsc{Hegel}.\textsuperscript{12}) I den nyeste
Tid er det nu ikke saa meget Diskontinuitetsforholdet mellem de individuelle
Bevidstheder, som Diskontinuitetsforholdet indenfor den enkelte Bevidsthed, der
tiltrækker sig Opmærksomheden. Og der er navnlig to Betragtninger, som her have
ført til at stille Problemet skarpere end før.

Fysiologien er kommen til fuld Bevidsthed om sin Metode og om sin
Selvstændighed som Videnskab, især efter den Reform, den undergik ved Midten af
det 19. Aarhundrede. Herefter kræver den nu konsekvent, at enhver
Hjernetilstand skal forklares ud fra Hjernens og Organismens foregaaende
Tilstande, --- at i det Hele fysiologiske Fænomener skulle forklares af
fysiologiske Aarsager, forsaavidt det ikke er Virkninger fra den fysiske
Omverden, der fortsættes ind i Organismen. Efter streng naturvidenskabelig
Metode maa enhver materiel Tilstand i og udenfor Organismen finde sin
Forklaring ved at ses som opstaaet ved Omsætning i ny Form af den Energi, der
ytrede sig
% --- p. 14 ---
\opage{14}i de foregaaende Tilstande. Naar to Hjernetilstande afløse hinanden,
vil altsaa Forklaringen først være naat, naar det paavises, at al den Energi,
der ytrede sig i den første Tilstand, uden Rest er bleven omsat til
Ernæringsvirksomhed, Varme, Elektricitet, motoriske Impulser o. s. v. Det er
vel vanskeligt at tænke sig et saadant Experiment udført; men det vilde betegne
Gennemførelsen af den fysiologiske Metode. Selve Nyvitalismen kender ingen
anden Metode til Undersøgelse af de fysiologiske Tilstande. Saaledes bekræftes
\textsc{Spinoza's} Ord, at naar man til Forklaring af en legemlig Tilstand
beraaber sig paa Sjælens Indgriben, er det det samme som at sige, at man ikke
véd, hvad Aarsagen er. Selv om et Forsøg skulde vise, at der ved en organisk
Tilstands Opstaaen opstod eller forsvandt Energi uden fysisk Ækvivalent, vilde
man sikkert snarere tro, at der var en Fejl ved Experimentet, end slaa sig til
Ro ved Resultatet. „Principet om Energiens Bestaaen har“, siger
\textsc{Maxwell}\textsuperscript{12b}), „faaet saa stor videnskabelig Vægt, at
ingen Fysiolog vilde have nogen Tillid til et Experiment, der viste en
betydelig Forskel mellem det af et levende Væsen udførte Arbejde og Opgørelsen
af modtagen og brugt Energi.“ --- Ved Gennemførelse af den fysiologiske Metode,
opfattet i denne Aand, vil der paavises stedse større Kontinuitet i de
organiske Processer, til hvilke Bevidsthedsfænomenerne faktisk vise sig
knyttede. Fysiologien er derfor langt heldigere stillet overfor
Kontinuitetsprincipet, end Psykologien nogensinde vil kunne blive det, og det
kunde derfor synes at være en velbegrundet Fordring, \textsc{Carl Lange}
opstillede, naar han erklærede, at alle psykologiske Definitioner burde afløses
af fysiologiske Definitioner\textsuperscript{13}). De psykiske Fænomener komme
da til at staa som blotte foreløbige Symptomer.

Den anden Betragtning slutter sig til den nu anførte, men er af almindelig
erkendelsesteoretisk Art. Den støtter sig til den Sætning, at en fuldkommen
Forstaaelse kun vindes, hvor Forholdet mellem Aarsag og Virkning kan gøres til
et Identitetseller Kontinuitetsforhold, saaledes at der bliver en kvantitativ
Ligning mulig mellem de Fænomener, der forbindes i Aarsagsforholdet. Det er det
ideale Aarsagsbegreb, man her hævder i
% --- p. 15 ---
\opage{15}Modsætning til det empiriske Aarsagsbegreb, ved hvilket Fænomener,
der ere kvalitativt forskellige, staa som forbundne i uundgaaelig Rækkefølge.
Den nærmere Undersøgelse af disse to Aarsagsbegreber opsætter jeg til
Erkendelsesproblemet. Her skal kun bemærkes, at der ingen Tvivl er om, at hvor
det ideale Aarsagsbegreb kan gennemføres, faa vi en fuldkomnere Forstaaelse,
end hvor vi maa stanse ved det empiriske Aarsagsbegreb og maaske endda have
Vanskelighed nok ved at anvende dette. Nu frembyde baade de materielle og de
psykiske Fænomener kvalitative Forskelligheder; men paa det materielle Omraade
er det muligt at opfatte disse kvalitative Forskelligheder som
Kvantitetsforskelligheder, nemlig formedelst Loven om Energiens og Stoffets
Bestaaen, medens noget Tilsvarende ikke er muligt paa det psykiske Omraade.
Derfor antages den rette psykologiske Metode at bestaa i at substituere de
fysiologiske Fænomener for de tilsvarende psykiske; det er den eneste Maade,
paa hvilken vi kunne sætte kvantitative Bestemmelser i Stedet for kvalitative
og derved gennemføre en stringent Aarsagssammenhæng. Strengest gennemført er
denne Opfattelse af det psykologiske Problem af \textsc{Richard Avenarius} i
hans skarpsindige Værk \textit{Kritik der reinen Erfahrung.} Han fordrer her
som Betingelse for fuld videnskabelig Forstaaelse en Tilbageførelse af „den
afhængige Vitalrække“, hvormed menes de psykiske Tilstande, til „den uafhængige
Vitalrække“, de tilsvarende fysiologiske Tilstande, hvorved det ene bliver
muligt at reducere de empiriske Forskelligheder til det mindst Mulige, eller,
som Avenarius udtrykker det, at naa til „et heterotisk Minimum“. I lignende
Retning gaar \textsc{Münsterberg's} Anskuelse, saaledes som han har udtalt den
i sit Skrift \textit{Psychology and Life.} Hans Hovedpaastand er, at da
psykiske Fænomener ikke ere kvantitative, kunne de ikke være Led i en
Aarsagsrække\textsuperscript{14}).

Overfor de Anskuelser, der saaledes fordre en Reduktion af Psykologi til
Fysiologi, for at en videnskabelig Psykologi skal blive mulig, --- altsaa
egentligt ville ophæve Psykologien for at gøre den til Videnskab, --- maa det
anerkendes, at de kunne støtte sig til den faktiske Diskontinuitet og den
kvalitative Forskellighed mellem psykiske Fænomener, der stedse mere
% --- p. 16 ---
\opage{16}eller mindre vil sætte Grænsen for en streng Gennemførelse at
Psykologien som Videnskab. En anden Sag er, om den Vej, de anvise, fører til
Maalet, og om den Metode, de anprise, ikke indeholder en haandgribelig
Selvmodsigelse.

Naar det fordres, at de psykologiske Definitioner skulle \textit{afløses} af
fysiologiske Definitioner, saa er det klart, at Forudsætningen herfor er, at de
psykologiske Definitioner først maa være færdige. Men at skaffe disse
Definitioner til Veje maa være Psykologiens Sag, og kan den ikke skaffe
bestemte og sikre Definitioner til Veje, kan Fysiologien ikke vide, hvad det
er, den skal søge Forklaring paa i Hjernen: hvis det, der skal afløses, er vagt
og vaklende, saa vil ogsaa Afløsningen blive derefter. Der vil ingen Sikkerhed
være for, at man virkeligt har fundet de Hjernetilstande, som svare til de
psykiske Fænomener, man har for Øje. Psykologiens Selvstændighed maa altsaa i
hvert Tilfælde anerkendes, da det her bliver den, der --- som en Art
Symptomatologi --- skal stille Fysiologien dens Opgaver. Nu er det for enhver
Erfaringsvidenskab et langt og vanskeligt Arbejde at finde afsluttende
Definitioner; saadanne ere først mulige, naar Erfaringsvidenskaben i det
Væsentlige er fuldendt; de komme i Slutningen, ikke i Begyndelsen af
Undersøgelsen. Der opereres altfor ofte med ufærdige psykologiske Definitioner,
der betragtes som tilforladelige Udgangspunkter for hjernefysiologisk
Undersøgelse. Saaledes betragtede \textsc{Descartes}, og i nyere Tid
\textsc{Lotze}, selve Begrebet „Sjæl“ for saa tydeligt og klart, at man trygt
kunde opfordre Fysiologer og Anatomer til at søge efter „Sjælens Sæde“. I
nyeste Tid har \textsc{Flechsig}\textsuperscript{15}) betragtet Begrebet
„Association“ som saa simpelt og klart og tillige for saa selvstændigt, at han
har kunnet mene at finde en særegen Plads i Hjernen for de Funktioner, dette
Begreb betegner. At det dog er et højst ufuldkomment Associationsbegreb,
Flechsig herved lægger til Grund, ses af, at psykologisk Erfaring aldrig viser
os absolut enkelte psykiske Elementer, der først ved en særegen Proces skulle
bringes i Forbindelse med hverandre. Man kan ikke stille Associationsprocessen
i bestemt Modsætning til de Processer, ved hvilke de enkelte psykiske Elementer
(Fornemmelser og Forestillinger) blive til.
% --- p. 17 ---
\opage{17}Det er en psykologisk Abstraktion, ikke en virkelig psykologisk
Definition, Flechsig opererer med. Der er mod Flechsigs Lære om særegne
Associationscentra fra flere Sider rejst anatomiske Indvendinger, jeg ikke kan
bedømme; men det psykologisk Ufærdige i hans Associationsbegreb viser
tilstrækkeligt, hvor vanskelig den Opgave er at „afløse“ psykologiske
Definitioner med fysiologiske\textsuperscript{16}). Et Hovedpunkt til Kritik af
Flechsigs Teori synes følgende Omstændighed mig at være. Naar der er psykiske
Elementer, der ikke indgaa en saadan Association, som der under sædvanlige
Forhold kunde ventes, saa vil ved nærmere Undersøgelse Aarsagen vise sig at
være den, at vedkommende Elementer hvert for sig ere Led i andre, faste
Associationsforhold, af hvilke de ikke kunne udløses. Det vil i hvert Tilfælde
stedse psykologisk set staa som en Opgave at finde de foregaaende
Associationer, der hindre senere i og for sig „naturlige“ Associationer. Ved at
isolere de to Begreber Element og Association fra hinanden vil man saaledes
baade faa en falsk Psykologi og en falsk Fysiologi. Til en udtømmende
Definition af et Begreb som Association vil der endnu behøves lange
Undersøgelser, og Fysiologien vil komme til at vente længe, hvis den vil
„afløse“ en psykologisk Realitet og ikke en blot Abstraktion.

Men lad os tænke os, at de psykologiske Definitioner ere færdige og fuldkomne,
ville de største Problemer saa ikke endnu staa uløste tilbage? Hvorledes kunne
fysiologiske Tilstande have psykiske Symptomer? Hvorledes kunne kvalitative
Forskelligheder svare til kvantitative, og hvorledes kan der til kontinuerlige
Processer være knyttet diskontinuerlige Fænomener? Langt fra, at saadanne
Spørgsmaal falde bort ved Reduktionen af Psykologi til Fysiologi, skærpes de
netop ved denne og fremtræde i en mere grel Form. Og i hvert Tilfælde staar det
som en uløst Gaade, hvorledes de kvalitative Forskelligheder og
Diskontinuiteten opstaar. Selv om det lykkes at paavise, at to Begreber
\textit{A} og \textit{B,} man hidtil har betragtet som forskellige, ere
identiske, saa at man kan sætte \textit{A = B,} har man derved ikke forklaret,
hvorledes \textit{A} og \textit{B} fra først af kunde staa som forskellige. At
kalde denne Forskellighed blot „subjektiv“ hjælper
% --- p. 18 ---
\opage{18}ikke; derved skaffes den ikke ud af Tilværelsen; det er jo netop den
Kendsgerning, at Tilværelsen ogsaa har en subjektiv Side, der gør Psykologi
mulig og nødvendig. Og erklærer man (med \textsc{Münsterberg}), at disse
subjektive Fænomener hverken kunne beskrives eller
forklares\textsuperscript{17}), er det ikke godt at forstaa, hvorledes man faar
Noget at sætte i „Kausalligning“. Man bærer sig da ad som Manden, der saver den
Gren over, paa hvilken han sidder.

Det er inkonsekvent at negte et Kausalitetsforhold mellem psykiske Fænomener
indbyrdes, naar man antager et Kausalitetsforhold mellem de tilsvarende
fysiologiske Tilstande Anerkender man en fysiologisk Lovmæssighed, maa man
ogsaa anerkende en psykologisk Lovmæssighed. Hvis det psykiske Fænomen αsvarer
til Hjernetilstanden \textit{A,} det psykiske Fænomen \textit{β} til
Hjernetilstanden \textit{B,} og hvis der er et Aarsagsforhold mellem \textit{A}
og \textit{B,} saa maa der ogsaa i psykologisk Erfaring fremtræde et
Aarsagsforhold --- i Betydning af uundgaaelig Rækkefølge --- mellem α og
\textit{β.} Og --- som det i det Foregaaende er paavist --- det er i
Virkeligheden dette Aarsagsforhold mellem αog \textit{β}, der gør, at man
kommer til at søge efter et Aarsagsforhold mellem \textit{A} og \textit{B.} I
\textsc{Avenarius} ' strenge Gennemførelse af Reduktionen af Psykologien (Læren
om „den afhængige Vitalrække“) til Fysiologi (Læren om „den uafhængige
Vitalrække“) ses det ogsaa paa ethvert Punkt tydeligt, at det er fra den
„afhængige Vitalrække“, han slutter til den „uafhængiges“ Beskaffenhed. For det
Meste er det kun en fysiologisk Skematisme, man konstruerer som „anskuelig“
Fremstilling af, hvad der gaar for sig i de psykiske Processer. Hverken
Fysiologien eller Psykologien er endnu saa fuldkommen, at man kan komme ud over
en saadan Skematisme.

De psykiske Fænomener frembyde nu dog ikke saa stor Diskontinuitet og saa rent
kvalitative Forskelligheder, som man ofte mener. Nøjagtig Iagttagelse fører
ofte til at opdage, at hvor der syntes at foreligge et psykisk tomt Rum, var
der i Virkeligheden et psykisk Indhold givet, skønt det ikke fattedes af
Opmærksomheden eller af Genkendelsen, eller skønt det glemtes strax efter
Oplevelsen. Og hvad de kvalitative Forskelligheder angaar, vil ogsaa her
nøjagtigere Iagttagelse finde flere og finere Nuancer end dem, man først
stansede ved og maaske trygt
% --- p. 19 ---
\opage{19}erklærede for aldeles disparate. Kontinuiteten i Bevidsthedslivet kan
altsaa føres videre, end man ofte antager, og der kan ligge en videnskabelig
Fare i for stærkt at proklamere Kvalitetsforskellighed og Diskontinuitet som
Bevidsthedslivets Særkende. Den indre Sammenhæng i Bevidsthedens Indhold er
ogsaa en Kendsgerning, og Funktioner som Erindring og Sammenligning ere nok saa
betydningsfulde til Karakteristik af Bevidsthedslivet som de Fænomener, der
bære Diskontinuitetens Præg. Psykologiens Opgave bliver derfor naturligt den,
saavidt muligt at paavise de enkelte Elementers Sammenhæng og Forbindelse,
saaledes at Helheden forstaas ved Delene, og Delene ved deres Forhold til
Helheden. At dette fører til en Antinomi, der tilsidst gør det psykologiske
Problem uløseligt for os, have vi allerede set. Men der er Arbejde nok at gøre,
før vi komme til denne Grænse (som desuden ikke kan bevises at være absolut).
\textsc{Leibniz} har den Fortjeneste med Energi at have hævdet
Kontinuitetsprincipet paa det psykiske, saavel som paa det fysiske Omraade,
især ved at henvise til de --- i Forhold til de klart bevidste Tilstande ---
forsvindende smaa psykiske Elementer, der først ved at opsummeres og kombineres
frembringe et for uøvet Selviagttagelse tilgængeligt Resultat. Allerede en
Sammenligning mellem moderne og antik eller middelalderlig Poesi viser, hvor
langt flere sjælelige Nuancer, Selviagttagelsen nu har faaet Øje paa; det er
neppe blot, fordi Differentiationen paa det psykiske Omraade nu er større end i
tidligere Tider, at denne større Rigdom findes. Især er Kendskabet til og
Forstaaelsen af det Umiddelbare og Uvilkaarlige i Sjælelivet tiltagen. Den
barnlige og primitive Psykologi har en Tilbøjelighed til at opfatte Alt, hvad
der foregaar i Sjælen, som klart bevidst og som hvilende paa Overlæg og Forsæt.
Bevisbyrden er i den moderne Psykologi --- baade i Praxis og i Teori --- bleven
flyttet: nu paahviler den den, der paastaar, at en Handling er øvet med Overlæg
og udtrykkelig Beslutning. Og selve Overgangen mellem det Ubevidste og det
Bevidste, og indenfor det Bevidste mellem det Uvilkaarlige og det Vilkaarlige,
foregaar kontinuerligt; det bliver, jo mere øvet Selviagttagelsen og
% --- p. 20 ---
\opage{20}Kritiken er, stedse vanskeligere at sætte Fingeren paa ét bestemt
Punkt, hvor Grænsen skulde ligge.

Leibniz var tilbøjelig til at reducere alle psykiske Forskelligheder til
Gradsforskelligheder i Henseende til Klarhed og Dunkelhed\textsuperscript{18}).
Dette Forsøg var karakteristisk for Oplysningsaarhundredets Intellektualisme.
Men der er andre Maader end en saadan Reduktion, paa hvilke den empiriske
Psykologi kan forestille sig en kontinuerlig Sammenhæng trods de kvalitative
Forskelligheder, de psykiske Processer frembyde i deres forskellige Stadier.
Der foregaar Sammensmeltninger og Sammensætninger paa Bevidsthedslivets
forskellige Trin, lige fra de simpleste Sansefornemmelser til de højeste
Følelsestilstande. Og der foregaar Motivforskydninger, ved hvilke det, der
først kun har middelbar Værdi, senere faar umiddelbar Værdi, eller omvendt.
Endeligt viser Bevidsthedslivet hos de helstøbte Karakterer, Psykologiens
højeste og ædleste Genstande, en saadan Afhængighed af Alt, hvad der rører sig
i Sjælen, over for ét Formaal, én ledende Tanke, at vi her finde en
Kausalsammenhæng, der er nok saa fast og inderlig som nogen paa det fysiske
Omraade. Hver enkelt Tanke, hver enkelt Stemning og hver enkelt Drift er i en
saadan Karakter tydeligt bestemt ved den Helhed, den hører til, og som den selv
bidrager til at forme.

Der bliver Tilfælde nok tilbage, hvor den psykiske Kontinuitet ikke kan
paavises. Spørgsmaalet er, om vi derfor have Ret til at negte den. Afset fra,
hvad fortsat Iagttagelse kan konstatere, opstaar her det store Spørgsmaal, om
de psykiske Fænomener kunne blive til af materielle Aarsager, naar man opfatter
Begrebet om det Materielle paa den Maade, Naturvidenskaben hidtil har gjort
det. Naar det naturvidenskabelige Begreb om det Materielle ikke indeslutter
nogen Mulighed for en Opstaaen af psykiske Fænomener, saa vil det være fuldt
berettiget at opstille Begrebet om potentiel psykisk Energi for at markere, at
vi ikke ville slippe Kontinuitetsprincipet selv der, hvor vi ikke i Erfaringen
kunne anvende det i det Enkelte. Det er jo af en lignende Grund, at Begrebet
potentiel Energi er indført i Fysiken, trods den Dunkelhed, hvormed det er
behæftet. Det indeholder overalt en Anerkendelse af en Grænse, hvor man
% --- p. 21 ---
\opage{21}dog ikke vil opgive den Sammenhæng, man hidtil har fundet. Allerede
den Kendsgerning, at psykiske Elementer kunne reproduceres efter et Mellemrum,
i hvilket de ikke have været i Bevidstheden, nødvendiggør Anvendelsen af
Begreber som „Disposition“, „Spor“, „Mulighed“ o. lign., der i Virkeligheden
udtrykke det Samme, som menes ved „potentiel Energi“. At vi ved Anvendelsen af
dette Begreb paa det fysiske Omraade ere heldigere stillede end paa det
psykiske Omraade\textsuperscript{19}), udelukker ikke Berettigelsen og
Nødvendigheden af at opstille det her. Man kan naturligvis i en rent
beskrivende Fremstilling lade sig nøje med at angive de forskellige psykiske
Fænomener, der opdukke i visse Øjeblikke, ofte uden at Iagttagelsen finder
nogen psykisk Sammenhæng mellem dem indbyrdes. De psykiske Fænomener staa da
som spredte Glimt, i stærk Kontrast til den kontinuerlige Sammenhæng, de
tilsvarende fysiologiske Fænomener frembyde. Da nu det, der fremtræder med den
højeste Grad af Kontinuitet, ogsaa let kommer til at staa for os som mere realt
end det Diskontinuerlige og Glimtvise, saa ville de fysiologiske Fænomener let
komme til at staa som det egentlige Reale, de egentlige Fænomener, og de
psykiske Fænomener ville komme til at staa som noget Overflødigt, et tilfældigt
Plus, eller, som man har udtrykt det, som Epifænomener. En saadan
„Epifænomenisme“ er dog vist ikke nogensinde opstillet som egentlig Teori; den
er blot en empirisk Konstateren af det Gaadefulde i de psykiske Fænomeners
Optræden i Sammenligning med den mere sluttede Kontinuitet i de materielle
Fænomeners Række. Den giver jo heller ingen Forklaring, ophæver tvertimod baade
enhver psykologisk og enhver fysiologisk Forklaring, naar den vil være Mere end
en blot Beskrivelse. Og selv som Beskrivelse bliver den let forhastet, idet den
let kommer til at stanse for tidligt ved det Diskontinuerlige uden at
undersøge, om der ikke er flere Overgange og Nuancer at finde, end der hidtil
har vist sig. Begrebet om potentiel psykisk Energi vil her --- netop ved
egentligt kun at være Udtryk for en uløst Opgave --- indeholde en stadig
Opfordring til at gaa videre i Iagttagelse og Analyse for om muligt at opdage
den Virkelighed, hvori „Muligheden“ tilsidst altid bestaar.

% --- p. 22 ---
\opage{22}3. Naar man baade fra den psykologiske og den fysiologiske Side
lægger Vægt paa at bevare Kontinuiteten i saa høj Grad som muligt, bliver
Identitetshypotesen den egentlige Arbejdshypotese ved Problemet om det Psykiske
og det Fysiologiske. Som enhver betydningsfuld Hypotese er den egentlig kun
Udtryk for en Metode. Vi have i de psykiske og de fysiologiske Fænomener to
Rækker af Tilstande eller Tilstandsformer, som erfaringsmæssigt variere i et
vist Forhold til hinanden, uden at det dog er muligt at \textit{udlede} den ene
Rækkes \textit{Existens} af den anden Rækkes \textit{Existens.} Opgaven maa da
videnskabeligt set være den at opfatte de to Rækker hver for sig saa
fuldstændigt og kontinuerligt som muligt og at paavise, hvilke bestemte Led i
den ene Række der svarer til bestemte Led i den anden Række. Vi kunne med Rette
betragte Led i den ene Række som Symptomer paa Led i den anden Række. Snart er
det de psykiske, snart de fysiologiske Tilstande, der ere de mest tilgængelige,
og som vi derfor tage vort Udgangspunkt fra. Den nøje Sammenhøren mellem de to
Rækker af Tilstande gør det umuligt at henføre disse til to forskellige
„Væsener“ eller „Ting“; det vil være naturligt at opfatte dem som forskellige
Ytringsformer af et og samme „Væsen“. Deraf at to Egenskaber ikke kunne udledes
af hinanden, har man ikke Ret til at slutte, at de tilhøre to forskellige Ting,
især naar de variere i et bestemt indbyrdes Forhold. Svovlets Tilstandsform
(fast, flydende o. s. v.) og Farve (gul, brun o. s. v.) variere i et bestemt
Forhold til hinanden og kunne dog ikke udledes af hinanden, skønt man, naar man
kender de to Forandringsrækker, kan slutte fra den ene til den anden. Ved
Svovlet kende vi nu den fælles Aarsag til de to Forandringsrækker, nemlig
Varmen; det er den, der ytrer sig paa disse to forskellige Maader. Ved
Forholdet mellem de psykiske og de fysiologiske Rækker er den analoge Kundskab
os negtet. Vor Erfaring er her utilstrækkelig til at løse Problemet. Det hænger
sammen med, at baade vor Kundskab om det Psykiske og vor Kundskab om det
Fysiske er empirisk begrænset. Det gælder da at arbejde paa begge
Kundskabsomraader saaledes, at man ikke paa noget Punkt vilkaarligt hugger
Traaden over.

% --- p. 23 ---
\opage{23}Den fysiologiske Kontinuitet er baade en Konsekvens af Loven om
Energiens Bestaaen (egentligt allerede af Inertisætningen efter den strenge
Fortolkning af denne) og et Udtryk for Fysiologiens Selvstændighed som
Videnskab. Den afgiver et frugtbart og uundværligt Forskningsprincip ved at
fordre stadig Paavisning af fysiologiske Aarsager til og Virkninger af
fysiologiske Tilstande. Som vi have set ovenfor, er det et saa væsentligt
Princip, vi her staa ved, at man --- naar et Experiment pegede i modsat Retning
--- snarere vilde antage en Fejl i Experimentet end et Brud paa Principet.
Kontinuitetsprincipet --- og med det Identitetshypotesen --- vilde være
modbevist, dersom det kunde vises, at den Energi, der indeholdes i en
Hjernetilstand, ikke stod i Ækvivalensforhold til forudgaaende og efterfølgende
Tilstande i Hjernen, i Organismen og i den fysiske Omverden. Det vilde allerede
være betænkeligt for Identitetshypotesen, dersom det kunde paavises, at de
psykiske Fænomener i Tiden gaa \textit{forud for} eller følge \textit{efter} de
tilsvarende fysiologiske Tilstande. Men ved et Forsøg paa at godtgøre dette, er
det ikke nok at vise, hvad allerede populær Iagttagelse viser, at Fornemmelsen
opstaar senere end det tilsvarende Sanseindtryk, eller tidligere end de
fysiologiske (motoriske, særligt vasomotoriske) Ytringer, der paa en eller
anden Maade hænge sammen med den. Ti imellem de psykiske Tilstande og de
periferiske Processer i Sanseorganer og Bevægelsesorganer ligge de centrale
fysiologiske Processer, der ikke saa let ere tilgængelige for nøjagtige og
direkte Experimenter. Og selv om et saadant direkte Experiment var muligt,
vilde det ikke være afgørende: fordi Fænomenet \textit{B} i Tiden følger efter
\textit{A,} kunne \textit{A} og \textit{B} meget vel være Virkninger af samme
Aarsag, der først afsætter Virkningen \textit{A} og saa Virkningen \textit{B.}
Der kunde være særegne Aarsager, som gjorde, at den ene Livsytring ændrede sig
hurtigere end den anden. --- Endeligt vilde det være et Bevis imod
Identitetshypotesen, dersom det kunde godtgøres, at der til en og samme
fysiologiske Tilstand kunde svare forskellige psykiske Fænomener. Man har f.
Ex. tænkt sig, at Hjernens Betydning for Bevidsthedslivet kunde bestaa i at
modvirke det organiske Livs Tendens til vanemæssig, automatisk
% --- p. 24 ---
\opage{24}Virken. Hjernen har en vidunderlig Evne til at indlede nye
Bevægelser, der, saasnart de ere indøvede, blive til Vanehandlinger, som det
overlades lavere Centra at udføre. Kun ved Hjernens Hjælp har Tanken kunnet
befri sig fra Automatismen og har kunnet benytte Organismen i sin Tjeneste uden
at synke under Vanens Magt. Men forskellige Tanker kunne til forskellige Tider
gøre Brug af samme „motoriske Skema“, samme Ophævelse af en automatisk Tendens.
Denne Mulighed har i den nyeste Tid \textsc{Bergson} fremført som en Indvending
mod „den strenge Parallelisme“\textsuperscript{20}). Denne Indvending rører ved
et Spørgsmaal, der har været meget lidet oppe i de hyppige Diskussioner om det
psykofysiske Problem. Hvad enten man gaar ind paa den specielle Teori, den
franske Filosof har opstillet, eller ikke, saa synes det klart, at det, jo mere
man rykker Problemet ind paa Livet, vil blive overordentligt vanskeligt at
finde de Led i de to Rækker af Fænomener, der kunne siges at være
„tilsvarende“. Hvor en kvalitativ Forskellighed er til Stede, kan man kun ad
Omveje begrunde, hvilke Led der svare til hinanden. Naar vi vide, hvorledes
Varmen virker paa et Stofs Farve og paa dets Form hver for sig, kunne vi
bestemme, hvilke Farveog Formforandringer der „svare til hverandre“. Men det er
jo netop en analog Viden, der ikke staar til vor Raadighed ved Drøftelsen af
det psykofysiske Problem. Vi ere da begrænsede til rent empiriske Slutninger.
Og saadanne ville her være vanskeligere at drage, end Bergson synes at antage.
Naar man nemlig lægger tilstrækkelig Vægt paa Kontinuiteten baade paa det
psykologiske og det fysiologiske Omraade, vil det vise sig vanskeligt at opløse
de to Rækker i Led, der fremtræde med en saadan Selvstændighed og
Ejendommelighed, at en virkelig Sammenligning kan finde Sted. Men hvilken
Hypotese vi saa lægge til Grund, maa vi være forberedte paa, at de Led i de to
Rækker, der kunne antages at svare til hverandre, ville frembyde
Forskelligheder, der ikke kunde sluttes ud fra det ene eller andet Led
\textit{for sig alene.} Det kunde derfor ogsaa meget vel være, at der til
„samme“ fysiologiske Tilstand kunde svare „forskellige“ psykiske Fænomener
(eller omvendt), ligesom et Sprog undertiden kun har ét Ord, hvor et andet
Sprog har to.
% --- p. 25 ---
\opage{25}Det er da Sammenhængen, der er det Afgørende, og den vil sikkert
endog --- baade paa det psykologiske og paa det fysiologiske Omraade --- være
saa afgørende, at det er kun med en ringe Grad af Tilnærmelse, man vil kunne
tale om „samme“ Fænomener eller Tilstande. „Parallelismen“ rammes altsaa ikke
af hin Indvending, naar vi stille tilstrækkeligt strenge Fordringer til
Godtgørelsen af, at Tilstandene virkeligt ere de samme.

Det er væsentligt som Arbejdshypotese, ikke som positiv Løsning, at
Identitetshypotesen (for hvilken „Parallelisme“ og lignende Udtryk ere uheldige
og misvisende Betegnelser) har sin Betydning. Jeg for mit Vedkommende har i
hvert Tilfælde stedse hævdet den som en „empirisk Formel“, der kan lede os i
vore Undersøgelser saaledes, at hverken Fysiologiens eller Psykologiens Ret
krænkes ved, at en Undersøgelse paa den ene eller anden Side stanses for
tidligt. Fysiologen kan fristes til at stanse for tidligt, naar han kan vente
paa et eller andet Punkt at støde paa „Sjælen“ som Aarsag til en Ændring i
Hjernens Tilstand\textsuperscript{21}), og Psykologen kan fristes til at stanse
for tidligt, naar han kan vente paa et eller andet Punkt at staa overfor et
psykisk Fænomen, der skulde have sin Aarsag i en „Nerveproces“ eller i „Nerveenergi“.

4. Det vil dog have sin Interesse, fra begge Sider, baade den psykologiske og
den fysiologiske, at nærme sig hinanden saa meget som muligt ved at gaa tilbage
til Grundbegreberne for de to Sider, under hvilke Tilværelsen her fremtræder i
vor Erfaring.

Fra den psykologiske Side vil \textit{Villiesbegrebet,} naar det tages i
videste Betydning som Begrebet om psykisk Aktivitet, staa som det fundamentale
Begreb. Det kunde synes at stride herimod, at flere psykologiske Forfattere
netop i den nyeste Tid søge at eliminere Begrebet Villie af Psykologien, ikke
fordi de negte, hvad den populære Sprogbrug kalder Villie, men fordi de mene,
at dette Begreb ikke betegner et saa fundamentalt Synspunkt som Erkendelse og
Følelse, idet de saakaldte Villiesfænomener skulle kunne føres tilbage til
Erkendelses- og Følelseselementer i ejendommelige Kombinationer. Man kan herved
% --- p. 26 ---
\opage{26}støtte sig til, at Villien som saadan, vor Aktivitet som bevidste
Væsener, ikke kan være Genstand for umiddelbar Selviagttagelse, saaledes som
Forestillinger og Følelser kunne det\textsuperscript{22}). Vi erfare Villiens
Motiver og Villiens Resultater, men ikke selve Villien --- ligesom vi paa den
materielle Naturs Omraade erfare Energiens Betingelser og Ytringer, men ikke
selve Energien. \textsc{Hume} har allerede paavist dette for al Kausalitets
Vedkommende, altsaa baade hvad psykisk og hvad materiel Kausalitet angaar.
Begrebet Villie dannes, ligesom Begrebet Energi, ved en Konstruktion. Men det
er en Konstruktion, vi ere nødte til at foretage. Og naar man ikke bestemmer
Begrebet psykisk Element saaledes, at det nødvendigvis skal være Noget, der kan
være Genstand for umiddelbar Selviagttagelse, men saaledes, at det betegner en
væsentlig og irreduktibel Egenskab ved Bevidsthedslivet, saa kan Villien meget
vel være et psykisk Element, og Villiens Begreb et psykologisk
Fundamentalbegreb. Grunden til, at vi ikke kunne gøre Villien til
Selviagttagelsens Objekt, saaledes som Fornemmelser, Forestillinger og Følelser
kunne blive det, kunde ligge i, at den strækker sig gennem alle
Bevidsthedslivets skiftende Tilstande og Former som den stadige Forudsætning.
Bevidstheden bestaar kun gennem et uafbrudt Arbejde for at sammenfatte de
enkelte Elementer til en Helhed. Et saadant Kombinations- og
Koncentrationsarbejde ytrer sig i den simpleste Sansefornemmelse saavel som i
enhver Forestillingsproces, enhver Følelse, enhver Drift og enhver Beslutning.
Paa ethvert Punkt ytrer der sig en Aktivitet, der er lige saa oprindelig en
Side ved Bevidsthedslivet, som de Elementer (Sider eller Egenskaber),
Iagttagelse og Analyse umiddelbart forefinde. Det forholder sig ikke saaledes,
at vi først have Fornemmelser, Forestillinger og Følelse, og at der saa ved
disses Kombination opstaar Noget, vi kunne kalde Villien. Uden en oprindelig
Kombination, en primær syntetisk Proces opstaa selve de Elementer, der bestemme
Villien i snevrere Betydning, ikke. I de specielle Villiesytringer
(Reflexhandling, Drift, Ønske, Forsæt, Beslutning) ytrer den primære
Koncentrationsevne sig paa særegen Maade under Indflydelse af visse bestemte
Erkendelses- og Følelseselementer\textsuperscript{23}). --- Der finder saaledes
en stadig Vexelvirkning Sted mellem Ak%
% --- p. 27 ---
\opage{27}tivitetselementet, det egentlige Villieselement, og de intellektuelle
og emotionelle Elementer. Vi møde atter her den ovenfor (1) omtalte Antinomi,
der gør sig gældende i alt Bevidsthedsliv, ja i alt Liv. Men hvad der her er
Hovedsagen, er, at vi allerede rent psykologisk kunne danne et Energibegreb,
idet der øves et psykisk Arbejde overalt, hvor et psykisk Fænomen optræder, da
dette stedse --- saavidt vi kunne spore --- forudsætter en Syntese. Det
psykiske Arbejde, Syntesen bestaar i, er des større, jo mere de enkelte
Elementer ere kvalitativt forskellige, og jo fjernere de ligge hverandre i Tiden.

Det kunde nu være fristende umiddelbart at gøre denne psykiske Energi til Et
med den Energi, der virker i Nervevævet. Men da ikke alle Nerveprocesser ere
forbundne med Bevidsthedsfænomener, maa der skelnes mellem bevidst og ubevidst
Nerveenergi, og Problemet stiller sig saaledes paany. Desuden vide vi intet
Nærmere om denne saakaldte Nerveenergi eller om dens Forhold til andre Former
for Energi, organiske og uorganiske. Ved --- som
\textsc{Ostwald}\textsuperscript{24}) o. A. --- at anse Gaaden løst ved selve
Opstillingen af Begrebet „Nerveenergi“ indfører man blot en qualitas occulta og
slaar sig til Ro ved den. Det naturvidenskabelige Begreb om Energi er, hvor man
saa møder det, stedse abstraheret ud af Fænomener med geometriske Egenskaber.
Naturvidenskaben kender kun Energi som Udtryk for Forholdet mellem Fænomener i
Rummet. Gaaden vilde altsaa kun være løst, hvis vi kunde danne et Energibegreb,
af hvilket det psykologiske og det naturvidenskabelige Energibegreb kunde
afledes som særlige Former; men et saadant Begreb mangle vi Midler til at
konstruere. I hvert Tilfælde vilde ethvert Forsøg i den Retning føre os ud over
det psykologiske Problems Omraade.

\begin{center}\rule{2cm}{0.4pt}\end{center}

% --- p. 28 ---
\opage{28}

\begin{center}II.\\[6pt]
\textbf{\Large Erkendelsesproblemet.}\end{center}

1. Ikke blot psykologisk Forstaaelse, men al Forstaaelse betinges ved Forholdet
mellem Kontinuitet og Diskontinuitet. Denne sidste betinger Fylden og
Mangfoldigheden af Forstaaelsens Indhold; hin betinger dets Sammenfatten og
dets Ordning. Forstaaelsen fremtræder under forskellige Hovedformer, til hvilke
der svarer forskellige Arter af Videnskaber.

Jeg forstaar, hvad Noget er, naar jeg kender det igen. Saaledes forstaar jeg,
eller véd jeg, hvem der kommer, naar jeg genkender den Kommende. Genkendelsen
beror paa, at der findes en Sammenhæng mellem nye og tidligere Erfaringer. I
Genkendelsesakten skydes alle mellemliggende Erfaringer til Side, og det nye
Fænomen identificeres umiddelbart eller middelbart, uvilkaarligt eller efter
nogen Overvejelse\textsuperscript{25}) med et tidligere forekommet Fænomen.
Kontrasten med det Nye virker her ogsaa med; Bekendthedskvaliteten bliver
stærkere udpræget, jo mere ubekendte Omgivelserne ere. --- Det er Genkendelse,
der gør \textit{Beskrivelse} og \textit{Klassifikation} mulige; dog virker
Forskelsopfattelse væsentligt med ved begge Operationer. Ved al Beskriven og
Klassificeren staa visse Typer fast, idetmindste foreløbigt fast, og de nye
Fænomener henføres til dem, enten som aldeles ens (dækningslige) eller som
lignende (kvalitetslige) eller som frembydende Analogi (Forholdslighed). De
positive Begreber, vi danne, udtrykke saadanne Typer. Hvor vi ingen Genkendelse
af nogen Grad kunne have, altsaa ingen Type
% --- p. 29 ---
\opage{29}opstille, samle vi foreløbigt Fænomenerne under et negativt Begreb,
d. v. s. vi give dem det fælles Mærke, at de ere forskellige fra de hidtil
bestemte Typer. De negative Begreber have i Klassifikationens Historie den
Betydning, at man ved deres Hjælp samler det hidtil Ubestemte i en fra det
Bestemte forskellig Gruppe, for saa at overlade til senere Forskning at finde
positive Bestemmelser indenfor den\textsuperscript{26}).

I Filosofiens Historie staar Platons Idelære som karakteristisk Udtryk for den
beskrivende og begrebsdannende Forskens Betydning. Selve Muligheden af
Genkendelse midt under Fænomenernes brogede Mangfoldighed opfylder Platon med
Begejstring, og den højeste aandelige Evne var for ham den, ved hvilken de
forskellige Typer, der ligge til Grund ved Genkendelsen, ordnes efter deres
indbyrdes Lighed og Forskel, saa at der dannes en Tankeverden, indenfor hvilken
man kan stige op og ned ad Trin, som danne en kontinuerlig Rækkefølge.

Naar der er dannet Begreber i forskellig Sammenhæng, kunne de forbindes i
Domme, og naar forskellige Domme have fælles Begreber, der kunne substitueres
for hverandre, kan der uddrages Slutninger. Ad \textit{Slutningens} Vej kan der
dannes Begreber og Domme, uden at man behøver at gaa tilbage til Erfaringen,
saaledes som de beskrivende og klassificerende Videnskaber paa ethvert Punkt
ere nødte til. Der gør sig her en Kontinuitet gældende af højere Art, idet
Tanker og Tankerækker af højst forskellig Oprindelse kunne bringes i
Forbindelse med hverandre. Medens Genkendelsen blot hviler paa Identitet,
hviler Slutningen paa Rationalitet: det er Forholdet mellem Grund og Følge, der
raader og gør en ny Art af Forstaaelse mulig. Jeg forstaar \textit{A} =
\textit{C,} naar jeg véd \textit{A} = \textit{B} og \textit{B} = \textit{C.}
Denne Art Forstaaelse vindes gennem de formale Videnskaber. Ejendommeligt for
dem er, at det tilsidst er ligegyldigt, hvorfra de første Begreber og Domme
skrive sig, naar de blot ere af den Beskaffenhed, at der kan bygges
Slutningsrækker paa dem. Saaledes bekymre de beskrivende Videnskaber sig heller
ikke om, hvorfra Fænomenerne komme, naar de blot kunne ordnes efter Ligheds- og
Forskelsforhold.

En tredie Art af Forstaaelse er karakteriseret ved, at man
% --- p. 30 ---
\opage{30}hverken blot stiger op fra Fænomener til Begreber, eller blot ved
Slutning danner nye Begrebsforbindelser af givne Begrebsforbindelser, men
udleder nye Fænomener af tidligere givne Fænomener. De nye Fænomener forstaas
ved at sættes i et analogt Forhold til tidligere Fænomener, som det, der
bestaar mellem Grund og Følge i en Slutning. Denne Art Forstaaelse er
karakteristisk for den exakte Realvidenskab (Naturvidenskab), saaledes som den
især har udviklet sig siden Renaissancetiden. Det er her
\textit{Aarsagsbegrebet,} som raader. Vi forbinde her Erfaringerne efter deres
uundgaaelige og lovbestemte Rækkefølge i Tiden.

Denne tredie Art af Forstaaelse har især beskæftiget den nyere Tids
Erkendelsesteori, efterat \textsc{Hume} og \textsc{Kant} havde indskærpet
Forskellen mellem logisk Slutning og real Aarsagsforklaring. Muligheden af at
danne og anvende Aarsagsbegrebet har i den nyere Tid vakt en lignende
Forundring, som hos Platon Muligheden af at danne Almenbegreber. Det er dog
ikke blot den intellektuelle Trang til at finde Sammenhæng mellem Erfaringerne,
der har ført til at stille Aarsagsbegrebet paa første Plads, men ogsaa Trangen
til skarpt at skelne mellem subjektive Forestillinger og objektiv Virkelighed,
idet Virkelighedskriteriet i Tvivlstilfælde stedse tilsidst er den faste,
ubrydelige Sammenhæng\textsuperscript{27}). Drømmenes og Tankernes Verden er
større end Virkelighedens (Eng ist die Welt, und das Gehirn ist weit), og
Menneskene have stedse mere følt Trangen til at faa deres Tanker saaledes
bestemte og forbundne, at de kunne staa som Udtryk for en Virkelighed, især da
de kun derved kunne vente at finde Midler til at opnaa deres Formaal.

Aarsagsbegrebet fremtræder i to Former, en foreløbig, elementær Form, ved
hvilken vi ofte nødes til at stanse, og en ideal Form, mod hvilken al Forsken
og alle Teorier stræbe hen. \textit{Det elementære Aarsagsbegreb} angiver blot
et uundgaaeligt Successionsforhold: naar Fænomenet \textit{A} er indtraadt,
indtræder derefter Fænomenet \textit{B} uundgaaeligt, og \textit{B} indtræder
kun, naar \textit{A} gaar forud. Det er ikke sagt at Aarsagsforholdet finder
Sted mellem \textit{A} og \textit{B} indbyrdes; det er muligt, at de begge ere
successivt fremtrædende Virkninger af en forudgaaende Aarsag. \textit{Det ideale
% --- p. 31 ---
\opage{31}Aarsagsbegreb} gaar et Skridt videre og ser i det Fænomen, der kaldes
Virkningen, en Fortsættelse af eller en Omsætning i ny Form af det Fænomen, der
kaldes Aarsagen. Det ideale Aarsagsbegreb gaar saaledes over i
Udviklingsbegrebet hvorfor det intet Under er, at dette Begreb ved Siden af
Aarsagsbegrebet spiller saa stor en Rolle i den nyere Tids Videnskab.

2. Alle tre Arter af Forstaaelse hviler paa visse Grundsætninger eller
Principer, som det uvilkaarlige Tankearbejde --- saaledes som det øves af den
praktiske Menneskeforstand, i Specialvidenskaberne, i den filosofiske
Spekulation og i den religiøse Tænkning --- ikke føler Trang til at opsøge, men
som Erkendelsesteorien fremdrager og prøver. Erkendelsesproblemet opstaar, naar
der spørges, hvorpaa \textit{Forstaaelsens Gyldighed} beror, og hvorvidt den
strækker sig. Det er Forholdet mellem Kontinuitet og Diskontinuitet, der atter
her stiller Problemet. For Platon opstod Problemet ved det irrationelle Forhold
mellem Begreb og Fænomen; i nyere Tid ved det irrationelle Forhold mellem
formal og real Videnskab. Kan Begrebet nogensinde være adækvat Udtryk for
Fænomenernes Mangfoldighed, og Loven for deres skiftende Forbindelser?

For \textsc{Platon} blev Konsekvensen den, at Begreberne („Ideerne“) maatte
være Mennesket bekendte fra en højere Tilværelse og nu dukkede op for Tanken i
deres hele Fuldkommenhed under Paavirkning af de ufuldkomne Efterligninger af
dem, Sanseverdenen frembyder. I nyere Tid har man ofte anset Tanken om Naturens
Lovmæssighed som en apriorisk Sandhed, en oprindelig Intuition, der i det
Højeste skyldte Erfaringen Anledningen til sin Fremkaldelse. Det er
Inkommensurabiliteten mellem Principerne (de logiske Principer og
Aarsagsprincipet) og Erfaringen, som har affødt denne spekulative
Erkendelsesteori. Denne Inkommensurabilitet har ogsaa affødt en Teori, man
kunde kalde den arbitrære Teori, fordi den stærkt betoner det Vilkaarlige i
Principernes første Opstilling, men som ikke mindre stærkt betoner, at
Principerne i deres rene Form ikke kunne vinde nogen Bekræftelse ved Erfaring,
ikke blive Resultater, stedse kun Postulater. Der kan --- siger
\textsc{Hobbes,} den mest udprægede Repræsentant for dette Standpunkt --- ikke
% --- p. 32 ---
\opage{32}være nogen Videnskab om selve Principerne for al Videnskab; de
opstilles ved konstruktiv Kunst (principia sunt artis sive constructionis, non
autem scientiæ et demonstrationis), og det er altsaa os selv, der frembringe
deres Sandhed (rationis prima principia vera esse facimus nosmet ipsi). Hos
\textsc{Fichte} opstilles al menneskelig Videns første og ubetingede
Grundsætning ved en fri Handling; vor Videnskab hviler ikke paa en „Thatsache“,
men paa en „Thathandlung“, nemlig den tænkende Bevidstheds Beslutning stedse at
forblive i Overensstemmelse med sig selv: Opstillingen af Identitetsprincipet.
I Lighed hermed antager \textsc{S. Kierkegaard} et Spring ved Principernes
Opstilling, og \textsc{Kromann} udleder Identitetsprincipet af
Selvopholdelsesdriften, der søger at opretholde Bevidsthedens Enhed; „af
Erfaringen er det paa ingen Maade fremgaaet“\textsuperscript{28}).

Baade den spekulative og den arbitrære Teori maa dog indrømme, at der maa være
bestemte empiriske Foranledninger, som fremkalde Intuitionerne eller
Postulaterne. Empirismen (der --- naar den tager Hensyn ikke blot til det
enkelte Individs, men til hele Slægtens Erfaringer --- gaar over til
Evolutionisme) lægger nu, hvor den udfolder sig i sin absolute Form,
Hovedvægten pas disse „Foranledninger“ og betragter dem som de fuldstændige
Aarsager. Den er klarest fremstillet af \textsc{Stuart Mill} og \textsc{Herbert
Spencer}. Den søger at vise, hvorledes almindelige Principer kunne have dannet
sig successivt i Bevidstheden under den stadige Paavirkning fra Omgivelserne,
og den hævder energisk, at intet Princip kan have nogen Gyldighed ud over den
Erfaringsbekræftelse (Verifikation), det kan vinde. Kun som Resultater altsaa
have Principerne Værdi.

Hvad Empirismen ikke kan forklare, er den Kendsgerning, at Principerne selv
fremkalde Erfaringer gennem de Spørgsmaal, de føre os til at stille. Heldigere
stillet er her en fjerde Teori, der i den nyeste Tid er bleven udviklet af
\textsc{Ernst Mach} og \textsc{Richard Avenarius}, og som man kunde kalde den
økonomiske. Den lader baade passiv Erfaring og aktiv Tankeudvikling komme til
deres Ret, idet den betragter Principerne som de „Begrebsreaktioner“, der ad
den korteste Vej gøre anskuelig Opfattelse og Oversigt mulige. Naar
Kontinuitetsprincipet spiller saa stor
% --- p. 33 ---
\opage{33}en Rolle, er det, fordi det er et afgjort økonomisk Princip. Der er
altsaa her kun Tale om praktisk Anvendelighed. Man danner de Begreber og
Principer, man behøver for at gøre sig til Herre over Fænomenerne. Enhver
praktisk og intellektuel Trang er tilfredsstillet, naar vore Tanker
fuldstændigt formaa at gengive de sanselige Kendsgerninger. Denne Gengivelse er
Fysikens Endemaal og Formaal, og Atomer, Kræfter og Love ere kun Midler til at
lette Gengivelsen. Deres Værdi gaar kun saa langt, som de hjælpe
os\textsuperscript{29}). I lignende Retning have \textsc{Maxwell} og
\textsc{Hertz} udtalt sig\textsuperscript{30}).

De fire Teorier, der henholdsvis opfatte Erkendelsens Principer som
Intuitioner, Postulater, Generalisationer og økonomiske Tankemidler, forudsætte
nu alle \textit{den analytiske eller regressive Erkendelsesteori,} som især
\textsc{Kant} har udviklet. Ti \textit{hvad} der skal være Genstand for
Intuition, eller for Postulat, eller for Generalisation, og \textit{hvad} der
svarer til Forskningens økonomiske Behov, det kan kun opdages ved \textit{at
slutte tilbage fra de Data, som Erfaringen frembyder, og finde de
Forudsætninger, deres Forstaaelse hviler paa.} En saadan Analyse maa ligge til
Grund for enhver Erkendelsesteori. Om denne Analyse er fuldført, vil man ikke
med fuld Sikkerhed kunne overbevise sig om. Videnskabens Historie viser, at der
her er et Arbejde, der atter og atter maa tages op. Snart mener man, at man
behøver flere, snart færre Principer end hidtil. En Garanti for, at man har
naat de absolut sidste Forudsætninger, vil man aldrig kunne faa. Hvad den
økonomiske Teori --- som jeg især vil holde mig til --- særligt betoner, er
dels, at der ikke maa opstilles flere Principer, end det Givne med streng
Nødvendighed kræver (det er dette, der menes med „Kritik af den rene
Erfaring“), --- dels at til forskellige Tider, i forskellige videnskabelige
Situationer kan der behøves forskellige Principer, saa at endog et Princip, der
en lang Tid har ført Forskningen videre, senere kan erkendes for ugyldigt, uden
at derved dets historiske Betydning miskendes. Den økonomiske Teori betoner
altsaa dels Sparsommelighedens Lov, dels den formaalsbestemte, nyttebestemte
Karakter af Principerne. Den skyldes især to Klasser af Motiver, for det Første
Ønsket om at reducere de Forudsætninger, der ikke selv kunne bevises,
% --- p. 34 ---
\opage{34}til det mindst Mulige, altsaa et antidogmatisk eller antimetafysisk
Motiv, --- for det Andet Erfaringer fra Videnskabernes Historie, der vise,
hvorledes Principer og Hypoteser kunne være berettigede og frugtbare i en vis
Periode, men senere maa afløses af andre, Noget, som særligt den nyeste Tids
Diskussion om den mekaniske Naturopfattelses Begrundelse og Gyldighed har
henledet Opmærksomheden paa.

Der er dog to Sider af Sagen, som den særlige Form, den analytiske
Erkendelsesteori har antaget ved at fremtræde som økonomisk Teori, ikke
tilstrækkeligt har betonet, og som blive af særlig Betydning, naar man ikke
betragter Erkendelsesproblemet isoleret, men ser det i dets Sammenhæng med de
andre filosofiske Problemer.

For det Første maa de Former, i hvilke den intellektuelle Trang
tilfredsstilles, hænge sammen med Bevidsthedens almindelige Natur.
\textit{Hvad} vi forstaa, og \textit{at} vi i det Hele forstaa Noget, beror
ikke blot paa Fænomenernes Beskaffenhed, men paa vor intellektuelle
Organisation, ligesaa vel som det afhænger af vort Synsorgans Beskaffenhed,
ikke blot af de ydre Genstande, hvilke Farver vi se. Der vil være en vis Type
for alle Principer og Hypoteser, som tilsidst viser tilbage til
Bevidsthedslivets inderste Natur. Og det vil her altid være Trangen til Enhed
og Kontinuitet, hvortil man kommer tilbage. Hvilke og hvormange Grundbegreber
(Kategorier) og Grundforudsætninger man maa stille op, det er et Problem, der
stedse maa tages op paany under Erkendelsens Kamp for at vinde fremad. Det var
en Illusion, naar Kant mente, at man én Gang for alle kunde give en Fortegnelse
over, hvad der i denne Henseende behøves; men deri tog den gamle Mester ikke
fejl, at det er \textit{Enheds- og Kontinuitetstrangen,} der ligger til Grund
ved alle Former, i hvilke Forstaaelse vindes eller menes at være vunden. Han
har selv vist, at alle hans Kategorier kunne føres tilbage til
Kontinuitetsbegrebet, og under alle Omskiftelser paa Principernes Omraade vil
det utvivlsomt være dette Begreb, der stadigt kommer igen. De logiske
Principer, Aarsagsprincipet og de fundamentale naturvidenskabelige
Grundsætninger føre alle tilbage til dette Begreb, der hænger saa nøje sammen
med Bevidsthedslivets Natur. En
% --- p. 35 ---
\opage{35}rent psykologisk Erkendelsesteori vil dog aldrig kunne gøre Fyldest;
ti fordi Trangen til Enhed og Kontinuitet --- den være nok saa væsentlig for
Bevidsthedslivet --- tilfredsstilles ved visse Principer, følger aldeles ikke,
at disse ere objektivt gyldige. Hin Trang kan skaffe sig Luft og
Tilfredsstillelse paa mange Maader og i mange Former, hvad Mytologiernes og
Spekulationernes Historie noksom godtgøre, --- Maader og Former, der i Regelen
ingenlunde tilfredsstille Økonomiens Fordringer, hverken i Henseende til
Sparsommelighed eller i Henseende til Anvendelighed. Her gælde
\textsc{Schiller's} Ord: Eng ist die Welt, und das Gehirn ist weit. Af al den
Tankerigdom, Menneskeaanden raader over, kan kun en snever og fast sluttet
Række af Tanker bruges, naar det gælder gyldig Forstaaelse. Hvad der er en
nødvendig Forudsætning, er kun, at de Forudsætninger, Forstaaelsen af det Givne
kræver, maa være psykologisk mulige, maa stemme med Bevidsthedslivets
almindelige Love, og kun være speciellere og nøjagtigere Udviklinger af, hvad
der ligger i disse. Der maa ske en Udvælgelse mellem de Tanker, der
uvilkaarligt trænge sig frem; men den medfører ingen Emancipation fra de
almindelige psykologiske Love. Hvad det gælder om, er, at der uddanner sig nye
Dispositioner af speciel Art, --- at der opstaar en intellektuel Habitus, der
stiller Spørgsmaalene og kritiserer Svarene paa strengere og bestemtere Vis,
end det uvilkaarlige Tankeløb føler Trang til. Ethvert omfattende Princip er
egentligt --- psykologisk set --- Udtryk for en saadan Habitus, der kan ligge
mere eller mindre dybt begrundet i Bevidsthedens Natur, det vil sige, snart
have et Instinkts Væsen, snart skyldes Øvelsens Indflydelse. De rent logiske
Principer nærme sig mest til det Instinktive. Trangen til Overensstemmelse med
sig selv, til Konsekvens i Forestillingsløbet kan ikke blot forklares ud fra
Sparsommeligheden og Hensigtsmæssigheden alene. Hvor streng Slutning fra de
hidtil foreliggende Erfaringer fører os til selvmodsigende Resultater, antage
vi langt snarere, at Erfaringerne ere ufuldstændige, end at Tilværelsen skulde
modsige sig selv. Hvad der i højeste Grad gælder for de rent logiske Principers
Vedkommende, gælder ogsaa for mere specielle Principer. Saaledes udtrykker
Aarsagsprincipet en Tilbøjelighed til, naar
% --- p. 36 ---
\opage{36}en Begivenhed er indtraadt, at se sig om efter andre Begivenheder, i
hvilke Betingelserne for deres Indtræden kunde findes. Der ytrer sig ogsaa her
en Trang til at bringe Enhed og Sammenhæng i Bevidsthedens Indhold. Jo mere
speciel og bestemt Tilfredsstillelsen af denne Trang skal være, des mere virker
det økonomiske Princip i sine to Former, som Sparsommeligheden, der gaar den
lige Vej til Maalet, og som Hensigtsmæssigheden, der følger Veje, som virkeligt
føre til Forstaaelse. Exempler herpaa haves i Keplers og Newtons Fordring om
veræ causæ, i Inertisætningen, i Energisætningen o. s. v. Hvad der rent
empirisk staar som en Hypotese, bliver erkendelsesteoretisk set et Princip, en
ledende Tanke, under hvis Førerskab Bevidstheden i Erfaringens Verden kan
tilfredsstille sin Trang til Enhed og Kontinuitet.

For det Andet maa Principer og fundamentale Hypoteser --- selv om de tilsidst
alle staa som \textit{Arbejds}hypoteseri den intellektuelle Økonomis Tjeneste
--- ikke opfattes som aldeles tilfældige eller vilkaarlige i Forhold til den
Tilværelse, som vi ved deres Hjælp ville forstaa. Selve Begrebet
Arbejdshypotese peger i dobbelt Retning: dels --- som allerede paavist ---
tilbage til den tænkende Bevidstheds Natur, da jo vor Bevidsthed ikke kan
udføre en Funktion, den være nok saa økonomisk, som er aldeles fremmed for dens
eget Væsen; dels over imod Tilværelsen, til hvilken de Fænomener, som skulle
forstaas, høre. Et Arbejdsredskab maa baade passe til Haanden, der skal bruge
det, og til Genstanden, der skal behandles. Tilværelsen maa da i sig selv
frembyde Sider, der svare til vor Erkendelses uundgaaelige Forudsætninger, selv
om disse nok saa meget betinges ved de indre og ydre Forhold, under hvilke
Erkendelsen arbejder, --- eller maaske netop derfor. Dette gælder for al gyldig
Erkendelse fra de mest elementære Former til de højeste. Det gælder allerede om
Sansekvaliteterne trods deres „Subjektivitet“, og det gælder om vor Tænknings
højeste Enhedsprinciper. \textsc{Kant} siger om Tendensen til at forudsætte en
Enhed bagved Naturfænomeners Forskelligheder\textsuperscript{31}): „Man kunde
maaske tro, at vi her havde et rent økonomisk Haandgreb af Fornuften, tjenende
til at spare Møje.... Men en saadan egoistisk Hensigt skelnes
% --- p. 37 ---
\opage{37}meget let fra den Ide, ifølge hvilken Enhver forudsætter, at denne
Fornuftenhed stemmer med Naturen selv, og at Fornuften her ikke tigger, men
befaler, skønt den ikke kan bestemme Grænserne for denne Enhed.“ Jeg føjer dog
hertil: selv hvor Fornuften befaler (eller spørger, forventer, anteciperer,
postulerer), nødes den (som alle Befalingsmænd) til at rette sine Befalinger
efter den Evne til at adlyde, der kan formodes at være til Stede. ---

Denne sidste Betragtning fører særligt til at muliggøre Sammenhængen mellem
Erkendelsesproblemet paa den ene Side og Tilværelsesproblemet paa den anden
Side. Det har saa meget des mere Betydning at fremdrage dette, som den
analytiske Erkendelsesteori --- ligesom den økonomiske, og paa sin Vis ogsaa
den arbitrære, --- fører til at sætte et nyt Sandhedsbegreb i Stedet for det
gængse, naive Begreb om Sandhed. Principernes Betydning er den at lede vort
Arbejde paa at vinde Forstaaelse. Deres \textit{Sandhed} bestaar i deres
\textit{Gyldighed,} og deres Gyldighed i deres \textit{Arbejdsværdi.} At et
Princip er sandt, betyder, at man kan arbejde med det, og det vil, naar Talen
er om Erkendelsens Principer, sige, at man ved dets Hjælp kan komme fremad i
Forstaaelse, --- i fast Ordning og Forbindelse af Fænomenerne. Sandhedens
Begreb er et \textit{dynamisk} Begreb, idet det udtrykker en bestemt Maade, paa
hvilken Tankeenergien anvendes, og et \textit{symbolsk} Begreb, idet det ikke
betyder Dækningslighed eller Kvalitetslighed med en absolut Genstand, men
Forholdslighed (Analogi) mellem Begivenhederne i Tilværelsen og de menneskelige
Tanker. Det gamle naive Sandhedsbegreb, ifølge hvilket en Erkendelse var sand,
naar den ligefrem afbildede eller afspejlede „Virkeligheden“, er uholdbart og
har været det lige fra det Tidspunkt af, da man begyndte at hævde
\textit{Sansekvaliteternes Subjektivitet.} At Sansekvaliteterne vare
subjektive, skulde jo ikke betyde, at de vare ugyldige, uanvendelige til at
orientere os i Verden. De staa stadigt som Tegn, Signaler, Symboler, hvis
indbyrdes Rækkefølge vi kunne tyde som Udtryk for en objektiv Række af
Begivenheder, skønt de ikke kunne bevises at være Afbilleder af disse. Paa
lignende Maade forholder det sig med de logiske Principer og de andre funda%
% --- p. 38 ---
\opage{38}mentale Forudsætninger for vor Erkendelse. \textit{Den kritiske
Filosofi} førte til det Resultat, som er en simpel Konsekvens af selve den
analytiske Metode, den anvendte, at Grundsætningernes Sandhed kun kan betyde,
at de gøre videnskabelig Erfaring mulig; de ere jo selv fundne ved Analyse som
en saadan Erfarings nødvendige Forudsætninger. En Sammenligning mellem vore
Tanker og en absolut Tingenes Verden er ikke mulig; vi kunne kun sammenligne
Tanker og Erfaringer indbyrdes. Denne Konsekvens saa Kant selv ikke saa klart
som nogle af hans Disciple (Maimon og Fries)\textsuperscript{31b}); Mesteren
selv var endnu hildet i Dogmatisme, som det ses af hans Lære om „Ding an sich“
som absolut bestaaende udenfor ethvert Subjekt. Men naar han betegner
Aarsagsloven som en „Erfaringsanalogi“ og dermed mener, at Begivenhederne i
Tiden forholde sig analogt som Grund og Følge i vor Tænkning, saa nærmer han
sig til at gøre den til en Arbejdshypotese. I den nyeste Tid er man --- som vi
snart skulle se --- \textit{fra naturvidenskabelig Side} mere tilbøjelig til at
anerkende det dynamiske og symbolske Sandhedsbegreb, end man var, saalænge den
mekaniske Naturopfattelse stod med en vis dogmatisk Karakter. --- Dette nye
Sandhedsbegreb, der er ved at trænge sig frem paa Videnskabens Omraade,
frembyder et Slægtskab med den Maade, paa hvilken --- som vi skulle se under
det fjerde Hovedproblem --- den religiøse Bevidsthed mere og mere stiller sig
overfor Dogmerne; det er ogsaa paa det religiøse Omraade mere og mere
Anvendeligheden og Arbejdsværdien, der spørges om. Det statiske
Sandhedsbegreb\textsuperscript{31c}) trænges overalt til Side for det dynamiske.

Selv om nu frugtbare Principer eller Arbejdshypoteser ere vundne, vil saa
Tilværelsen udtømmende kunne gengives ved Hjælp af dem, eller vil der stedse
vedblive at bestaa \textit{et irrationelt Forhold mellem de Principer, vor
Bevidsthed formaar at forme, og selve Tilværelsen, hvorfra vore Erfaringer
stamme?} --- Det vil vise sig, at der især under tre forskellige Former stedse
vil blive et Irrationelt tilbage: i Forholdet mellem Kvalitet og Kvantitet, i
den Betydning, Tidsforholdet har for Aarsagsbegrebet, og i Forholdet mellem
Subjekt og Objekt. Vi betragte disse tre Punkter hvert for sig.

% --- p. 39 ---
\opage{39}3. Ved Forsøgene paa Reduktion af alle givne Forskelligheder til
Identitet og Kontinuitet har især Bestræbelsen for at føre Artsforskelligheder
tilbage til Gradsforskelligheder været af stor Betydning. I Videnskaben om den
materielle Natur ytrer den sig som en Bestræbelse for at forklare al Forandring
som Bevægelse i Rummet. Bevægelse fra et Sted til et andet er den simpleste
Forandring; derfor vilde det være et stort Fremskridt i Klarhed, hvis den kunde
vises at være den Forandring, som under forskellige Former foregaar overalt i
Naturen, hvor den umiddelbare Erfaring viser os kvalitative Forandringer.
Allerede \textsc{Aristoteles} lærte, at den rumlige Bevægelse laa til Grund for
al anden Forandring, ogsaa for al Opstaaen og Forgaaen; men han ansaa det dog
ikke for muligt --- saaledes som Atomisterne havde villet --- at føre Alt i den
materielle Natur tilbage til den. Først med den moderne Naturvidenskabs
Opstaaen kommer denne Tanke frem for Alvor. \textsc{Galilei} staar som den
store Grundlægger af, hvad man (i snevreste Betydning) plejer at kalde
\textit{den mekaniske Naturopfattelse.} Foruden den Simpelhed, som derved
opnaas, --- og Galilei var sikker paa, at Naturen følger de simpleste Veje, ---
opnaas ogsaa den Fordel, at man kan operere med bestemte Kvantiteter. Galilei
sagde derfor: maal Alt, hvad der er maaleligt, og gør det maaleligt, som ikke
er det! Ved denne Reduktion bliver dels Enighed mulig, medens
Kvalitetsforskelle kun kunne skønnes, ikke maales, dels bliver nøjagtig
Verifikation mulig\textsuperscript{32}).

De Principer, hvorpaa den mekaniske Naturopfattelse hviler, betragtedes hos
dens vigtigste Repræsentanter som absolute Sandheder, fundamentale
Tilværelseslove; forsaavidt man endnu mente, at de behøvede en Begrundelse,
hentede man denne umiddelbart fra Guds Væsen og Villie. Heri vare Kartesianere
og Newtonianere enige. Og Materialisterne skilte sig kun fra dem ved at anse
den teologiske Begrundelse for overflødig og umulig. Gennem en Række af
storartede Opdagelser og Forklaringer har denne Naturopfattelse vist sin
Frugtbarhed. Spørgsmaalet er for os, om den fra et erkendelsesteoretisk
Synspunkt kan siges at betegne en absolut Afslutning.

Selv om vi nu tænke os, at Alt i den materielle Natur
% --- p. 40 ---
\opage{40}kunde forklares ved den mekaniske Naturopfattelses Principer (der er,
som vi snart skulle se, i den nyeste Tid Tvivl derom indenfor Naturforskernes
Kreds), saa er det for det Første klart, at Kvaliteterne ikke skaffes ud af
Verden, fordi de „reduceres“ til Kvantiteter eller skrives paa det sansende
Subjekts Regning. De blive staaende som umiddelbare Fakta, der maa konstateres
empirisk. Det kemiske Produkts Egenskaber kunne ikke udledes af Elementernes
Egenskaber, og naar en Art af fysisk Energi faar sit Ækvivalent i en anden Art
af fysisk Energi, saa har Ækvivalentet dog nye Egenskaber, der ikke kunne
udledes af den første Energiforms Egenskaber. Man kan dog ikke lade det
sansende Subjekt skabe disse Kvaliteter af Intet, i hvert Tilfælde faa vi
uløselige psykologiske Vanskeligheder i Stedet for de fysiske og kemiske
Vanskeligheder, vi have skudt til Side.

For det Andet er Udstrækning og Bevægelse tilsidst selv Kvaliteter i videste
Forstand, faktiske Egenskaber, der i og for sig ligesaavel kunde behøve en
Forklaring som de specielt saakaldte Sansekvaliteter. Dette er efter
\textsc{Berkeley's} og \textsc{Leibniz's} Tid ofte blevet gjort gjældende fra
filosofisk Side. Der er ingen Grund til at antage, at de kvantitative
Egenskaber skulde udtrykke Tingenes inderste Væsen. Grunden til, at Videnskaben
med Forkærlighed opsøger dem og holder sig til dem, er egentligt kun den, at
man ved deres Hjælp kan give nøjagtige Beskrivelser af de materielle Fænomener
og uddrage bestemte Slutninger af dem. Dette er i den nyeste Tid stærkt
fremhævet af Forskere som \textsc{Maxwell} og \textsc{Hertz}, for denne Sidstes
Vedkommende under udtrykkelig Tilslutning til den ovenfor omtalte økonomiske
Erkendelsesteori. „De exakte Videnskabers Fremgang“, siger Maxwell, „beror paa
Opdagelsen og Udviklingen af hensigtsmæssige (appropriate) og nøjagtige
Forestillinger, ved hvis Hjælp vi kunne danne os en aandelig Repræsentation af
Kendsgerningerne, der paa én Gang er tilstrækkeligt omfattende til at udtrykke
(stand for) hvert enkelt Tilfælde, og tilstrækkeligt nøjagtig til at yde
Sikkerhed for de Slutninger, vi ad matematisk Beregnings Vej drage fra dem“. Og
Hertz siger: „Vi danne os saadanne indre Symboler (Scheinbilder oder Symbole)
af de ydre Genstande, at de logisk nødvendige Følger
% --- p. 41 ---
\opage{41}af Billederne stedse igen ere Billeder af de afbildede Genstandes
naturnødvendige Følger.... Anden Overensstemmelse med Tingene fordrer Formaalet
med vore Forestillinger ikke“\textsuperscript{33}). --- Det dynamiske og
symbolske Sandhedsbegreb er her udtrykkeligt sat i Stedet for det naivt
dogmatiske, som den mekaniske Naturopfattelse tidligere havde hyldet. Hvad det
gælder om, er kun at finde en Kreds af Symboler, der kan anvendes med fuld
Konsekvens, og fra hvilken der kan drages Slutninger, som bekræftes ved nye
Erfaringer, der selv igen kunne udtrykkes ved den samme Kreds af Symboler. Men
den Mulighed vil her aldrig kunne afvises, at en anden Kreds af Symboler
ligesaa godt eller bedre vil kunne udtrykke de positive Erfaringer og egne sig
til Grundlag for exakte og verifikable Deduktioner. Der vil aldrig kunne føres
Bevis for, at nogen Kreds af Symboler skulde være den eneste rette og den
eneste nødvendige eller mulige.

De erkendelsesteoriske Refiexioner, til hvilke nyere Naturforskere saaledes ere
blevne førte, fremkaldtes ved den Vanskelighed, der frembød sig ved at bringe
de elektriske Fænomener ind under den mekaniske Naturopfattelse eller ogsaa
udlede dennes Principer af de for hine Fænomener gældende Love. Den sidste
Mulighed holder \textsc{Hertz} nærmest paa, medens \textsc{Boltzmann} nærmest
holder paa den første\textsuperscript{34}). Men af størst Interesse for
Erkendelsesteorien er \textsc{Maxwell's} Kritik af det gængse Materiebegreb,
ifølge hvilket man ved Materie forestiller sig en udstrakt Masse. Fejlen ved
dette Begreb er ifølge Maxwell den, at det leder til at tænke sig Materien som
noget Hvilende. Men det er Bevægelse, der gør os Hvilen forstaaelig, ikke
omvendt. Bevægelseslæren maa derfor gaa forud for Ligevægtslæren, Dynamiken for
Statiken. Det er fra Bevægelsen, vi vinde Begrebet Kraft eller Energi, ved hvis
Hjælp Ligevægten bliver os forstaaelig. Naar vi nu ved Materie forstaa det
Konstante, det, som bestaar under alle Forandringer, saa kan det kun være selve
Bevægelsens Lov, og „Materien“s Væsen bestaar altsaa i Bevægelse. Det er ifølge
Maxwell en Fordom at betragte Materien som udstrakt og Molekulerne som haarde
og solide: man maatte da spørge, hvad der holder Molekulens Dele sammen, og
% --- p. 42 ---
\opage{42}man vilde komme til Molekuler af anden Orden. Dog maa vi stadigt
operere baade med geometriske og med dynamiske Begreber, hvad enten vi anse de
sidste Elementer for udstrakte eller ikke. Det er kun Udstrækning som hvilende
Egenskab, Maxwell erklærer sig imod; selv ved en Linie, vi drage paa en Tavle,
er Bevægelsen, hvorved den bliver til, det Væsentlige\textsuperscript{35}).---
Det er et overordentligt vigtigt erkendelsesteoretisk Synspunkt, som her er
fremdraget: at hvilende Tilstande stedse indeholde Problemer, der kun kunne
løsesved, at Hvilen afløses af Bevægelse. Potentiel Energi forstaa vi kun ved
aktuel Energi, Evner og Anlæg kun ved, hvad der kommer ud af dem. Denne Lov
gælder overalt, paa det psykiske saavel som paa det fysiske Omraade. Der er en
halvt mystisk, halvt materialistisk Tilbøjelighed til at finde Fuldendelsen i
Kontemplation over --- eller Stirren paa --- et hvilende
„Noget“\textsuperscript{36}). Maxwell har ydet et vigtigt Bidrag til at faa
denne Tilbøjelighed udryddet. Men det store Spørgsmaal er, om Tanken om
Bevægelsens eller Aktivitetens Kontinuitet kan gennemføres paa alle Omraader.
Selv om det Dynamiske erkendelsesteoretisk set hævder sin Prioritet for det
Statiske, kan dette dog ikke skaffes bort; det stiller sig stedse paany frem
med sine Opgaver for Tanken, og Maxwell anerkender det selv, naar han erklærer
baade geometriske og dynamiske Begreber for nødvendige i Naturforklaringen. I
Modsætning til det Dynamiske betegner det Geometriske Samtidigheden. Paa den
materielle Naturs Omraade fremtræder det Samtidige i Rummets
Form\textsuperscript{37}); paa det psykiske Omraade antager
Samtidighedsforholdet ikke denne geometriske Form, men fremtræder alligevel som
noget „Statisk“, der stiller Forskningen nye Opgaver, naar den lykkeligt og vel
har fundet de psykiske Processers Love. Efterat Kontinuiteten i
Udviklingsprocesserne er paavist, gælder det at paavise, at der er Kontinuitet
mellem Processerne og de hvilende Tilstande. --- Vi møde altsaa her igen det
stadige Problem; det vilde fremstille sig, selv om Kvaliteterne vare fuldt
forklarede ved at repræsenteres af Kvantiteter, og selv om de fysiske
Grundbegreber, hvormed Videnskaben opererer, vare mere end den fuldkomneste og
hensigtsmæssigste Kreds af Symboler, det hidtil er lykkedes at forme og anvende.
% --- p. 43 ---
\opage{43}Muligheden af et irrationelt Forhold mellem Tilværelsen og vor
Erkendelse kan derfor ikke afvises.

4. Til et lignende Resultat vil Undersøgelsen af Forholdet mellem \textit{det
elementære og det ideale Aarsagsbegreb} føre os. I sin elementære Form angiver
Aarsagsbegrebet, som vi have set, kun det Forhold mellem to Begivenheder, at
naar den ene er indtraadt, indtræder uundgaaeligt den anden. Men fra denne
uundgaaelige Rækkefølge skrider Forskningen gennem Iagttagelser, Forsøg og
Hypoteser saavidt muligt frem til Paavisning af en Kontinuitet i
Begivenhedernes Række, ved hvilken Forskellene mellem de enkelte Led i denne
blive forsvindende, saa at tilsidst endog selve Tidsforskellen mellem
Begivenhederne reduceres til det mindst Mulige. Fra ydre, skønt uundgaaelig
Rækkefølge arbejder Tanken sig gennem Kontinuiteten hen imod en fuldstændig
Identitet. Medens Tidsforholdet spiller en væsentlig Rolle ved det elementære
Aarsagsbegreb, er det saa godt som elimineret i det ideale Aarsagsbegreb. Hvis
denne Proces kunde gennemføres helt, vilde vi komme til det paradoxe Resultat,
at den fuldkomne Aarsagsforklaring vilde være en Ophævelse af selve
Aarsagsbegrebet: Aarsagsforholdet er jo kun forskelligt fra det rent logiske
Forhold mellem Grund og Følge ved Tidsforskellen mellem Forholdets Led.

Det 17. Aarhundredes Filosofi (\textsc{Spinoza} og \textsc{Leibniz}) skelnede
ikke mellem Aarsag (causa) og Grund (ratio); Aarsagsforholdet mellem to
Fænomener var for dem det samme som Forholdet mellem Præmisser og Konklusion i
en Slutning, et blot Identitetsforhold. Det stod for dem som selvfølgeligt, at
der ikke kunde være Mere i Virkningen, end der var i Aarsagen; Tidsforholdet
mellem Leddene var ligegyldigt, og Tiden hørte i det Hele kun hjemme i den
sanselige Forestillings (Imaginationens) dunkle Sfære. Da \textsc{Hume} rejste
Aarsagsproblemet, skete det netop ved at betone, at der intet Bevis kan føres
for Berettigelsen af at slutte fra Fortid til Fremtid. Naar der henvises til
Fortidens Erfaringer, siger Hume, da kunne disse kun vise, at en Ting engang, i
et enkelt Øjeblik, har været i Besiddelse af en vis Evne, ikke at den stadigt
besidder den og vil besidde den\textsuperscript{38}). \textsc{Kant} mente vel
at kunne føre et rationelt Bevis for
% --- p. 44 ---
\opage{44}Aarsagssætningens Gyldighed, men han skelnede mellem Aarsagsforholdet
og Forholdet mellem Grund og Følge, saaledes at der kun var en Analogi, ikke en
Identitet, imellem dem: Aarsagsloven betød for ham, at Begivenheder følge efter
hinanden i Tiden paa en Maade, der er analog med den Maade, paa hvilken
Konklusionen udspringer af Præmisserne. Derfor giver Aarsagssætningen os en
Regel, ved hvis Hjælp vi kunne opnaa Enhed i vore Erfaringer\textsuperscript{39}).

Fra to forskellige Standpunkter har man i den nyeste Tid forsøgt at eliminere
Tidsforholdet og derved ophæve ikke blot det elementære Aarsagsbegreb, men
tilsidst Aarsagsbegrebet i det Hele.

Fra et spekulativt Standpunkt har man hævdet, at Tidsforholdet altid tyder paa
en Ufuldkommenhed, paa en Ufuldstændighed i Udviklingen. Saalænge Tidsforholdet
er bestemmende for vor Opfattelse af Tilværelsen, er vor Opfattelse derfor
stedse ufuldkommen, og Tidsforholdets Realitet er uforenelig med Ideen om en
fuldkommen Erkendelse, om en absolut Sandhed. Vor Erkendelse arbejder sig jo
ogsaa mere og mere bort fra Tidsforholdet; jo klarere vi forstaa Noget, des
mindre Rolle spiller Tidsforskellen, --- des mere gaar den reale Viden op i den
formale, Aarsagen i Grunden. Hvad vi begynde med at kalde Aarsag og Virkning,
kommer under Erkendelsens Fremskriden til at staa som Led i en Helhed, Led, der
staa i et lovmæssigt Forhold til hinanden, et rationelt Forhold, for hvilket
Tidsfølgen er uvæsentlig. At vi først iagttage det ene Led og derefter det
andet Led, bar ingen Betydning: i vor Erkendelse af Loven staar hele Processen
for os som en Enhed. Det er ikke Tidsforholdet, men Enheden bag Tidsforholdet,
der binder det, vi kalde Aarsag, til det, vi kalde Virkning. Desuden er der
heller intet Mellemrum mellem Slutningen af den Begivenhed, vi kalde Aarsag, og
Begyndelsen af den Begivenhed, vi kalde Virkning. Det er de engelske Filosofer
\textsc{Francis Bradley} og \textsc{Bernard Bosanquet}\textsuperscript{40}),
som have gjort denne Betragtning gældende, i hvilken en spekulativ Interesse
for Helhedsopfattelse og for afsluttende Begreber ligger til Grund.

Fra et helt andet Standpunkt, der kan betegnes som „den
% --- p. 45 ---
\opage{45}rene Erfarings“ Standpunkt, og hvis vigtigste Repræsentanter ere
\textsc{Richard Avenarius} og \textsc{Ernst Mach}, arbejdes der i en lignende
Retning, men af andre Motiver. Interessen gaar her ud paa at udskille alle de
Biforestillinger og Hjælpebegreber, med hvilke vi uvilkaarligt eller af
metodiske Grunde supplere selve Erfaringen, det umiddelbart Givne. Erkendelsens
Fremskriden bestaar i en Reduktion af Forskellighederne (til et „heterotisk
Minimum“) og i Tilnærmelse til en ren Beskrivelse af en kontinuerlig Proces.
Under denne Fremskriden begrænses Erkendelsesindholdet mere og mere til
deskriptive Udsagn, saavidt muligt med analytiske Overgange, og
Forskellighederne reduceres fra kvalitative til kvantitative, saavidt muligt
med stadig Paavisning af Ækvivalensforhold. Det er ikke længere den strenge
Videnskabs Opgave at „forklare“; hvem der vil have „Forklaringer“ henvises til
Mytologien og Metafysiken; Videnskabens Maal er at give en nøjagtig, metodisk
Beskrivelse af alle Forhold og Overgange. Hvad særligt Aarsagsforholdet angaar,
lægger man dels Vægten paa, at enhver Anvendelse af Begreberne Aarsag og
Virkning vilkaarligt drager to Elementer ud af hele den Sammenhæng, i hvilken
de forekomme, og stiller dem i Modsætning til hinanden, dels derpaa, at
Tidsforholdet meget vel kan vendes om, saasnart man har paavist et
Ækvivalensforhold, og saasnart man ser bort fra Forandringernes
Retning\textsuperscript{41}).

De erkendelsesteoretiske Anskuelser, der saaledes med forskellig Motivering
ville eliminere det elementære Aarsagsbegreb og negte Tidsforholdets
principielle Betydning, kunne karakteriseres ved, at de sætte et Erkendelsens
Ideal i Stedet for den virkelige Erkendelse, vi til nogen given Tid kunne
opnaa. Enhver bestemt Undersøgelse maa begynde paa et bestemt Punkt, som ligger
der, hvor Problemet stiller sig. Og Problemet stilles, naar to Led i
Begivenhedernes Række tildrage sig Opmærksomheden og vække Formodninger om en
indre Sammenhæng mellem dem, eller naar et enkelt Led træder frem i sin
Særegenhed eller Pludselighed og derved vækker Trangen til at søge Mellemled,
gennem hvilke det kan bringes i Sammenhæng med hele den øvrige Række. Det er
forsaavidt et til%
% --- p. 46 ---
\opage{46}fældigt, vilkaarligt Udgangspunkt, der gribes; men det vil efter
Sagens Natur ikke kunne være anderledes. Vor Erkendelse udvikler sig historisk,
fordi den menneskelige Opmærksomhed kun vækkes under visse bestemte
Betingelser. Hvis den uden Forskel til alle Tider lige meget rettedes mod alle
Led i Begivenhedernes Række, vilde der ingen Erkendelse opstaa. Naturligvis er
det af Vigtighed, at man er sig det Tilfældige eller Subjektive i
Udgangspunktet og Udsnittet bevidst, og hvis Erkendelsesarbejdet har Fremgang,
kommer man jo ud over det, idet Udsnittet netop gennem Forstaaelsen indarbejdes
i en stor kontinuerlig Sammenhæng.

Den samme historiske Karakter af vor Erkendelse bliver klar, naar vi besinde os
paa, at vi stedse staa mellem Fortidens Erfaringer og Fremtidens Muligheder.
\textit{I hvert enkelt Øjeblik} maa vi skelne mellem, hvad der foreligger
givet, og hvad vi vente; hint er Aarsagen, dette Virkningen; og Forventningen
kan aldrig blive Mere end en Hypotese. Det enkelte Øjeblik, i hvilket der til
den ene Side gælder et „ikke længere“, til den anden Side et „endnu ikke“,
stiller os Problemet i hele dets Spænding, en Spænding, som kun Vanens Magt
forringer. Ved nye Problemer vil Spændingen ikke udeblive --- og saalænge
Erkendelsen gaar fremad, vil den søge og finde nye
Problemer\textsuperscript{42}). Den fulde Sammenhæng foreligger stedse først
bagefter, og indtil den foreligger, staar Antagelsen af den som en Hypotese. Vi
ville derfor aldrig kunne undvære Begreber som Kraft, Energi, Aarsag eller
Mulighed, der med forskellige Nuancer udtrykke et og samme Forhold, nemlig de
senere Tilstandes Afhængighed af de foregaaende Tilstande; denne Afhængighed
adskiller sig fra rent logiske og matematiske Afhængighedsforhold derved, at
det paa én Gang er et Tidsforhold og et rationelt Forhold, idet den følgende
Tilstand baade følger efter og af den foregaaende. Selv om Kontinuiteten
paavises ved Hjælp af nok saa mange Mellemled og Nuancer, bliver dette Forhold
ved at gælde for enhver lille Overgang mellem to Led eller Nuancer.

Og herved ændres Intet, selv om der paavises Ækvivalens mellem de hinanden
afløsende Tilstande. Humes Problem er ingenlunde, som det undertiden er sagt,
blevet løst ved Op%
% --- p. 47 ---
\opage{47}dagelsen af Energiens Bestaaen. Et Ækvivalensforhold udelukker ikke
en kvalitativ Forskellighed, men forudsætter den netop: jeg har ingen Grund til
at opstille Dommen \textit{A = B,} naar \textit{A} og \textit{B} ikke først
have staaet for mig som forskellige, før jeg ved nærmere Undersøgelse fandt, at
de kunde sættes i Stedet for hinanden. Naar et saadant 1dentitetsforhold %
% PRINTED AS IS: the figure 1 for I
 er fundet, er det ligegyldigt, om vi gaa fra \textit{A} til \textit{B,} eller
fra \textit{B} til \textit{A}; men den Gang, da det skulde findes, begyndte vi
med \textit{A} eller \textit{B} og kom ad Undersøgelsens Vej over til det andet
Begreb. Derfor maa vi ved en Dom skelne mellem den psykologiske Proces, ved
hvilken Dommen bliver til, og i hvilken der gør sig en bestemt Forskel gældende
mellem Udgangsforestilling (Subjekt) og Slutningsforestilling (Prædikat), og
den færdige Dom, der kan formuleres som et Identitetsforhold, ved hvilket
Forskellen mellem Subjekt og Prædikat bliver ligegyldig\textsuperscript{43}).
Og ligesaalidt som i Tankebevægelsen er ved den ydre Skeen Leddenes Orden
ligegyldig. Har jeg erkendt, at der er et Ækvivalensforhold mellem Varme og
Bevægelse, er det rent abstrakt set ligegyldigt, om jeg gaar fra Varme til
Bevægelse eller fra Bevægelse til Varme; men „Forandringernes Retning“ er i
Virkeligheden et Spørgsmaal paa Liv og Død: Tilværelsen bliver faktisk en
anden, alt efter som Forandringerne overvejende foregaa i den ene eller den
anden Retning. Tiden kan da i Virkeligheden ikke vendes om, og derfor kan
Paavisningen af Ækvivalensforhold ikke danne Afslutningen af vor Erkendelse. Vi
maa tillige have Besked om Forandringernes (Omsætningernes) faktiske Retning,
og den vindes kun ved stedse ny Erfaring.

Dertil kommer endnu, at ethvert Ækvivalensforhold --- saavel som ethvert
Aarsagsforhold i det Hele --- kun træder i Kraft, naar visse Betingelser,
særligt Udløsningsbetingelser, ere til Stede. Med Hensyn til disse Betingelsers
Tilstedeværelse opstaar da et nyt Problem: kan man, ogsaa hvad dem angaar,
paavise, at Forholdet til deres Aarsag kan fremstilles som et
Ækvivalensforhold? Vi komme her ind i en uendelig Række, hvor Svarene (de
bestemte, konkrete Svar) høre op, længe før Spørgsmaalene høre op. Ogsaa her
have vi en rent logisk Analogi; ti ethvert absolut Identitetsforhold (\textit{A
= B})mellem to
% --- p. 48 ---
\opage{48}Begreber hviler ved nærmere Betragtning paa bestemte Betingelser (saa
at det --- efter \textsc{Jevons}' Formulering --- maa udtrykkes ved \textit{AC
= BC})\textsuperscript{44}); det bliver da et nyt Spørgsmaal, hvorledes det
forholder sig med disse Betingelser (\textit{C}), og saaledes gaar Tanken
stedse videre. ---

Der er saaledes ingen Udsigt til at slippe for det historiske Element i vor
Erkendelse. Maalestokken for Erkendelsens Udvikling bestaar i den Udstrækning,
i hvilken det elementære Aarsagsbegreb (uundgaaelig Rækkefølge) kan anvendes i
Sammenligning med den blotte Konstateren af faktiske simultane eller successive
Forskelligheder, og dernæst i den Udstrækning, i hvilken det elementære
Aarsagsbegreb kan afløses af det ideale Aarsagsbegreb (Ækvivalens eller
Identitet). Men Erkendelsens Proces vil til enhver Tid bestaa i en Opstigen
gennem de her antydede tre Stadier, --- en Opstigningsproces, der stedse paany
maa gentages ud fra nye Udgangspunkter. Og det er Tidsforholdets Realitet, der
betinger denne Nødvendighed. Spekulative Opfattelser --- hvad enten de
fremtræde i idealistisk eller i realistisk Form --- have derfor stedse en
Tilbøjelighed til at skyde Tidsforholdet til Side eller betragte det som blot
„empirisk“, om ikke som illusorisk\textsuperscript{45}). Skulde det være en
Illusion, saa er det i hvert Tilfælde en Illusion i anden Potens, naar man
mener at kunne frigøre sig for denne Illusion. Tilværelsen gaar for os aldrig
op i Tanken.

5. Endnu fra en tredie Side viser Erkendelsesproblemet sig i hele sin Skarphed,
samtidigt med at en stadig Udvikling af Erkendelsen viser sig mulig. I enhver
Erkendelsesakt kan der skelnes mellem \textit{et subjektivt og et objektivt
Element,} mellem det Erkendende og det, som erkendes, --- men begge Elementer
ere kun givne i Forhold til hinanden, skønt de indenfor dette Forhold kunne
gøre sig gældende i forskellig Grad.

Hvad er nu subjektivt, og hvad er objektivt i vor Erkendelse? --- Allerede af
det Foregaaende vil det ses, at dette Spørgsmaal besvares paa forskellig Maade.
Der gør sig paa Videnskabens Omraade en Tendens gældende til at skrive alle
Kvalitetsforskelle, Alt hvad der bryder Kontinuiteten, og tilsidst maaske endog
Alt hvad der bryder Identiteten, paa Sub%
% --- p. 49 ---
\opage{49}jektets Regning. Sansekvaliteterne, Rums- og Tidsforskellene skulle
da blot være subjektive. Det Berettigede i denne Synsmaade beror paa, at de
Forskelligheder, vi statuere i kvalitativ, extensiv, intensiv og protensiv
(tidslig) Henseende, bero paa vore psykiske Dispositioner. Vor Fornemmelse er
en Skelnen, hvis Resultat beror paa det fornemmende Subjekts Organisation og
Forhistorie. De fundne Forskelligheder gælde altsaa kun i Relation til det
erkendende Subjekts Stade i hine forskellige Henseender. Dertil kommer, at vor
Opmærksomhed er diskontinuerlig, bevæger sig i Spring til og fra sin Genstand;
ogsaa derved kommer der Forskelligheder og Afbrydelser frem, der ikke kunne
skrives paa Objektets Regning.

Den dogmatiske og spekulative Retning i Filosofi og Naturvidenskab har været
tilbøjelig til at gaa denne Vej. Emancipationen fra det blot Subjektive blev
her tillige en Tilnærmelse til Kontinuiteten og Identiteten som Sandhedens
Væsen og Norm.

Den kritiske Filosofi fremhæver i Modsætning hertil, at de kvalitative,
extensive, intensive og protensive Forskelligheder udgøre det for vor
Erkendelse givne Materiale og stille os Opgaverne for vor Forsken. Ligesom vor
Personlighed arbejder paa at samle sine sporadisk optrædende Elementer, at
harmonisere de stridende Tendenser og frigøre sig for det Dunkle og
Selmodsigende, saaledes arbejder vor Forstaaelse paa at forvandle de givne, for
os faktiske Forskelligheder til Stadier i en og samme kontinuerlig
Udviklingsproces eller til Former for et og samme Indhold. Trangen til
Kontinuitet og Identitet ligger dybt i den menneskelige Bevidsthed, og denne
søger derfor at genfinde dem i Erkendelsens givne Indhold. Tilmed kan
Bevidstheden ikke selv frembringe de Forskelligheder, der danne Materialet for
dens Virken, saa meget end den Form og den Grad, i hvilke de fremtræde i
Bevidstheden, bestemmes ved dennes egne uvilkaarligt virkende Betingelser. ---
I den nyeste Tids Diskussioner om Naturvidenskabens erkendelsesteoretiske
Grundlag viser der sig ogsaa en Tendens til overvejende at udlede alle
Enhedssynspunkter af Erkendelsessubjektets Økonomi.

I Virkeligheden have vi aldrig nogensinde det rene Subjekt
% --- p. 50 ---
\opage{50}med dets Former givet i Modsætning til et rent Objekt, eller rettere
til en „Ting i sig“, fra hvilken Erkendelsesindholdets Mangfoldighed skulde
skrive sig. \textsc{Kant} sluttede her for tidligt af, idet han mente, at
Erkendelsens subjektive Former én Gang for alle skulde kunne bestemmes, saa at
„Stoffet“ i Erkendelsen maatte være at udlede af „das Ding an sich“. I den
speciellere Udvikling af sin Erkendelsesteori kunde han dog ikke undgaa at
komme til at spørge om, hvorfra Formerne skrev sig, og til at antyde, at ogsaa
disse (der jo ligesaa vel som Stoffet ere faktisk givne, idet de jo maa
udfindes ved psykologisk Analyse) tilsidst ere bestemte ved „Ding an sich“, og
da „Formerne“ betegne det Konstante i vor Erkendelse, bliver den sidste
Forudsætning for Kants Filosofi egentligt Noget, han ikke kan give Begrundelse
af, nemlig at „Ding an sich“ virker paa konstant Maade, da Formerne ellers ikke
kunne bestaa eller i hvert Tilfælde ikke anvendes\textsuperscript{46}).
Saaledes kommer Kant til at svinge mellem begge de to ovenfor beskrevne
Anskuelser, skønt det aabenbart var hans Mening at holde sig til den sidste.
Kontinuiteten var for ham Kategoriernes fælles Grundform, og Syntesen var for
ham Grundloven for den erkendende Bevidstheds Virken.

Problemet er mere indviklet, end Kant saa. Naar vi skelne mellem Subjekt og
Objekt i vor Erkendelse, er det i Virkeligheden et objektivt bestemt Subjekt
(S\textsubscript{0}), vi stille i Modsætning til et subjektivt bestemt Objekt
(O\textsubscript{s}). %
% PRINTED AS IS: subscript s for 0 (also p. 51)
 De Egenskaber eller „Former“, vi tillægge Subjektet, kunne ikke forklares ud
fra selve Begrebet Subjekt (det rene S); de staa som Fakta, ligesaa vel som
alle andre Egenskaber, vor Erkendelse faar med at gøre. Og de Egenskaber eller
Bestemmelser, vi tillægge Objektet, tilkomme dette stedse kun i Relation til et
Subjekt og, nærmere bestemt, til et Subjekt af en vis speciel Beskaffenhed.
Derfor begynder Problemet stedse paany: hvorfra faar Subjektet sine objektive
(faktiske) Egenskaber, og hvilket Forhold er der mellem Objektets subjektive
Bestemmelser (Kvaliteter o. s. v.) og dets Væsen, som dette vilde fremtræde
overfor Subjekter af anden Art end vi?

Atter her møde vi det Irrationelle, og her se vi maaske
% --- p. 51 ---
\opage{51}klarest, hvor uudtømmelig Tilværelsen er i Forhold til vor
Erkendelse. Det Berettigede i Kants Opstilling af Begrebet „Ding an sich“ laa
i, at der behøves et Grænsebegreb til at udtrykke det irrationelle Forhold
mellem, hvad han kaldte vor Erkendelses „Form“ og dens „Stof“. Hvis vi ville
beholde Begrebet „Ting i sig“, kunne vi anvende det i Kants Aand og dog undgaa
at falde i de Modsigelser, med hvilke det i hans Filosofi er behæftet, naar vi
bruge det til at udtrykke det Faktum, at Forskellen mellem Subjekt og Objekt
stedse træder i Kraft paany, hvor ofte vi end have fundet en objektiv
Forklaring af Subjektets Ejendommeligheder (af S\textsubscript{1} ved
O\textsubscript{1}) og en subjektiv Forklaring af Objektets Ejendommeligheder
(af O\textsubscript{1} ved S\textsubscript{2}). Det Irrationelle viser sig i,
at en stadig Rækkedannelse (af Typen: S\textsubscript{1} \{O\textsubscript{1}
\{S\textsubscript{2} \{O\textsubscript{2}....) er mulig og nødvendig. Tanken
maa stedse paany i Gang med at finde Prædikater til Bestemmelse af Tilværelsen,
fordi Kilden, der gør Tankestrømmen mulig, er uudtømmelig. „Tingen i sig“ er
den dunkle, underforstaaede Udgangsforestilling, der stedse fremtræder paa nye
Maader og kræver nye Bestemmelser\textsuperscript{47}). Maaske vilde det rette
Symbol for vor Erkendelses Forhold til Tilværelsen ikke være et irrationelt
Tal, men et imaginært Tal, idet Tilværelsen kunde have Attributer, der slet
ikke kunne naas eller bestemmes ved de Dimensioner, i hvilke vor Tanke kan
bevæge sig. Muligheden heraf kan i hvert Tilfælde ikke modbevises. Og heller
ikke den Mulighed, at Tilværelsen kun er rationel i meget ringe Omfang, og at
den engang kunde vende en anden Side ud imod os, der ikke muliggjorde Bygningen
af sammenhængende og anvendelige Tankerækker; der vilde da, som jeg et andet
Sted har udtrykt det\textsuperscript{48}), begynde en logisk Istid for os.
Forholdet mellem Subjekt og Objekt vilde slet ikke træde i Kraft; hverken et
S\textsubscript{0} eller O\textsubscript{s} %
% PRINTED AS IS: subscript s for 0 (as on p. 50)
 kunde bestaa.

Naar jeg i en tidligere Sammenhæng anvendte Schillers Ord: Eng ist die Welt,
und das Gehirn ist weit, saa viser det sig nu, at Sætningen kan vendes om:
Gross ist die Welt, und das Gehirn ist klein! Erkendelsen er, saa rig og
kraftig den end monne være, stedse dog kun en Del af Tilværelsen, og
Erkendelsesproblemet vilde kun kunne løses, hvis Tilværelsen
% --- p. 52 ---
\opage{52}som Helhed (forsaavidt den kan tænkes som afsluttet Helhed!) kunde
udtrykkes ved Hjælp af en enkelt Del af den. Symbolsk maa Udtrykket i hvert
Tilfælde stedse blive, Vor Erkendelse giver os, selv hvor den naar højest, kun
et Udsnit af noget mere Omfattende. Af alle Tankens Muligheder er kun en eneste
repræsenteret i den Virkelighed, vi erkende; men den Virkelighed, vi erkende,
er selv stedse kun en Del af et større Hele --- og Forholdet mellem Del og
Helhed kunne vi her ikke bestemme. Vi kunne intet udtømmende Virkelighedsbegreb
danne.

\begin{center}\rule{2cm}{0.4pt}\end{center}

% --- p. 53 ---
\opage{53}

\begin{center}III.\\[6pt]
\textbf{\Large Tilværelsesproblemet.}\end{center}

1. Baade Bevidsthedsproblemet og Erkendelsesproblemet vise ud over sig selv:
Bevidsthedslivet er en Del af Tilværelsen, og det er Erkendelsens Opgave at
forstaa Tilværelsen. Det er da naturligt at spørge, hvilken Plads
Bevidsthedslivet indtager i Tilværelsen, og hvilket Billede Erkendelsen formaar
at give os af Tilværelsen. Det Problem, vi saaledes føres over til, kunne vi
kalde det \textit{kosmologiske,} da det er Problemet om, hvorvidt en
afsluttende Verdensanskuelse er mulig: \textit{Kosmos} betyder Tilværelsen,
betragtet som et Hele, og \textit{Logos} betyder Lære, Anskuelse. Ordet
Metafysik vil det ogsaa være berettiget at bruge her, skønt det er blevet
anvendt i mange forskellige Betydninger og derfor er mindre tydeligt.

Hovedvanskeligheden ved dette Problem vil være den, at de kosmologiske
Principer, de Principer, der skulde ligge til Grund for Verdensanskuelsen, ikke
uden videre kunne hentes fra noget specielt Erfaringsomraade, da det jo netop
skulde gælde om at forbinde alle Erfaringsomraader til et Hele og give en
positiv Karakteristik af dette.

Et Forsøg paa at behandle Tilværelsesproblemet vil baade have en formal og en
real Karakter: dels maa der spørges om, hvilke Fordringer der maa stilles til
det Begreb, der skal kunne omfatte Tilværelsen som Helhed, og om det er muligt
at fyldestgøre disse Fordringer, dels maa der spørges om, hvilke positive
Karaktermærker der kan angives om en saadan Helhed, hvis vi kunne danne et
Begreb om den.

% --- p. 54 ---
\opage{54}De \textit{formale} Motiver til Spekulation over Tilværelsens Natur
ligge dybt i Bevidsthedens og Erkendelsens Beskaffenhed. De hænge sammen med
Trangen til Kontinuitet, den Trang, i hvilken Personligheden og Videnskaben
mødes. Selve Tænkningens Væsen ytrer sig paa alle Trin og i alle Former som en
Sammenfatten, en Syntese, og det er da ikke at vente, at Tænkningen godvilligt
skulde stanse med sit Forsøg paa at sammenarbejde det sporadisk Givne. Speciel
og praktisk Betydning faar denne Stræben derved, at en fast og kontinuerlig
Sammenhæng er det eneste Kriterium, vi i Tvivlstilfælde have, naar det gælder
at skelne mellem Drøm eller Illusion og Virkelighed. Jo mere omfattende og jo
inderligere sammenhængende et Virkelighedsbegreb vi kunne danne, des større
Sikkerhed indeholder det. Der vil derfor baade af teoretiske og praktiske
Grunde være en Tilbøjelighed til at gaa til Grænsen og forsøge paa at fortsætte
og afslutte, hvad allerede teoretisk og praktisk Orientering fører ind paa.
Idealet vilde være naat, hvis der kunde paavises en fuldstændig Harmoni mellem
alle vore Erfaringer, en kontinuerlig Helhed, til hvilken alle specielle
Erfaringsomraader efter deres egne Love slutte sig sammen.

Men allerede Drøftelsen af Erkendelsesproblemet maa have lært os, at en saadan
afsluttende Verdensanskuelse er umulig, og at den i en vis Forstand indeholder
en Selvmodsigelse. Intet af de specielle Erfaringsomraader foreligger afsluttet
for os; der kommer stedse nye Erfaringer og nye Gaader op; vor sammenfattende
Tænkning faar stedse nye Opgaver. Og da nu vor Erkendelse stedse virker ved at
sammenholde og sammenligne, vilde ethvert Totalbillede, dersom det skulde være
Genstand for fuldendt Erkendelse, være at sammenholde og sammenstille med
Noget, der var forskelligt fra det: kun da kunde det faa sin fulde Bestemthed;
men hvis der var noget fra det Forskelligt, vilde det jo ikke være et
\textit{Total}billede! --- Det er her ligegyldigt, om vi blive staaende ved
selve den Totalitet, vi have forsøgt at danne, eller om vi gaa tilbage til
Principet for en saadan Totalitet og kalde det „Gud“ i Modsætning til selve
Totaliteten som „Verden“; Antinomien bliver den samme i begge
Tilfælde\textsuperscript{49}). Det Irrationelle melder sig her igen, ligesom ved
% --- p. 55 ---
\opage{55}Erkendelsesproblemer. Og der er her sikkert en indre Sammenhæng. At
Erkendelsen ikke kunde afsluttes, kunde hænge sammen med, at selve Tilværelsen
ikke er afsluttet, ikke er færdig, men er i stadig Vorden, ligesom den enkelte
Personlighed er det, og ligesom Erkendelsen er det. Den skjuler maaske ogsaa
samtidige Disharmonier i sig, der gøre det umuligt for den at udgøre et
harmonisk Hele. Analogien mellem de forskellige Problemer vilde i saa Fald være
særligt tydelig. Man har i de forskellige Systemer været for sikker paa, at
Tilværelsen i sig selv var et sluttet og konstant Hele, og at det kun var vor
Villie og Tanke, der stedse maatte kæmpe for at bestaa og for at vinde Harmoni.

Et \textit{realt} Motiv til Spekulation ligger i den fremtrædende Rolle, et
Fænomen eller et Erfaringsomraade eller en Side ved Tilværelsen kan spille for
os. Derved kan opstaa en Tendens til at lægge dette Fænomen eller denne Side
til Grund ved en Tydning af Tilværelsen og forsøge paa at udlede de andre
Fænomener eller Sider deraf, eller i hvert Tilfælde at føre dem tilbage dertil.
De forskellige kosmologiske Systemer ere ligesaa mange Forsøg paa at lodde
Tilværelsens Dyb, at prøve, hvor langt en Tanke kan kaste Lys over Tilværelsen.
Det er en Række af Tankeexperimenter, ved hvilke vore vigtigste Tankers
Bærekraft, deres Evne til at tjene som Grundlag for en omfattende
Verdensanskuelse er bleven prøvet.

Ethvert saadant Forsøg gaar --- som \textsc{Leibniz} først med genial Klarhed
har set --- nødvendigvis \textit{Analogiens} Vej. --- I al Videnskab kan
Analogi blive af Betydning. Alle Betegnelser for psykiske Fænomener ere
oprindeligt dannede paa Grundlag af en Analogi med fysiske Fænomener, og vi ere
i Psykologien stadigt nødte til at operere med fysiske Analogier; det bliver da
et vigtigt Spørgsmaal, hvor langt deres Gyldighed naar. Aarsagssætningen
udtrykker ifølge Kant, at der mellem de virkelige Begivenheder i Verden er et
analogt Forhold som mellem Grund og Følge i vor Tænkning. Atomteorien og i det
Hele den saakaldte mekaniske Naturopfattelse betragtes i den nyeste Tid
erkendelsesteoretisk som en stor Analogi, ved hvis Hjælp de kvalitative
Forandringer i Naturen kunne beskrives og beregnes.
% --- p. 56 ---
\opage{56}Ogsaa ved Opdagelser spiller Analogien en stor Rolle; Analogien
mellem Kemi og Fysik hjalp Robert Mayer til hans Opdagelse af Energiens
Bestaaen. Men naar Analogien benyttes metafysisk eller kosmologisk, er det ikke
\textit{et enkelt} Tilværelsesomraade, der tjener til at belyse \textit{et
andet enkelt}; det er \textit{et enkelt} Omraade, der benyttes til at udtrykke
Tilværelsen \textit{som Helhed.} Denne Symboliseren bliver af anden Art og
Gyldighed end den, der finder Anvendelse paa de specielle Omraader: den kan
ikke gennemføres med fuld Konsekvens, og den kan ikke verificeres. Hverken
Analogiens Berettigelse eller dens Begrænsning kan strengt paavises. Derved ere
de kosmologiske eller metafysiske Symboler forskellige fra de videnskabelige.
--- Som jeg i min „Religionsfilosofi“ har søgt at vise, dele de religiøse
Symboler Skæbne med de metafysiske. I begge Tilfælde gøres der et Forsøg paa at
danne afsluttende Begreber med real Gyldighed; Forskellen beror kun paa det
Motiv, der fører til disse Forsøg.

I Erkendelsesproblemet kom vi ogsaa til et irrationelt Forhold mellem Del og
Helhed. Dér drejede det sig om, \textit{hvorvidt Helheden kunde udtrykkes i og
for en enkelt Del.} Her, ved det kosmologiske Problem, drejer det sig om,
\textit{hvorvidt der fra en enkelt Del kan hentes Bestemmelser, som gælde for
Helheden som saadan.} Det er et og samme Forhold, der i de to Problemer
fremtræde fra forskellig Side. ---

En Verdensanskuelses Karakter og Værdi beror ikke blot paa, hvor klart og
konsekvent Analogien er gennemført, men ogsaa paa, hvorfra Analogien hentes.
Det Fænomen (den Del eller Side af Erfaringen), hvorpaa Analogien bygges, kunne
vi kalde for \textit{Urfænomenet.} Dette Udtryk stammer fra \textsc{Goethe}.
Men Goethe forstod ved Urfænomen ikke blot et Fænomen af typisk Art, der kan
tjene til at belyse andre Fænomener for os; det betyder for ham et Faktum, der
tillige er Lov, saa at man blot behøver at udtale det for at have forklaret
det, og det maa ikke betragtes som sammensat, da det selv skal være Symbol for
alt Andet, være en „Grænse for vor Skuen“, som paa én Gang vækker „den store
Undren“ og Følelsen af, at vi staa ved vor Evnes Grænse!\textsuperscript{50})
--- Det er bekendt, hvor skæbnesvanger en Rolle dette Begreb spillede i Goethes
Farvelære. Ved selve sin
% --- p. 57 ---
\opage{57}Definition af det vilde han udelukke enhver videre Undersøgelse og
Forklaring af det Fænomen, han valgte til „Urfænomen“. Men denne Mangel behøver
ikke at klæbe ved det. Det kan meget vel betegne det Punkt i vor Erfaring, ud
fra hvilket vi forsøge at orientere os i alle Retninger, men som derfor godt
selv kan være Genstand for videregaaende Undersøgelse og Forklaring. Hovedsagen
er, at vedkommende Fænomen er saa ejendommeligt og betydningsfuldt, at det kan
faa en typisk Karakter for os. Tydningen af Tilværelsen maa stedse foregaa ud
fra et enkelt Omraade, og den kan foregaa, uden at dette Omraade selv unddrages
speciel videnskabelig Behandling.

Men hvor skulle vi nu søge Urfænomenet? Derom staar Kampen mellem de
forskellige Verdensanskuelser. Snart er det Livet, snart Tanken, snart
Materien, der gøres til Urfænomen og lægges til Grund ved Tydningen. Førend vi
gaa over til at betragte de vigtigste Urfænomener, der have været tagne i den
kosmologiske Tydnings Tjeneste, maa vi dog dvæle lidt ved denne Tydning selv.

2. Metafysiken kan blive dogmatisk og derved virke skadeligt i videnskabelig
Henseende paa to Maader --- dels ved at afskære enhver speciel videnskabelig
Forklaring af det valgte „Urfænomen“, dels ved at opfatte sin analogiserende
Metode som mere videnskabelig, end den er, og at glemme den stadige
Erfaringsbekræftelses Nødvendighed. Den kan ikke gribe ind i de specielle
Videnskabers Husholdning, i Aandsvidenskabernes ligesaa lidt som i
Naturvidenskabernes. Dens filosofiske Sted er paa den videnskabelige Tænknings
Grænse; dens Opgave er at give en sidste Tydning. Erkendelsesteorien førte os
til et Grænsebegreb om det, der betinger den stadige Brydning i vor
Verdensopfattelse mellem Kvalitet og Kvantitet, elementært og idealt
Aarsagsbegreb, Subjekt og Objekt. Hvis vi ville kalde det, der indeholder
Grunden til disse Modsætninger, „Tingen i sig selv“, saa betegner dette Udtryk
Metafysikens eller Kosmologiens filosofiske Sted.

Den sidste Tydning vil have Mere af Kunstens end af Videnskabens Præg.
Overgangen fra Filosofiens empiriske og kritiske Dele (Psykologi og
Erkendelsesteori) til Kosmologien frembyder
% --- p. 58 ---
\opage{58}en vis Analogi til Overgangen fra historisk Kritik til
Historieskrivning. Historieskriveren former det kritisk behandlede Stof til et
Helhedsbillede af Begivenheder og Karakterer; det Fragmentariske udfyldes, og
der gives en Afrunding, ved hvilken Fremstillerens Personlighed nødvendigvis
spiller en Rolle med. Den kosmologiske Filosof vil saaledes give os et
Helhedsbillede af Tilværelsen paa Grundlag af vor Viden om den; han vil
sammenarbejde de spredte Træk til en Totalitet, indenfor hvilken et enkelt
Element (Urfænomenet) særligt virker bestemmende, saa det Heles Klangfarve
afhænger af det. I Valget af Urfænomen og i Analogiens Gennemførelse gør et
afgjort personligt Element sig gældende. Et stort filosofisk System er et
Kunstværk, et Drama, i hvilket de enkelte Repliker og Scener slutte sig sammen,
snart mere med logisk, snart mere med æstetisk-etisk Harmoni. Aristoteles',
Spinozas og Schopenhauers Systemer frembyde særligt denne Karakter.

Denne Kunst vil blive vanskeligere og vanskeligere at øve. Ti den vil for
Fremtiden ikke blot forudsætte, at et rigt Stof staar til Raadighed og kan
formes med konstruktiv Kraft, men maa ogsaa forudsætte Evne til at undgaa
Dogmatiseren, til at hævde Ideernes og Tankebygningernes Betydning uden at
forvexle dem med absolute Sandheder. Det gælder, som \textsc{Lessing} sagde, at
tænke gymnastisk, ikke dogmatisk. Kunsten vil være den: at forene dristig
Tankeudvikling med vaagen kritisk Bevidsthed. Det er denne Side af Sagen, jeg
særligt har betonet i min Afhandling \textit{Filosofi som Kunst.} Tidligere er
det kunstneriske Element i filosofisk Tænkning blevet betonet af
\textsc{Schopenhauer}, som dog mest tænker paa den \textit{uvilkaarlige}
Tilbliven af den afgørende Tankekonstruktion, hvorved Kunstens „Hvad?“ afløser
Videnskabens „Hvorfor?“ Naar ogsaa \textsc{Albert Lange} kaldte filosofisk
Spekulation for Kunst, tænkte han nærmest paa den \textit{idealiserende}
Tendens, Trangen til i ideale Billeder at se et Udtryk for den højeste
Virkelighed\textsuperscript{51}). --- Ved det religiøse Problem ville vi møde
et Forhold, der er beslægtet med den Overgang fra Videnskab til Kunst, vi her
have talt om; kun vil det der dreje sig om en Afslutning af
\textit{Livs}anskuelsen, medens det her drejer sig om en Afslutning af
\textit{Verdens}anskuelsen. ---

% --- p. 59 ---
\opage{59}Det er, som Filosofiens Historie viser, ikke til alle Tider, der er
Betingelser til Stede for Udøvelsen af den her omtalte Kunst i stor Stil. En
lang Tids Ophoben af Stof og af Synspunkter, en stor Koncentration i Sindet og
en ved Skærpelse af Problemerne vakt Tankeenergi maa være til Stede. Vi finde i
Filosofiens Historie Systemerne ligge samlede i visse enkelte Perioder, endog
tillige bundne til enkelte Steder. Det 5. og 4. Aarhundrede f. Kr. skabte de
store græske Systemer, og her var Aten det filosofiske Centralsted. I det 17.
Aarhundrede opstod den nyere Tids Hovedsystemer, der bleve til paa Grundlag af
Renaissancens og den nys grundlagte Naturvidenskabs mangfoldige Stof og nye
Synspunkter; her indtager Holland en central Stilling. Ved Begyndelsen af det
19. Aarhundrede grundlagdes alle den filosofiske Idealismes Systemer i Jena,
hvor Kants Disciple mødtes med Goethes Disciple og med Romantikens Forkæmpere.
Det er i saadanne Tidspunkter, at Filosofien staar som et af Symptomerne paa
Aandslivets Gang; i Mellemtiderne gaar den speciellere Drøftelse af
psykologiske, erkendelsesteoretiske og etiske Problemer sin mere stille Gang. ---

Jeg gaar nu over til at fremdrage de vigtigste Urfænomener, der kunne tjene til
systematisk Grundlag og historisk have ydet saadan Tjeneste.

3. Den første Kendsgerning, det ligger nær for Filosofien at betragte som
Urfænomen i Tilværelsen, er den, at Tilværelsen i stort Omfang er
\textit{forstaaelig} for os. Vi kunne genkende Fænomener, slutte fra det ene
Fænomen til det andet og finde Kontinuitet imellem dem. Selv om vore Hypoteser
egentligt ere Arbejdshypoteser, og selv om det gamle, naive Sandhedsbegreb mere
og mere maa vige Pladsen for det dynamiske eller symbolske Sandhedsbegreb (se
III, 2), saa maa dog selve det Faktum, at vi med vore Evner og vore Metoder til
en vis Grad kunne orientere os i Tilværelsen, hænge sammen med dennes Væsen.
Metodens Anvendelighed indeholder stedse et Vidnesbyrd om Emnets Beskaffenhed.
Tilværelsens Forstaaelighed for os tyder paa en indre Enhed i dens Væsen, der
lægger sig for Dagen i den Lovmæssighed, Fænomenernes Gang frembyder. Selve
vort Virkelighedskriterium bestaar i den faste Sammenhæng, og det
% --- p. 60 ---
\opage{60}er en naturlig Udvidelse af, hvad Anvendelsen af dette Kriterium
godtgør, naar vi i Tilværelsen finde en \textit{Enhedsmagt,} ved hvilken de
enkelte Elementer og Begivenheder forbindes. Filosofien har en naturlig Tendens
til at gaa tilbage til et saadant Princip; Platon og Giordano Bruno, Spinoza og
Hegel, Fechner og Lotze arbejdede i denne Retning, og selv Kant har, trods sin
store kritiske Betænkelighed, aflagt Vidnesbyrd om denne Tankes Betydning.
Medens den populære Tankegang er tilbøjelig til at søge en Afslutning ved at
stige op ad Aarsagernes Stige for at ende i en første Aarsag, --- en Vej, der
kun fører til en endeløs Række, --- søger den filosofiske Tænkning hellere sin
Afslutning i Dybden end i Længden, idet den besinder sig paa, hvad der er
Forudsætningen for, at et rationelt Verdensbillede, hvor hypotetisk det end
maatte være, i det Hele kan dannes. \textit{Kausaliteten} staar her som
Urfænomenet, og det er ikke mindst af denne Grund, at Diskussionen om
Aarsagsprincipet fylder saa stor en Plads i den nyere Filosofis Historie.

Mod den Monisme, som saaledes trænger sig frem, kunde man indvende, at det
Karakteristiske for den videnskabelige Forstaaelse af Fænomenerne ikke blot er
Antagelsen af en indre Sammenhæng mellem disse, men ogsaa Antagelsen af en
Mangfoldighed af Elementer, mellem hvilke den Vexelvirkning finder Sted, som de
af Videnskaben fundne Love gælde for. Hvorfor ikke betragte denne Mangfoldighed
som Urfænomen, saa at vor Metafysik bliver pluralistisk, ikke monistisk?

Svaret paa dette Spørgsmaal er givet i den Kendsgerning, at alle de Egenskaber
eller Kræfter, med hvilke vi udstyre Tilværelsens Elementer (hvad enten vi
forestille os dem som materielle Atomer eller psykiske Monader), hentes fra den
lovmæssige Sammenhæng, i hvilken disse Elementer ere Led. Tingenes Egenskaber
ere de konstante Maader, paa hvilke de paavirke hverandre og paavirkes af
hverandre. „Kraft“ er her det Element, der for os indeholder Grunden til en
Forandring i et eller flere andre Elementer. Saaledes gaar det med de
beslægtede Begreber Energi, Evne, Individualitet. Alle disse Begreber ere
afledede i Forhold til Begrebet Lov. For Individualitetens Vedkommende maa her
ikke blot tænkes paa den Lov, der gælder
% --- p. 61 ---
\opage{61}for Individets Forhold til Omverdenen, men ogsaa paa Loven for
Forholdet mellem Individets indre Tilstande.

I samme Retning peger en anden Betragtning. Al Erkendelse begynder som Analyse
af givne Anskuelser (Iagttagelser eller Erindringer). For at en saadan
Anskuelse skal kunne blive Genstand for Analyse, maa den fremtræde som en
Totalitet, der omfatter et Indbegreb af Elementer (Dele eller Egenskaber).
Enhver Erkendelsesakt beror paa et Udsnit af et større Hele: enhver
Begrebsbestemmelse (Definition) er en Begrænsning; enhver Dom hviler paa
Forudsætninger, der pege ud over selve Domsakten; enhver Slutning forudsætter
flere Præmisser, hvis Gyldighed ofte maa begrundes ad højst forskellige Veje.
Vi bevæge os altid indenfor en mere omfattende Helhed, i hvilken Betingelserne
maa søges for de enkelte Resultater, vi komme til.

Naar saaledes baade i Virkelighedens og i den rene Tankes Verden det Enkelte
faar sin Beskaffenhed og sin Gyldighed ved den Enhedssammenhæng, i hvilken det
forekommer, vil Monismen staa som den, der tager et mere fundamentalt
Udgangspunkt end Pluralismen.

Imidlertid viser jo Erkendelsesteorien, at der er Grænser for Forstaaelsen.
Empirikeren og Skeptikeren ville stedse kunne stanse den monistiske Metafysiker
ved at henvise til de faktiske Grænser for Erkendelsen. End ikke
Virkelighedskriteriet kunne vi anvende overalt, og saaledes ikke gennemføre den
strenge Forskel mellem Drøm og Virkelighed. Med samme Ret, man slutter fra
rationel Erkendelses Mulighed til en Enhedsmagt i Tilværelsen, vilde man da,
synes det, kunne slutte til et irrationelt Element i Tilværelsen, til et
kosmologisk Princip, der hindrede Tilværelsens Elementer i at staa i et
rationelt bestemmeligt Forhold til hverandre.

Fra monistisk Side vil dertil igen kunne svares, at dersom Enhedsmagten ikke
raader overalt, kunde det tyde paa, at Tilværelsen selv er i Vorden, i
Udvikling, og at det, der for os staar som en lovmæssig Sammenhæng, forudsætter
en Udvikling i Tilværelsens Inderste, der endnu ikke er afsluttet. Det vilde da
ogsaa fra dette Synspunkt være \textit{Tiden,} der betinger det Irrationelle.
Og saalænge Tanken, Erkendelsen ikke er færdig,
% --- p. 62 ---
\opage{62}men er stedt i Vorden, saalænge \textit{kan} Tilværelsen i det Hele
heller ikke være færdig. Tanken, Erkendelsen er jo dog en Del af Tilværelsen!
Og naar \textit{den} er i Vorden, kunde der jo være Mere, som er i Vorden. Det
er de store rationalistiske Systemers forunderlige Modsigelse, at selv om de
kunde forklare alt Andet, kunde de dog ikke forklare selve den stræbende og
arbejdende Tanke, hvis Produkter de ere. Hos Platon og Spinoza, hos Hegel og
Boström fremtræder denne Modsigelse.

\textit{Den kritiske Monisme,} som fastholder Tidens Gyldighed og dermed
Tilværelsens Uafsluttelighed saavel som Erkendelsens, kan meget vel i sin
Verdensanskuelse bygge paa Kausaliteten og Rationaliteten som Urfænomener. Den
finder da i Tilværelsen en kæmpende Enhedsmagt, der gennem fremadskridende
Udvikling føres ud over det Sporadiske og Modstridende. Maaske opstaar der
stedse nye Elementer, der saaledes skulle sammenarbejdes, og Udviklingen maa da
stedse begynde paa ny. Tankens eget Arbejde bliver herved set i en kosmisk
Belysning: en konstant Sammenhæng mellem vore Metoder og Hypoteser og
Tilværelsens reale Processer er jo det Maal, Tanken sætter sig (selv naar den
ombytter det statiske Sandhedsbegreb med det dynamiske); lykkes det Tanken at
nærme sig til dette Maal, vil Tilværelsen derved være bleven mere rationel, end
den forud var, idet et nyt konstant og harmonisk Forhold er realiseret. Den
Tænker, hvem det gives at arbejde fremad i denne Retning, kan med Rette sige:
„Ort für Ort \textit{sind} wir im Innern“ --- hvor fjernt og ophøjet Tankens
højeste Ideal end staar for ham.

4. Kunne vi ikke forsøge en real, positiv Bestemmelse af det Enhedsprincip, som
efter den nys udviklede Hypotese holder Tilværelsen sammen i dens Inderste? ---
Ethvert Forsøg i denne Retning maa i særligt høj Grad gøre Brug af Analogi. Vi
have her ikke saa fundamental en Kendsgerning at vise hen til som ved
Opstillingen af hin Hypotese, hvor vi kunde støtte os til Tilværelsesproblemets
Sammenhæng med Erkendelsesproblemet. Ved en nærmere Bestemmelse af
Tilværelsesprincipet vil Tanken henledes paa den mest fundamentale Forskel, de
os tilgængelige Fænomener frembyde: Forskellen mellem det Aandelige og
% --- p. 63 ---
\opage{63}det Materielle. Der synes her at være bundet Valg for ethvert Forsøg
i kosmologisk Retning,

Hvorledes jeg rent metodisk og empirisk stiller mig til Problemet om Forholdet
mellem Aand og Materie, har jeg allerede udviklet (II, 3). Men Spørgsmaalet
kommer her igen fra et andet Synspunkt. Ti hvad enten vi metodisk og empirisk
ere Dualister, Materialister eller Monister, kan der spørges, hvorledes vi
tænke os Tilværelsens fundamentale Væsen, --- hvilket det til Grund liggende
Attribut er, det, under hvilket vi tilsidst --- naar vi tænke vor sidste Tanke
--- tænke os Tilværelsen.

Lægger jeg Analogien med det i Rummet Udstrakte og Bevægelige til Grund, saa
bliver min Metafysik, ikke blot min Metode, materialistisk. Som afsluttende
Hypotese har \textit{Materialismen} det Fortrin, at den giver os Billedet af en
stor Kontinuitet, og at vi ikke behøve at forlade det Anskuelige. Materialismen
er en barnlig og naiv Opfattelse; den er Menneskets første Filosofi. Men
stadigt øver Indtrykket af den materielle Verdens Sammenhæng og Omfang en
overvældende Magt, saa at Forsøg i materialistisk Retning atter og atter komme
frem, selv om de efter den kritiske Filosofis Fremtræden ikke kunne møde med
den dogmatiske Sikkerhed som tidligere. De ville stedse strande dels paa
Umuligheden af at føre det Aandelige tilbage til det Materielle, dels paa den
erkendelsesteoretiske Betænkelighed, at det Materielle have vi stedse kun
\textit{som Objekt} for Bevidstheden, og at der, hvis Materialismen havde Ret,
ikke kunde existere Noget, \textit{for hvilket} det Materielle kunde blive
Objekt eller Fænomen. \textsc{Hobbes}, hvis Tænkning gik stærkt i
materialistisk Retning, stansede dog netop ved den Betragtning, at af alle
Fænomener er \textit{selve det, at Noget kan blive Fænomen for os,} det vigtigste.

I Modsætning til Materialismen lægger \textit{den metafysiske Idealisme}
Analogien med de aandelige Fænomener til Grund. Den kan empirisk hylde
Dualismen (som \textsc{Lotze}) eller Materialismen (som \textsc{Schopenhauer})
eller den spinozistiske Monisme (som \textsc{Leibniz, Fechner} og
\textsc{Wundt}). Den sidste Tanke er i alle Former den, at kun Analogien med de
aandelige Tilstande, vi kende fra os selv, giver os en Nøgle til Forstaaelse af
Tilværelsens
% --- p. 64 ---
\opage{64}inderste Væsen. Kun os selv kende vi jo indefra, alt Andet blot
udefra! Allerede psykologisk Erfaring lærer os, at der gives mange Grader og
Arter af aandelig Tilværelse, og ville vi benytte den til Fortolkning af hele
Tilværelsen, dens inderste Grund og dens sidste Elementer, saa maa vi
naturligvis tænke os disse Grade- og Kvalitetsrækker fortsatte i det
Ubegrænsede. Denne idealistiske Kosmologi kan varieres paa mange Maader, og
fremtræder historisk i mange Former, der især betinges ved de andre Motiver,
der bestemme Verdensanskuelsen (f. Ex. Tendensen til Optimisme eller
Pessimisme, Hovedvægt paa Tanke eller Villie o. s. v.). Allerede i de indiske
„Upanishad“ fremtræder den: i Læren om, at Bráhman (Verdensprincipet) er Atman
(Sjæl).

Undertiden negter man, at det er en Analogislutning, man her støtter sig til,
og paastaar at være naat til den metafysiske Idealisme (eller den idealistiske
Kosmologi) ad den logiske Konstruktions, den dialektiske Metodes strenge Vej.
Men selv den stolteste Tankebygning, man saaledes har trot at kunne opføre,
\textsc{Hegel's} System, gaar dog egentligt ud paa at vise, at Alt i
Tilværelsen hænger ligesaa inderligt sammen som Tankerne i Menneskets Aand: det
er i Virkeligheden denne, der benyttes som Grundlag for en Analogiseren, idet
den staar som bedste Exempel paa en inderlig Totalitet\textsuperscript{52}). ---

Den idealistiske Analogisering vilde kun kunne føre til en afsluttende Løsning,
dersom vi havde Midler til paa positiv Maade at bestemme de forskellige Grader
og Arter af aandelig Existens, der maa findes i Verden, hvis den metafysiske
Idealisme har Ret. Men vi kunne, som det allerede har vist sig i en tidligere
Sammenhæng (II, 3), her ikke komme videre end til det negative Begreb om
potentiel psykisk Energi. En Verifikation af Idealismen er altsaa umulig. Og
selv om den kunde gennemføres fuldt ud, vilde endnu den Vanskelighed staa
tilbage, at det Materielle ligesaa lidt kan udledes af det Aandelige som dette
af hint.

Men det er ingenlunde sikkert, at vi have bundet Valg mellem materialistisk og
idealistisk Løsning af Tilværelsesproblemet. Forskellen mellem Materie og Aand
er ganske vist en Hoved%
% --- p. 65 ---
\opage{65}forskel i vor Erfarings Indhold; men der vil intet Bevis kunne føres
for, at der i Tilværelsen ikke findes andre Grundformer („Attributer“) end
disse to. \textit{Skal} Tilværelsesproblemet løses ud fra menneskelig Erfaring,
saa maa en af hine to Muligheder vælges. Men Spørgsmaalet er, om vor Erfaring
frembyder os Elementer nok til en virkelig Løsning. Tilværelsens Rigdom kan jo
være meget større end vor Erfarings Muligheder. Atter her kan det gælde, at
Verden er stor og vor Hjerne lille; atter her støde vi paa det Irrationelle.
Kunde det bevises, at Forskellen mellem Aand og Materie var kontradiktorisk,
ikke kontrær, altsaa at der forelaa et absolut Enten-Eller, vilde
Tilværelsesproblemet være simplere, end det er --- og dog mere indviklet, end
Mennesketanken, der er tilbøjelig til at mene sig fuldt udrustet til religiøse
og metafysiske Spekulationer, ofte mente. \textit{Den kritiske Monisme,} som
stræber at fastholde Enhedstanken uden at dogmatisere, maa indskærpe, at vi her
mangle Betingelser for en fuldstændig Løsning. At der er Mulighed for flere
Former, end vor Erfaring frembyder, kan tyde paa, at hele Problemet ligger
dybere, end man har formodet. Der kunde jo f. Ex. gives en Grundegenskab ved
Tilværelsen, af hvilken baade Aand og Materie vare Følger, og Problemets
Uløselighed for os kunde komme af, at vi ikke kende denne Grundegenskab.

5. Et tredie Urfænomen have vi, naar vi have truffet vort Valg mellem Bestaaen
og Udvikling (Væren og Vorden). Den antike Tænkning var gennemgaaende
tilbøjelig til at holde paa bestaaende Væren; den var en Begrebsfilosofi, der
fremfor Alt søgte at føre Fænomenerne tilbage til faste Artsbegreber. Platons
Idelære danner her Højdepunktet (Smig. II, 1). Og den --- eller den samme
Aandsretning, som har fundet Udtryk i den --- kommer ikke sjeldent frem ogsaa i
den nyere Forskning, og det ikke blot i Filosofien, men overalt i Videnskaben,
hvor der kæmpes for faste, uforanderlige Arter eller Typer. Selv hvor den
antike Tænkning antog en Udvikling, var den særligt tilbøjelig til at antage en
rytmisk Proces, der atter og atter førte de samme Tilstande og Begivenheder med
sig. Udviklingen som en stedse fremadskridende Række af Forandringer er en
moderne Ide, som skyldes udvidet Erfaring baade paa Historiens
% --- p. 66 ---
\opage{66}og paa Naturens Omraade, men som tillige har formet sig under
Indflydelse af den persisk-kristelige Religionstype\textsuperscript{53}).

At Begrebet Udvikling kommer til at spille saa stor en Rolle for den nyere Tids
Tænkning, hænger sammen med, at Aarsagsbegrebet er kommet saa stærkt frem. Jo
mere det elementære Aarsagsbegreb ved fortsat Forsken nærmer sig til det ideale
Aarsagsbegreb, des mere kommer Aarsagsforholdet til at betegne en kontinuerlig
Proces, i hvilken de følgende Led ere bestemte ved de foregaaende. Og
Udviklingsbegrebet fremtræder, saasnart det viser sig, at \textit{Retningen} er
væsentlig, saa at Tiden ikke kan vendes om (Smlgn. III, 4). Medens de formale
Videnskaber, der hvile paa Identitetsprincipet, kunne gaa frem og tilbage i
deres Tankerækker, faar Tidsforholdet, særligt Tidsretningen en afgørende
Betydning for de reale Videnskaber, og disse beholde derved en historisk
Karakter. Her er selve Begrebet Begivenhed et Urfænomen. Men i Begrebet
Udvikling ligger, at der ikke blot sker Noget, men at der gennem
Begivenhedernes Række naas Tilstande, som have et vist Helhedspræg, idet en
Mangfoldighed af Elementer ere forbundne paa en saadan Maade, at de trods deres
Forskelligartethed virke i fast Forbindelse med hverandre med en vis
Afsluttethed overfor Omgivelserne. Udførligst og med en stor Rigdom af Exempler
har \textsc{Herbert} \textsc{Spencer} analyseret Begrebet Udvikling og som det
væsentlige Træk hævdet Forbindelsen af Differentiation og Koncentration. Det er
ved dette Begreb som Maalestok, vi bestemme, hvad der --- fra et rent teoretisk
Synspunkt --- skal kaldes „lavere“ og „højere“ indenfor en skiftende Række af
Tilstande eller Former.

Trods sin Sammenhæng med Aarsagsbegrebet er Udviklingsbegrebet et selvstændigt
Begreb, der ikke kan udledes af det almindelige Aarsagsbegreb (skønt det
forudsætter det). Aarsagssætningen vilde ogsaa gælde, selv om der kun fandt en
rytmisk Svingning Sted uden fremskridende Dannelse af nye Helhedstilstande.
Udviklingen staar som et Erfaringsfaktum, der kaster Lys over Tilværelsens
Natur. Det var i og for sig meget vel muligt, at de forskellige Aarsagsrækker i
Tilværelsen enten slet
% --- p. 67 ---
\opage{67}ikke mødtes, eller ogsaa kun mødtes for at støde sammen paa
disharmonisk Maade. Nu viser Erfaring, at de under visse Forhold kunne mødes
saaledes, at de forbindes til mere sammensatte Processer og afføde
ejendommelige Helheder. Stjernesystemernes Opstaaen, Livets Opstaaen, Bestaaen
og Udfoldelse, Sjælelivets og det sociale, historiske Menneskelivs Existens ere
Vidnesbyrd om en individualiserende og totaliserende Tendens i Tilværelsen.
Disse Fænomener stille stadigt Forskningen de største Problemer; naar vi her
henpege til dem som Urfænomener, afgøre vi intet om, hvorvidt disse Problemer
kunne løses eller ikke, men betragte dem blot som Karakteristika for
Tilværelsen, saaledes som denne nu engang er. Det er sikkert --- dette maa dog
tilføjes her --- en stor Misforstaaelse, naar man ofte tror, at disse Fænomener
vilde være mærkeligere, hvis de \textit{ikke} skulde kunne forklares ved Hjælp
af Love, Videnskaben kan finde. Tvertimod, hvis der kan findes en „naturlig“
Forklaring af dem, vil deraf kunne sluttes med større Sikkerhed end før, at den
individualiserende og totaliserende Tendens ligger dybt grundet i Tilværelsens Kærne.

Men hvor dybt? Erfaringen viser os ikke blot ofte et rent udvortes og
ligegyldigt Forhold mellem forskellige Aarsagsrækker, men ogsaa ofte Sammenstød
mellem dem eller dog disharmoniske Forhold, der hæmme eller hindre
Individualiteters og Totaliteters Opstaaen eller Bestaaen. Og efter Udvikling
følger Opløsning, og Spørgsmaalet bliver, om der gennem den rytmiske Skiften af
disse Processer kan spores en fremskridende Udviklingsgang, eller om vi skulle
stanse ved den antike Tanke om Rytmen som den sidste Tanke.

Her fremtræder igen det Irrationelle, og det fremtræder her i skarpere Form end
paa de tidligere Punkter, hvor vi mødte det; det betegner her det Hæmmende og
Opløsende, Tendensen til at blive staaende ved eller vende tilbage til de mest
elementære Tilværelsesformer. Tilværelsen giver os her, jo mere Forskningen
trænger ind i den, Billedet af en Kamp, den store Kamp, som alle
Tilværelsesformer, der bære Individualitetens og Totalitetens Præg, maa kæmpe
for at bestaa. Selve Kampen er igen
% --- p. 68 ---
\opage{68}dobbeltsidig: den kan være et Middel til Udvikling, men den kan ogsaa
være en Vej til Undergang, og hvilken af disse Muligheder er den overvejende?

Tillige viser sig her endnu tydeligere end paa de tidligere Punkter Umuligheden
af at danne et absolut afsluttet Begreb om Tilværelsen som Helhed: dersom
Brydningen mellem Elementer og Totaliteter er et \textit{væsentligt}
Karaktertræk ved Tilværelsen, saa er det et Karaktertræk, der stedse kun kan
findes ved begrænsede Udsnit af Tilværelsen, ikke ved Tilværelsen som et Hele
betragtet, da et absolut Hele ingen ydre Modstand kan møde, ingen Kamp kan være
stedt i. Vi maa lade os nøje med at sige, at overalt hvor vi træffe
Tilværelsen, finde vi, at der indenfor den foregaar en saadan Brydning mellem
Elementer og Totaliteter. Det store Spørgsmaal er, om det gennem denne Brydning
er Elementerne eller Totaliteterne (Solsystemerne, Organismerne, Sjælene,
Samfundene), der gaa af med Sejren. Empirisk staa vi her midt i den store
Verdensproces og kunne ikke naa videre end til at hævde, at det ikke er nogen
rent subjektiv Maalestok, vi anlægge, naar vi betegne en Tilværelsesform som
højere eller lavere end en anden, da Udvikling er et for al den Tilværelse, vi
kende, karakteristisk Fænomen, --- et Urfænomen. ---

Ved det Resultat, ved hvilket vi saaledes ende for Tilværelsesproblemets
Vedkommende, bliver en naturlig Overgang til det sidste Problem,
Vurderingsproblemet, mulig.

Dersom Tilværelsen var færdig, harmonisk og uforanderlig, vilde ingen Etik være
mulig. Al Etik kræver et Arbejde udført; men der vilde ikke være Plads for
noget Arbejde, dersom Alt bestod i evig og aktuel Fuldendthed. Netop det, at
alle Individualiteter og Totaliteter maa kæmpe for at bestaa, --- at der
stadigt er Disharmonier at overvinde, sporadiske Tendenser at forbinde, ---
netop det gør en Etik mulig. Etiken undersøger nemlig Principerne for en
Videreførelse af den individualiserende og totaliserende Tendens, som
Erfaringen paaviser i Tilværelsen. Etik beror paa en Selvmodsigelse, dersom
Verdensløbet ikke for en Del foregaar eller kan foregaa gennem men%
% --- p. 69 ---
\opage{69}neskelig Villie, ligesom det for en Del foregaar eller kan foregaa
gennem menneskeligt Tankearbejde (III, 5). Etiken tager da
Kontinuitetsproblemet op, hvor de tre første Problemer slippe det.

Endnu tydeligere finder det religiøse Problem en Tilknytning i den foregaaende
Betragtning, da det --- som jeg skal forsøge at vise --- angaar Værdiernes
Kontinuitet under Tilværelseskampen.

\begin{center}\rule{2cm}{0.4pt}\end{center}

% --- p. 70 ---
\opage{70}

\begin{center}IV.\\[6pt]
\textbf{\Large Vurderingsproblemet.}\end{center}

1. At Vurderingsproblemet opstilles som et særskilt Problem, forudsætter en
Sondring mellem Forstaaelse og Vurdering, som ikke altid har været anerkendt,
og som endnu ikke er trængt igennem. Der har i Filosofien været en
Tilbøjelighed til at lade dem gaa i Et; saaledes ere Platons „Ideer“ og
Spinozas „Substans“ paa én Gang Udtryk for de to Tænkeres Forstaaelse af
Tilværelsen og for deres Beundring for den. I hvert Tilfælde har man været
tilbøjelig til at lade Vurderingen være en blot Følge af Forstaaelsen, medens
omvendt mystiske og teologiske Retninger have været tilbøjelige til at lade
Forstaaelsen være afhængig af Vurderingen. Aldeles uafhængige af hinanden ere
de ganske vist ikke: Forstaaelsen har sin Værdi og vilde ikke være mulig, hvis
det ikke var saa, da Menneskene ellers ikke vilde stræbe efter den; for Nogle
er endog Erkendelsesglæden det Højeste; --- og Vurderingen kan blive Genstand
for psykologisk og historisk Forstaaelse, da al Værdi beror paa Oplevelsers og
Tilstandes Forhold til Livet paa de forskellige Trin, til Livets Bestaaen og
dets Udvikling. Vurderingsproblemet vil da ogsaa vise sig at frembyde en
Analogi til de tre første, teoretiske Problemer. Det er paa Vurderingens
Omraade, ligesom paa Personlighedens, Erkendelsens og Tilværelsens Omraade,
Kontinuitetsprincipet, der fører til de afgørende Problemer.

\textit{Værdi} har Alt, hvad der medfører en Tilfredsstillelse eller afhjælper
en Trang. Undertiden blive vi først ved, at en Til%
% --- p. 71 ---
\opage{71}fredsstillelse opstaar, opmærksomme paa, at der har været et Hul i
vor Existens; undertiden mærkes dette Hul forud og volder et Savn eller
fremkalder Drift og Begær. Naar det Værdifulde ikke umiddelbart falder i vor
Lod, gøre vi det til \textit{Formaal} og søge Midler til at opnaa det. Hvad der
staar som Middel til at opnaa et umiddelbart Værdifuldt, faar middelbar Værdi
for os. I vore Værdier og vore Formaal kundgør sig vor Følelses og vor Villies
inderste Væsen. Ligesom Begrebet Formaal afhænger af Begrebet Værdi, saaledes
afhænger ogsaa Begrebet \textit{Norm} af Begrebet Formaal: Normen er Regelen
for den Virksomhed, der er nødvendig for Formaalets
Opnaaelse\textsuperscript{54}). --- Det fik en skæbnesvanger Indflydelse paa
Vurderingsproblemets Behandling, at \textsc{Immanuel} \textsc{kant} vilde vende
Forholdet om og udlede Begreberne Formaal og Værdi af Begrebet Norm („Lov“).
Det er en psykologisk Umulighed.

Nu viser Erfaring, at der for forskellige Individer, og for samme Individ til
forskellige Tider, gælde forskellige Værdier. Jeg har allerede nævnt
Erkendelsens Værdi; foruden den er især at nævne de Værdier, der ere knyttede
til Selvopholdelsestrangen og til Livets Bestaaen og Udvikling i større eller
mindre Kredse, --- til Bevægelse og Virksomhed, --- til Billeder
(Virkelighedens eller Fantasiens) o. s. v. Skal der anstilles en Sammenligning
mellem forskellige Værdier, --- og enhver bevidst Vurdering bestaar i en saadan
Sammenligning, --- saa maa der forudsættes en \textit{Grundværdi,} ud fra
hvilken de forskellige Værdiers Rangfølge kan fastsættes. En bestemt Værdi maa
lægges til Grund og afgive Maalestokken for alle andre Værdier, dersom
konsekvent Tænkning og Drøftelse skal blive mulig paa Værdiernes Omraade. Ud
fra en given Grundværdi kan der saa --- under Forudsætning af tilstrækkelig
Erfaring --- konstrueres et Vurderingssystem, i hvilket hver enkelt Værdi faar
sin Plads efter sit Forhold til Grundværdien. Tænkning, „praktisk Fornuft“
behøves for at bestemme dette Forhold, og tillige for at finde Midler og Veje
til at frembringe eller udfinde de enkelte Værdier under visse givne Forhold.
Men paa bar Grund er ingen Vurdering (eller „Omvurdering“) mulig. Begrebet
Grundværdi er indenfor Vurderingsproblemet (altsaa i Etik og Religions%
% --- p. 72 ---
\opage{72}filosofi), hvad Begrebet Urfænomen er indenfor Tilværelsesproblemet
(altsaa i Metafysiken). ---

Vurderingsproblemet deler sig i to Problemer, det etiske og det religiøse.
Etisk Vurdering angaar menneskelige Handlinger, Egenskaber og Livsordninger
(Institutioner); religiøs Vurdering strækker sig videre og vurderer Tilværelsen
efter Værdiernes Skæbne i Virkelighedens Verden. Indenfor begge Problemer vil
--- ligesom ved de tidligere Problemer --- Forholdet mellem Kontinuitet og
Diskontinuitet vise sig at være af afgørende Betydning.

\begin{center}\textbf{\large A.\ \ Det etiske Problem.}\end{center}

\begin{center}\textit{a. Det etiske Arbejde.}\end{center}

2. Der er, som Drøftelsen af Tilværelsesproblemet har vist, Plads for et
Arbejde, hvorved Tilværelsen udvikles videre. Et saadant Arbejde udføres i al
menneskelig Kultur, men ganske særligt i den etiske Stræben. Denne kan betegnes
som en Stræben efter at frembringe større Kontinuitet dels indenfor den enkelte
Personlighed, dels mellem de forskellige Personligheder indbyrdes. Maalestokken
for etisk Stræben, altsaa Principet for etisk Vurdering vil vise sig at være
bestemt ved Kontinuitetsprincipet.

Personlighed forudsætter Kontinuitet, og denne betinges igen ved, at der er én
Grundværdi, som bestemmer de enkelte Øjeblikkes, Livsperioders, Evners og
Drifters Værdi. Udviklingen til sand Personlighed forudsætter en Stræben ud
over det Momentane og det Sporadiske, --- en Overvinden af det enkelte Øjebliks
eller den enkelte Trangs Tendens til Isolation og Eneherredømme. Det gælder at
indarbejde de enkelte Momenter og Elementer paa harmonisk Maade i
Personlighedslivet som Helhed. Her er et Arbejde at gøre, en Kamp at bestaa,
der hos de forskellige Individer kræver forskellig Grad af Energi.

Der er etiske Standpunkter, som hævde de enkelte Øjeblikkes og Drifters, de
successive og simultane Forskelligheders Ret overfor Livshelheden. Men dette
strider ikke imod, at Maalestokken hentes fra Kontinuitetsprincipet. Saadanne
Standpunkter kunne være berettigede, forsaavidt de kræve Begrun%
% --- p. 73 ---
\opage{73}delse af Underordningen af Øjeblikke og specielle Drifter under
Livshelheden. Kontinuitet betyder ikke Forskelsløshed, men Indordning af
Forskellighederne i en Trinrække. Livet som Helhed kan stedse kræves til
Regnskab af det enkelte Livselement. Det vil stedse staa som en Ufuldkommenhed,
naar et Øjeblik, en Livsperiode, en Evne, en Drift gøres til \textit{blotte}
Midler, uden at faa selvstændig Værdi. Kunsten vil være at lade dem paa én Gang
have umiddelbar og middelbar Værdi. --- Og forsaavidt Øjeblikkets eller
Specialitetens Etik optræder i absolut Form og proklamerer det enkelte Moments
eller Elements Suverænitet, vil man ofte finde den Forudsætning ligge til
Grund, at Livet som Helhed repræsenteres af eller er koncentreret i det enkelte
Øjeblik eller i den enkelte Trang, saa at alle andre Hensyn blegne. Saaledes er
det i Opofrelsens store Øjeblikke, hvor Mennesket --- som \textsc{Aristoteles}
siger --- „hellere vil udføre én stor og skøn Handling end mange smaa“, og
saaledes er det, hvor Mennesket i den enkelte Evnes specielle Udvikling ser sit
hele Livs Kald. Eller man tror paa en Harmoni mellem alle Livets forskellige
Øjeblikke, saa at fuldkommen Opgaaen i hvert af dem bliver mulig, uden at det
enkelte Øjebliks Tilfredsstillelse berøver de andre Øjeblikke Noget. Saaledes
synes \textsc{Aristippos} at have tænkt sig det. Eller det Højeste sættes i en
Svæven over de enkelte Øjeblikke og Sysler med Mulighed for vilkaarligt og uden
at bindes at slaa ned paa det enkelte Punkt for en Stund. Saaledes i den af
\textsc{S. Kierkegaard} i 1. Del af Enten-Eller skildrede „æstetiske“
Livsanskuelse. Under alle disse Former ligger tydeligt et Henblik til
Livshelheden bagved.

Erfares der en Modsigelse mellem det Momentane eller Sporadiske og
Helhedstrangen, vil der opstaa en mere eller mindre bevidst og mere eller
mindre energisk Stræben efter at udvikle Personligheden til et Kunstværk, hvor
hvert enkelt Øjeblik og hver enkelt Evne faar sin Plads og kommer til sin Ret.
Intet Element i det personlige Liv betragtes da som blot Middel, men stedse
tillige som Formaal. Hvad der er Middel eller Gennemgangspunkt for
Personlighedens Udvikling, skal saavidt muligt selv have umiddelbar Værdi. Det
er \textit{den græske} (\textit{platonisk-aristoteliske})\textit{ Etiks
uforglemmelige Grundtanke} om Sjælelivets harmoniske
% --- p. 74 ---
\opage{74}Udfoldelse som den store Opgave, en Grundtanke, der ganske særligt
har sin Betydning under den moderne Kulturs Tendens til at isolere eller at
mekanisere de enkelte Livselementer\textsc{. Rousseau} og \textsc{Schiller}
have fornyet denne antike Grundtanke --- overfor moderne Arbejdshelotisme og
overfor Leflen og Decadence.

Men Kontinuitetsproblemet kommer igen, og det i meget skarp Form, ved
Spørgsmaalet om Muligheden af, at den enkelte Personlighed kan danne en
afsluttet Verden for sig selv, og om det vilde være det mest Værdifulde. Ikke
blot maa den Enkelte stedse træde i Vexelvirkningsforhold til andre
Personligheder for at faa Midler for sin egen Udvikling. Men der gives ogsaa en
Hengivelsestrang, som kan fremtræde i forskellige Former, og som kan føre til
at tillægge andre Personligheder umiddelbar Værdi. Derved drages den Enkelte
ind i Personlighedernes store Rige, og ligesom Spørgsmaalet før var, hvorvidt
de enkelte Livselementer kunde indordnes paa harmonisk Maade i den enkelte
Personligheds Livshelhed, saaledes bliver det nu et Spørgsmaal, hvorvidt de
enkelte Personligheder kunne udvikle sig selvstændigt og dog i Harmoni med
hverandre, saa der kan bestaa \textit{en social Livshelhed} i Analogi med den
individuelle Livshelhed. Det er Kontinuitet i Kontinuitet, der her stræbes hen
til. Maalestokken for et Samfunds Fuldkommenhed vil --- i Kraft af
Kontinuitetsprincipet, der her tydeligt viser sig i sin Sammenhæng med
Velfærdsprincipet --- være den: i hvilken Grad ethvert personligt Væsen stilles
og behandles saaledes, at det ikke staar som blot Middel, men stedse tillige
som Formaal. Det er \textsc{Kants} berømte Sætning, med en anden Begrundelse
end den, han gav\textsuperscript{55}). Den stoiske og den kristelige Etik og
den moderne sociale og politiske Udvikling føre i samme Retning.
„Menneskerettighedernes“ Erklæring (hvad enten vi betragte den som et Symptom
eller som et Program) hviler paa denne Forudsætning, og Brodden i det sociale
Spørgsmaal er bleven til ved, at denne Forudsætning ikke er sket Fyldest. Ogsaa
ved Drøftelsen af speciellere etiske og retsfilosofiske Spørgsmaal, f. Ex. om
Monogamiets Begrundelse og om Straffevæsenets Udvikling, afgiver dette Princip
Maalestokken. ---

% --- p. 75 ---
\opage{75}Det etiske Arbejde viser sig saaledes som en ejendommelig
Fortsættelse af den store Tilværelsesproces. Men den hele Tankegang, ved
hvilken jeg har søgt at begrunde dette Synspunkt, støder paa Vanskeligheder,
der skærpe det etiske Problem. Vi møde atter her det Irrationelle.

\begin{center}\textit{b. Den etiske Vurderings Rationalitet.}\end{center}

3. Etiken vilde være en fuldkomnere Videnskab, end den er, dersom
Kontinuitetsprincipets Gennemførelse ikke stødte paa store Vanskeligheder.

Ethvert etisk Ræsonnement gælder kun ud fra Anerkendelsen af en bestemt
Grundværdi, som bestemmer alle specielle Vurderinger. Man kan stille sig paa
det enkelte Øjebliks eller den enkelte Drifts Standpunkt, eller paa den
isolerede Personligheds Standpunkt, eller paa Familiens, eller Klassens, eller
Statens eller Menneskehedens Standpunkter. Spørgsmaalet er, om alle saadanne
Standpunkter virkeligt kunne bringes i Harmoni med hverandre, saaledes som det
ovenfor er antaget, og Spørgsmaalet er især, om det er muligt ad Ræsonnementets
Vej at føre den, der konsekvent og med fornøden Hensynsløshed fastholder et
enkelt af disse Standpunkter, over paa et af de andre. Mellem Grundværdien og
de specielle Værdier kan der maaske paavises et rationelt Forhold, saaledes at
den, der vedkender sig en Grundværdi og er tilstrækkeligt bekendt med de
faktiske Forhold, under hvilke den skal fastholdes, kan logisk tvinges til at
indrømme de Slutninger, man drager ud fra Grundværdien. Saaledes have
Selvhævdelsen og Hengivelsen hver sin Logik, ligesaa Familiefølelsen og
Nationalfølelsen. Men naar det nu gælder \textit{Overgang fra en Grundværdi til
en anden?} Her er den indre Konsekvens ikke nok. Konsekvensen kan kun udfolde
og bringe til Bevidsthed, hvad der --- under visse givne, faktiske Forhold ---
følger af Grundværdien. \textsc{Sokrates} blev Etikens Grundlægger ved sin
Fordring om Selverkendelse, der i Virkeligheden var en Fordring om at blive
klar over sin Grundværdi, og hvad der konsekvent følger af den. Men hvorledes
Grundværdien skaffes til Veje, og om der ikke faktisk
% --- p. 76 ---
\opage{76}gives forskellige Værdier, der gøre Fordring paa at være
Grundværdier, drøftede han ikke nærmere.

En Grundværdi danner sig ved en psykologisk og historisk Proces, der
forudsætter flere virkende Elementer end logisk Konsekvens og Kundskab om
Fakta. Gennem Motiv- og Værdiforskydning\textsuperscript{56}) kan en Grundværdi
gaa over til en anden; men \textit{medens} en saadan fundamental Forandring paa
Følelses- og Villieslivets Omraade gaar for sig, vil det ikke nytte at
ræsonnere ud fra de Forudsætninger, som først komme til Stede, naar Processen
er endt. Ved Opdragelsen kan man ikke uden videre argumentere med Barnet ud fra
de Voxnes Forudsætninger; Barnet skal jo netop først erhverve sig disse
Forudsætninger. Og saaledes gaar det ogsaa med den store Opdragelsesproces, der
gaar for sig i Historien. \textit{Under} Opdragelsen og Udviklingen virke
naturligvis andre Grundmotiver end de, der først blive Resultatet af den hele
Proces. Den, der opdrages, vil derfor ikke rationelt kunne forstaa det System,
han opdrages efter. „Wer kann der Raupe, die im Staube kriecht, von ihrem
künftigen Futter sprechen?“

Der foregaar nu ikke en jevnt fremskridende Opdragelse i Historien fra
primitive Grundværdier til højere (det vil sige: til saadanne, som de, der have
erhvervet dem, kalde højere i Sammenligning med de primitive). Historien er
Grundværdiernes store Kampplads. Her staar Individ mod Individ, Individ og
Samfund mod hinanden, og Samfund staar mod Samfund, og det er ofte først efter
stærke Brydninger, at en ny Grundværdi kan fæstnes i Menneskenes Sind. Et
Exempel herpaa er den Maade, paa hvilken Alexanders og Romernes Erobringer have
banet Vejen for den almindelige Humanitetsfølelses Udvikling. Under denne store
Kamp kan den etiske Tænkning kun indirekte spille en Rolle med: ved at drage de
bestemte Konsekvenser af hvert enkelt Stade og ved at belyse disse
Konsekvensers Betydning ved saa mange Erfaringer som muligt, --- altsaa ved at
fremme Selverkendelsen i Sokrates' Aand. I selve Kampen eller Opdragelsen
gælder det om en Kunst; Videnskaben slaar ikke til, saa betydningsfuldt et
Bidrag den end ofte kan give.

Men Kontinuitetsprincipet forlade vi ikke, fordi vi staa ved
% --- p. 77 ---
\opage{77}den videnskabelige Etiks Grænse. Dér, hvor det ikke er muligt at
forfølge Kontinuiteten længere i rent etisk Form, søge vi at efterspore den i
psykologisk og historisk Form. Etiken gaar paa dette Punkt over i Psykologien
og i Sociologien.

Medens jeg i min Etik har søgt at vise, at Retfærdigheden, opfattet som et
indre harmonisk Forhold mellem Selvhævdelse og Hengivelse og tillige mellem
Følelse og Tænkning, er den højeste Karakteregenskab, er dette blevet bestridt
af filosofiske Tænkere, der hævde, at der i Selvhævdelsen og Hengivelsen
kundgøre sig saa forskellige Tendenser (eller, som jeg vilde sige:
Grundværdier), at det ikke er muligt at forene dem i et omspændende
Begreb\textsuperscript{57}). Og i den populære Literatur er Hengivelsen
(Altruismen) gjort gældende ligesaa hensynsløst af \textsc{Leo Tolstoy} som
Selvhævdelsen af \textsc{Friedrich Nietzsche}. Her have vi et Exempel paa en
Strid mellem Grundværdier, der samtidigt trænge sig frem. Der vil sandsynligvis
stedse i det virkelige Menneskeliv foregaa Svingninger mellem disse to modsatte
Poler. I Slægten som Helhed have baade Selvhævdelsen og Hengivelsen deres
Funktion at udøve, Men Vurderingen af den enkelte Svingning i den ene eller den
anden Retning vil stedse bero paa, om den fører fremad i Retning af en
Livsordning, indenfor hvilken hver enkelt Personlighed kan udvikle sig saa
ejendommeligt og selvstændigt som muligt, saaledes at den netop derigennem yder
den bedste Hjælp til Andres lige ejendommelige og selvstændige Udvikling. Hos
de enkelte Individer ville Selvhævdelsen og Hengivelsen staa i højst
forskelligt Forhold til hinanden, og den Harmoni mellem begge, hvori efter min
Opfattelse Retfærdigheden bestaar, vil derfor fremtræde med højst forskellig
Klangfarve hos de Forskellige og til de forskellige Tider. Men ogsaa det hører
med til Livets Rigdom --- og til dets Kontinuitet. Nuancernes Mangfoldighed
betinger Sammenhængen og gør den forskellig fra død Ensformighed, der hverken
fra et etisk Synspunkt eller fra noget andet Synspunkt har nogen Værdi.

4. En anden Vanskelighed, der hænger sammen med Diskontinuiteten, opstaar
derved, at de Slutninger, der under visse Forhold drages af en Grundværdi,
skulle anvendes paa \textit{Individer,
% --- p. 78 ---
\opage{78}der møde med højst forskellige indre og ydre Betingelser} for at
fyldestgøre den Fordring, en saadan Slutning begrunder. Anlæg og Drifter ere
ikke de samme hos alle Individer, hverken i Art eller i Grad. En og samme
Fordring rettet til forskellige Individer paalægger dem altsaa et højst
forskelligt etisk Arbejde. Begyndelsespunktet og Begyndelseshastigheden ere
forskellige. Nogle Individer ere maaske allerede paa Vej til uvilkaarlig
Udførelse af Fordringens Indhold, før Fordringen bliver dem bevidst; andre maa
møjsommelig kæmpe for at lægge den første Grund. Loven, Fordringen maa derfor
stilles forskelligt for de Forskellige, hvis den virkeligt skal blive den samme
for Alle. Enhver maa beskattes efter Evne. Der maa foretages en gennemgaaende
Individualisering af den etiske Fordring, for at ikke selve Etiken skal
forsynde sig mod den Sætning, at Personligheden stedse er Formaal, aldrig blot
Middel. Den etiske Fordring skal ikke være et abstrakt eller udvortes Bud, men
skal svare til de etiske Muligheder i den enkelte Personlighed og være i Stand
til at udvikle dem videre. Lovgivning og Pædagogik kunne her ikke holdes
absolut ude fra hinanden. Men det gør etiske Afgørelser vanskelige i de enkelte
Tilfælde\textsuperscript{58}). Her er igen Verden --- Personlighedernes Verden
--- stor, og vor Hjerne lille, Den etiske Tænkning kan ikke formulere en Lov,
der uden videre kan anvendes i Livets Mangfoldighed. Og dog maa vi antage, at i
hvert enkelt Tilfælde kan kun en enkelt Afgørelse være den rette (Eng ist die
Welt, und das Gehirn ist weit). Vore etiske Tanker kredse om den snevre Vej,
der fører til den etiske Sandhed, og arbeider sig fremad til den gennem stadig
Kamp med det Irrationelle, som de søge at overvinde ved at finde saa stor
Tilnærmelse som muligt. Her som indenfor Erkendelsesproblemet ende vi ikke med
en absolut Afslutning, men høre et Excelsior! --- selv hvor vor Tanke har
anstrengt sig til det Yderste.

\begin{center}\textbf{\large B.\ \ Det religiøse Problem.}\end{center}

5. Allerede selve det, at der existerer et religiøst Problem, viser
Kontinuitetens Betydning. Ti at Religionen er bleven
% --- p. 79 ---
\opage{79}problematisk, hænger sammen med, at der er indtraadt en
Arbejdsdeling, en Differentiation paa det aandelige Livs Omraade. I Religionens
klassiske Tider staar den som den ene, koncentrerede Form, i hvilken alle
Aandslivets Fornødenheder tilfredsstilles; Religionen er som saadan ikke blot
Hjælp og Trøst, men er Poesi, Moral og Videnskab, eller den formaar i hvert
Tilfælde at optage disse Omraader i sig paa organisk Maade. Efterat disse
forskellige Omraader have emanciperet sig, og hvert for sig udvikler sig efter
sine egne Love, bliver der Spørgsmaal om, hvorvidt Aandslivet --- hvad Værdi
angaar --- har bevaret sin Kontinuitet under denne Overgang fra Koncentration
til Differentiation. At der er en psykologisk og historisk Kontinuitet til
Stede, er der ingen Grund til at tvivle om; men er der ikke for det
menneskelige Liv gaaet Værdi tabt ved hin Overgang?

Svarene paa dette Spørgsmaal lyde højst forskelligt. Nogle lade sig nøje med,
at de samme Dogmer læres endnu som i gamle Dage. Derimod hævdes fra anden Side,
at denne dogmatiske Kontinuitet Intet betyder: det Væsentlige er, om Livet
føres paa samme Maade, --- om der er etisk Kontinuitet, --- om „det nye
Testamentes Kristendom“ endnu er til. Og endeligt er der dem, der mene, at med
den kulturhistoriske Arbejdsdeling er Religionens Tid forbi, og at dette intet
Værditab betyder, men tvertimod er en Vinding for Aandslivet. Foruden disse
Hovedstandpunkter, der selv kunne fremtræde i forskellige Nuancer, gives der en
Række forskellige mindre karakteristiske Standpunkter.

Spørgsmaalets Besvarelse afhænger naturligvis af, hvad man anser for det
Væsentlige i Religionen. Kontinuiteten kan ikke fremtræde i saadanne Træk, der
skille de forskellige Religioner og religiøse Standpunkter fra hverandre; man
maa opsøge de dybest bevægende Tendenser, der kunne aabenbare sig under højst
forskellige Former, og som maaske ogsaa kunne vedblive at virke, efterat den
kulturhistoriske Arbejdsdeling er traadt i Kraft. Dette er Religionsfilosofiens
Opgave, som den især søger at løse ad to Veje. For det første anstiller den en
Sammenligning mellem Religionen og andre Sider af Aandslivet for at
% --- p. 80 ---
\opage{80}finde, i hvilken Egn af Aandslivet Religionen hører hjemme; den søger
at bestemme Religionens psykologiske Sted. Og for det Andet anstiller den en
Sammenligning mellem de vigtigste historiske Former for Religion for at se, om
den fundne psykologiske Bestemmelse vinder Bekræftelse i
Erfaringen\textsuperscript{59}).

Ved begge Undersøgelser ville vi faa Lejlighed til at se, at det ikke blot er
Kontinuitetsprincipet, der fører til at stille det religiøse Problem, men at
ogsaa selve Begrebet Religion hænger meget nøje sammen med dette Princip.

6. Den religiøse Følelse kan stilles i et klart og simpelt Forhold til de andre
Følelser, naar den opfattes som bestemt ved de Erfaringer, Mennesket gør
angaaende Værdiernes Skæbne i Virkelighedens Verden. Derved ses den paa en Gang
i sin Forskel fra og i sin Sammenhæng med de andre Følelser. Til enhver Følelse
svarer en Værdi. Livsfølelsen, den intellektuelle, den æstetiske og den etiske
Følelse kundgøre saaledes forskellige Arter af Værdi. Livets Opretholdelse og
Udvikling, Sandheden, Skønheden og Godheden ere Værdier, Mennesket kan have Del
i uden religiøs Følelse. Men naar det viser sig, at Liv, Sandhed, Skønhed og
Godhed maa kæmpe for at bestaa i Verden, opstaar der en ejendommelig Følelse,
som ikke længere bestemmes ved selve hine Værdier i og for sig, men ved deres
Bestaaen, Fremgang eller Tilbagegang i Tilværelsen, saavidt Mennesket formaar
at overskue denne. De Erfaringer, Mennesket gør herom, ville kunne afføde en
Trang til Tro paa, at Værdierne bestaa, selv hvor de i den Virkelighed,
Mennesket kan overskue, ikke længere kundgøre sig, eller selv om de ikke
længere fremtræde i de samme Former som hidtil. Mest tilbøjelig vil Mennesket
være til at antage, at den Værdi, der bestaar, er den samme, han hidtil har
været delagtig i; men han vil kunne udvikle sig til at fastholde, at de
empiriske Værdier maa gaa under, for at en højere og mere omfattende
Værdisammenhæng kan ske Fyldest. Af Menneskenes Forestillinger om en Gudeverden
og om et hinsidigt Liv lærer man at kende, hvilke Værdier der har ligget deres
Hjerter nærmest, og i hvilken Grad de have forsonet sig med den Tanke, at de
empiriske Værdier maa undergaa en mere eller mindre væsentlig Metamorfose, for
% --- p. 81 ---
\opage{81}at en Værdiens Bestaaen skal være mulig. Tilsidst bliver det, der
bestaar, Noget, som intet Øje har set og intet Øre har hørt. Den religiøse Tro
hævder en Kontinuitet paa Værdiernes Omraade ud over enhver mulig Erfaring.
Denne Værdikontinuitet kan være af forskelligt Indhold og Omfang. Men der vil
være en Tendens til at udstrække den ud over den menneskelige Verden og opfatte
Tilværelsen som et Hjemsted for Værdiers Udvikling og Bestaaen. Denne religiøse
Kontinuitetstro danner en Analogi til den intellektuelle Forudsætning om
Tilværelsens Rationalitet (III), og de Vanskeligheder, den har at kæmpe med,
ere ogsaa analoge.

Den religiøse Tro vil særligt have en vis Lighed med (og en vis Sympati med)
den metafysiske Idealisme (IV, 4). Men den staar og falder ikke med denne. Det
Væsentlige i den religiøse Tro er ikke den intellektuelle Tilfredsstillelse,
den kan yde. Den Kreds af Forestillinger, i hvilken Religionen udtrykkes, hører
ikke til dens inderste psykologiske Væsen. Det Centrale i Religionen er ---
hvis den givne psykologiske Bestemmelse er rigtig --- en Følelses- og
Villiesinteresse. Intellektuelt spørge vi kun om Genkendelse (Identitet),
Konsekvens (Rationalitet) og Aarsagssammenhæng (Kausalitet) (III, 1). Troen er
stedse kun \textit{Genstand} for Videnskab, ikke selv Videnskab. Den opstaar
ved det harmoniske eller disharmoniske Forhold mellem det Billede af
Virkeligheden, Erkendelsen giver Mennesket, og de Værdier, der for Mennesket
staa som de højeste.

Muligvis gælder den Forskel, vi ere nødte til at gøre mellem Værdi og
Virkelighed, kun fra et menneskeligt Standpunkt. Men --- andet Standpunkt kunne
vi nu engang ikke indtage.

Troen paa Værdiens Bestaaen kan selv faa Værdi ved at holde Modet oppe under
Kampen for Livet og ved at anspore til at finde nye Værdier som Ækvivalenter
for de Værdier, der forsvinde. Selv den, der antager, at Religionens Tid er
forbi, --- en Antagelse, der, naar den skal være mere end en Paastand eller et
Ønske, maa begrundes erkendelsesteoretisk, psykologisk og etisk, --- selv han
vil dog sikkert føle Trang til at finde Ækvivalenter for det Værditab, der
lides ved Religionens Forsvinden. I denne Forstand existerer der et religiøst Problem
% --- p. 82 ---
\opage{82}ogsaa for den, der mener, at Religionens Kærne forsvinder med dens
Skal. Religionen faar nemlig ved at hævde Værdiens Bestaaen selv Værdi, Værdi i
anden Potens. Og den betegner særligt en af de mest koncentrerede Former for
Aandsliv, Erfaringen frembyder: den staar, hvor den er ægte og oprindelig, som
det samlede Resultat af Menneskets Følelses- og Villiesliv, og den skaffer sig
Udtryk i de vældigste og ædleste Billeder, Menneskets Forestillingsliv raader
over. Sandsynligvis vil denne koncentrerede Livsproces stedse ytre sig under
nye Former --- hvad enten man saa vil vedblive at bruge Ordet Religion eller ikke.

7. En Bekræftelse af den fremsatte Karakteristik af det Væsentlige i Religionen
kan søges deri, at Forskellighederne mellem Religionerne eller de religiøse
Standpunkter indbyrdes simpelt og let kunne forklares ud fra den. Disse
Forskelligheder bero nemlig dels paa Forskelligheder i de Værdier, paa hvis
Bestaaen der tros, dels paa Forskelligheder i den Opfattelse af Virkeligheden,
der ligger til Grund, dels paa Forskelligheder i Erfaringen angaaende Forholdet
mellem Værdierne og Virkeligheden.

Forskellene med Hensyn til \textit{Værdierne} fremtræde i de historiske
Religioner i det, man betragter som Gudernes Opgave, --- det, man tænker sig
Guderne kæmpe for, --- eller, for at bruge et Udtryk af \textsc{Platon}, i
\textit{det, der gør Gud guddommelig} \textit{\textsuperscript{60})}. Det kan
være den rent fysiske Livsopholdelse og de til den knyttede Nydelser, eller det
kan være det Gode, Skønne og Sande. Religionshistorikerne fremhæve Overgangen
fra Naturreligion til etisk Religion som den betydningsfuldeste indenfor
Religionshistorien; men denne Overgang beror netop paa en Overgang fra de mere
elementære Værdier til de ideelle Værdier som det, der ligger til Grund ved
Forholdet til Guderne, --- det, Guderne værne om, og som de selv skylde deres
Guddommelighed. Tillige falder det mere udvortes Forhold mellem de Magter,
Mennesket tror sig i Forhold til, og det, disse Magter værne om, bort, og
Guderne selv blive mere umiddelbare Repræsentanter for det Værdifulde, de
hævde, ja blive Et dermed. En stedse inderligere Sammenhæng mellem det
Religiøse og det Etiske kan spores i Religionens Udvikling, og Religionen
tenderer mere og mere til at blive en Projektion af det Etiske.

% --- p. 83 ---
\opage{83}Naar \textit{Virkelighedsbilledet} ændres ved fremadskridende
Erkendelse, sker der en ikke mindre afgørende Ændring i Religionens Karakter,
end naar de til Grund liggende Værdier ændres. De mest iøjnefaldende religiøse
Kriser skyldes ny Virkelighedsopfattelse. Religionen har her gennemgaaende
forholdt sig receptivt til den ad anden Vej udviklede Erkendelse, men har
bagefter søgt at optage den og bruge den som Grundlag for sine Billeder.
Saaledes have successivt Animismen, Stjernekundskaben, Sjælevandringslæren,
indisk og græsk Filosofi, Kopernikanismen og nyere Naturvidenskab, spekulativ
og kritisk Filosofi grebet ind, snart som velkomne Former for religiøst
Forestillingsliv, snart som udvidende Impulser, snart som fjendtlige Tendenser.

Men det Afgørende vil stedse være \textit{Forholdet mellem Værdierne og
Virkeligheden.} Alt efter som dette Forhold fremtræder i Menneskets
Livserfaring som harmonisk eller disharmonisk, vil Religionen faa en forskellig
Karakter, og i de højeste Religioner ville baade Livets Disharmonier og dets
Harmonier finde deres Plads; den religiøse Stemning vil dér først naa Højde og
Fasthed, naar den gennem Brydning og Lidelse har arbejdet sig frem til
Saligheden. Disharmoniernes Styrke bliver da kun et Vidnesbyrd for den sejrende
Harmonis Styrke. Men dertil kommer et andet Forhold. Religionen vil faa en
forskellig Karakter, alt efter som det højeste Værdifulde tænkes i evig
Bestaaen, hævet over al Vorden og Forandring, saa at Livet i Tiden tilsidst er
en Illusion, eller om det Værdifulde selv udfolder sig i Tiden og kæmper for
sin Bestaaen under Forandringernes Forløb. Paa denne Forskel beror Modsætningen
mellem den indisk-græske Religionstype og den persisk-kristelige Religionstype.

Det synes saaledes ikke at være nogen fremmed Maalestok, vi anlægge paa
Religionerne, naar vi karakterisere og vurdere dem efter den Maade, hvorpaa, og
den Grad, hvori Principet om Værdiens Bestaaen fremtræder i dem. Det er en
Maalestok, der naturligt frembyder sig for en sammenlignende Betragtning af
Religionernes og de religiøse Standpunkters Historie.

Den anden af de nævnte Religionstyper staar i Historien
% --- p. 84 ---
\opage{84}som den sejrende. Den har sikkert ikke naaet sin definitive Form.
Hvis den vil komme til at antage nye Former, vil det som i Fortiden skyldes
profetiske Personligheder. Kun de formaa at væve Guddommen et nyt Klædebon.

For Filosofen er den væsentlige Interesse den at paavise, at saalænge der
bestaar et religiøst Problem, hænger dette nøje sammen dels med
Kontinuitetstrangen, dels med Tidsforholdet, dels med de etiske Værdier.
Slægtskabet og den indre Sammenhæng mellem Menneskeaandens store Problemer, de
praktiske som de teoretiske, træder her tydeligt frem.

Og ved alle Problemerne ende vi --- naar vi se Modsætningen i dens skarpeste
Form --- i den evige Strid, en Strid, i hvilken Menneskenes aandelige Kræfter
stedse paany skulle prøves. Men fordi vi ikke kunne løse de store Problemer,
kunne vi dog se Vejen, der fører fremad saaledes, at baade Tankens og Livets
Ret hævdes. Problemernes Uløselighed betyder egentligt kun, at hvor langt vi
end naa i vor Forsken og Tænken, danner der sig stedse nye Horizonter for os,
nye Maal og nye Opgaver. 

\begin{center}\rule{2cm}{0.4pt}\end{center}
\clearpage

% ===== NOTER (pp. 85--90) =====
\begin{small}
% --- p. 85 ---
\opage{85}

\begin{center}{\LARGE Noter.}\\[6pt]\rule{1.5cm}{0.4pt}\end{center}

\textsuperscript{1}) (S. 2). \textsc{Ch. Renouvier}: \textit{Les dilemmes de la
métaphysique pure.} Paris 1901, p. 202.

\textsuperscript{2}) (S 3). %
% PRINTED AS IS: no point after S
 Disse fire Problemer indeholde den antike, fra Platons Skole (ifølge
\textsc{Sextus} \textsc{Empiricus}: \textit{Adv. mathematicos} VII, 16 fra
Xenokrates) stammende Tredeling i Logik, Fysik [ɔ: Kosmologi] og Etik med
Tilføjelse af det psykologiske Problem, der i den nyere Tid har skudt sig frem
som selvstændigt Udgangspunkt.

\textsuperscript{3}) (S. 5). Smlgn. foruden Renouvier's ovenfor (Note 1) nævnte
Skrift \textsc{Emile Boutroux's} interessante Skrifter: \textit{De la
contingence des lois de la nature.} 2. éd. Paris 1895 og \textit{De l'idée de
loi naturelle dans la science et la philosophie contemporaines.} Paris 1895.
--- Beslægtet er \textsc{Rudolph Eucken's} Standpunkt (\textit{Der Kampf um
einen geistigen Lebensinhalt.} Leipzig 1896).

\textsuperscript{4}) (S. 6). \textsc{Ph.} \textsc{Lipps}; \textit{Psychologie,
Wissenschaft und Leben.} München 1901. Om Fries og Beneke se mit Værk
\textit{Den nyere Filosofis Historie.} II. p. 221 f; 236; 240.

\textsuperscript{5}) (S 8). %
% PRINTED AS IS: no point after S
 Se min \textit{Psykologi} V A, 2.

\textsuperscript{6}) (S. 10). Smlgn. min Afhandling \textit{Det psykologiske
Grundlag for logiske Domme} (Vid. Selsk. Skr. 6. Række). p. 16 f; 59, og min
Afhandling \textit{Om Vitalisme} (optrykt i „Mindre Arbejder“).

\textsuperscript{7}) (S. 11), Af Interesse for den moderne Positivismes
Standpunkt overfor dette Problem i Modsætning til det Standpunkt, Comte og
Stuart Mill indtog, er \textsc{Roberto Ardigo's} Skrift \textit{L'unità della
coscienza.} Padova 1898.

\textsuperscript{8}) (S. 11). Smlgn. min \textit{Psykologi} I, 8 d.

\textsuperscript{9}) (S. 11). \textsc{Münstehrerg}: \textit{Psychology and
Life.} Westminster 1899. --- Münsterberg slutter sig til den Betragtning, som
Windelband og Rickert have gjort gældende, ifølge hvilken der er en Kløft
mellem Naturvidenskab (til hvilken ogsaa Psykologien regnes) og
Kulturvidenskab. Se om denne Betragtning de træffende Bemærkninger hos
\textsc{Wundt;} \textit{Einleitung in die Philosophie.} Leipzig 1901. p.
65---74, og hos \textsc{Guido Villa}: \textit{Psychology and History.} (The
Monist. January 1902.)

\textsuperscript{10}) (S. 12) \textsc{Münsterberg} anf. Skr. p. 282: „If we
force the system of science upon the real life, claiming that our life is
really a psychophysical
% --- p. 86 ---
\opage{86}phenomenon, we are under the illusion of psychologism. If, on the
other hand, we force the views of the real life, the personal categories, upon
the scientific psychophysical phenomena, we are under the illusion of
mysticism. The result on both are the same. We lose the truth of life and the
truth of science.“

\textsuperscript{11}) (S. 13). \textsc{Descartes} (\textit{Synopsis
meditationum}) lærer, at alle Substanser ere uforgængelige, men at det kun er
corpus in genere sumptum, der er Substans, ikke f. Ex. corpus humanum; enhver
Sjæl derimod er pura substantia, som ikke forandres, hvormeget end dens enkelte
accidentia (Tanker, Følelser og Villiesytringer) skifte. --- I sit Skrift
\textit{Cogitata metaphysica} (II c. 11) siger \textsc{Spinoza} (ud fra det
kartesianske Standpunkt, hvorpaa han i dette Skrift væsentligt stiller sig): Si
ad totam Naturam materiæ attendamus, illi nihil novi accedit; at respectu rerum
particularium aliqvo modo potest dici, illi aliquid novi accedere. Qvod an
etiam locum habet in rebus spiritualibus, non videtur: nam illa [ɔ:
spiritualia] ab invicem ita dependere non apparet [Sjælene afhænge ikke
saaledes af hverandre, som Legemerne afhænge af hverandre]. Senere kom Spinoza
til en anden Opfattelse.

\textsuperscript{12}) (S. 13). Støttet til sin Fortolkning af Aristoteles'
„aktive Fornuft“ lærer \textsc{Averroes} (se om ham \textsc{Renan}:
\textit{Averroës et l'Averroïsme.} Paris 1852), at medens de enkelte Sjæle
opstaa og forgaa, bestaar intellectus universalis, den Verdenstanke, der virker
i de enkelte Sjæles Tænkning. I Middelalderen fornyedes denne Lære bl. A. af
\textsc{Siger Af Brabant} i hans nyligt af Mandonnet udgivne \textit{Qvæstiones
de anima intellectiva.} (Mandonnet: Siger de Brabant et l'Averroïsme latin au
13' siècle. Fribourg 1899). Den hele Lære om Aandens aktuelle Bestaaen gaar
tilbage til Neoplatonismen. (Se \textsc{Plotinos}: \textit{Ennead.} V, 9, 6).
--- Om \textsc{Spinozas} Lære om intellectus infinitus og idea dei se
\textit{Den nyere Filosofis Historie.} I. p. 290; om \textsc{Hegel} ib. II, p 165.

% PRINTED AS IS: the Noter label this note „12)“; its anchor in the text (p. 14) is 12b
\textsuperscript{12}) (S. 14). \textsc{Maxwell}: \textit{Scientific papers.}
Vol. II. Cambridge 1890. p. 759. Kort forud for den citerede Ytring bemærker
Maxwell, at medens man nu til Dags i Reglen har opgivet den Tanke, at Sjælen
skulde kunne findes et Sted i Hjernen ad anatomisk Vej, har den Tanke holdt sig
længere, at man ved at forfølge de materielle Processer tilstrækkeligt langt
tilbage vil komme til en materiel Proces, der er bevirket af Sjælen. Det er mod
denne Mulighed, den citerede Ytring er rettet. Smlgn. \textsc{Spinoza}:
\textit{Ethica} III, 2. Schol. --- Til et lignende Resultat som Maxwells
Opfattelse af Energisætningen fører ogsaa hans Opfattelse af Inertisætningen i
hans Skrift \textit{Matter and Motion} § 41. --- Smlgn. min
\textit{Psykologi}\textsuperscript{4} p. 36 og mine \textit{Psykologiske
Undersøgelser.} p. 82---86.

\textsuperscript{13}) (S. 14). \textsc{Carl Lange}: \textit{Nydelsernes
Fysiologi.} København 1899 p. 45.

\textsuperscript{14}) (S. 15). \textsc{Münsterberg}: \textit{Psychology and
Life.} p. 127: Mental facts, as they are not quantitative, cannot enter into
any causal equation.

\textsuperscript{15}) (S. 16). \textsc{Flechsig}: \textit{Gehirn und Seele} 2.
Leipzig 1896. p. 24 f.

\textsuperscript{16}) (S. 17). Smlgn. Kritiken af Flechsigs Teori hos
\textsc{O. Vogt} (\textit{Flechsigs Associationscentrenlehve, ihre Anhänger und
Gegner.} Zeitschrift für Hypnotismus. V, 6) og \textsc{Alb. Adamkiewicz}
(\textit{Die Grosshirnsrinde als Organ der Seele.} Wiesbaden 1902. p. 75 f.).

\textsuperscript{17}) (S. 18). \textsc{Münsterberg}: \textit{Psychology and
Life.} p. 172: The reality of the will and feeling and judgment do not belong
to te describable world. --- p. 208; The subjective attitude is never object;
it is never perceived.

\textsuperscript{18}) (S. 20). \textit{Den nyere Filosofis Historie.} I p. 334---336.

\textsuperscript{19}) (S. 21). Smlgn. min \textit{Psykologi} V B, 6.

\textsuperscript{20}) (S. 24). Se den interessante Diskussion om \textit{Le
parallélisme psycho%
% --- p. 87 ---
\opage{87}physique et la Métaphysique positive} i Bulletin de la société
française de Philosophie. Juin 1901. Bergson bemærkede under denne Diskussion
bl. A. (p. 51).: Étant donné un état psychologique, la partie jouable de cet
état, celle qui se traduirait par une attitude du corps ou par des actions du
corps, est représentée dans le cerveau: le reste en est indépendant et n'a pas
d'équivalent cérébral. De sorte qu'à un même état cérébral donné peuvent
correspondre bien des états psychologiques différents, mais non pas des état
quelconques. Ce sont des états psychologiques qui ont tous en commun le même
„schéma moteur“.

\textsuperscript{21}) (S. 25). Et Exempel herpaa har jeg anført i min
\textit{Psykologi.} 4. Udg p. 67 Note.

\textsuperscript{22}) (S. 26). Smlgn. min \textit{Psykologi} VII B, 4.
\textit{Psykologiske Undersøgelser.} p. 67---81. --- \textit{Søren Kierkegaard
som Filosof.} p. 74---81. --- \textsc{Ebbinghaus}: \textit{Grundzüge der
Psychologie.} I p. 168 (smlgn. ogsaa hans Artikel i Zeitschrift für Psychologie
IX p. 201) og \textsc{Ehrenfells}: \textit{Werttheorie.} I p. 245---249, slutte
af Ikkeiagttageligheden af Villien, at Villien ikke er et særligt psykisk
Element. Smlgn. min Kritik af sidstnævnte Skrift i Göttinger gel. Anzeigen.
1900 p. 742 f.

\textsuperscript{23}) (S. 26). Se nærmere herom i min \textit{Psykologi} II, 5;
IV, 7 e; V B, 5; VII A, 1; B, 4---5.

\textsuperscript{24}) (S. 27). \textsc{Ostwald} søger i sin
\textit{Naturphilosophie} (Leipzig 1902) at vise, at Begrebet Energi er det
fundamentale naturvidenskabelige Begreb, og da nu ogsaa Bevidsthedsfænomenerne
(med Kant) maa opfattes som Virksomhedsytringer (Bethätigungen), saa bliver det
ogsaa det fundamentale psykologiske Begreb. Men om Forholdet mellem de to Arter
af Energi udtaler han sig ikke med Klarhed. Snart kaldes selve
Bevidshedsprocesserne „energetiske“ (p. 394), snart siges Bevidstheden at være
betinget af Nervesystemets Energi (p. 396), snart defineres aandelig Energi som
„bevidst \textit{eller} ubevidst Nerveenergi“ (p. 398).

\textsuperscript{25}) (S. 28). Se om de forskellige Arter af Genkendelse mine
\textit{Psykologiske Undersøgelser.} p. 35 f.

\textsuperscript{26}) (S. 29). Se \textit{Det psykologiske Grundlag for logiske
Domme.} p. 32 f.

\textsuperscript{27}) (S. 30). Se \textit{Det psykol. Grundlag.} Kap. 6.

\textsuperscript{28}) (S. 32). \textsc{Hobbes}: \textit{Logica.} Cap. 3 (§§
8---9); \textit{Physica.} Cap. 25 (§ 1). (Dog lærer Hobbes et andet Sted, at
det er Analysen ud fra det Givne, der fører os til Principerne: Analytica est
ars ratiocinandi a supposito ad principia, id est ad propositiones primas vel
ex primis demonstratas, qvot sufficiunt ad suppositi veritatem vel falsitatem
demonstrandam. \textit{De rationibus motuum et magnitudinum.} Cap. 20 (§ 6).)
--- \textsc{Fichte}: \textit{Grundlage der gesammten Wissenschaftslehre.} § 1.
--- \textsc{S. Kierkegaard}: \textit{Uvidenskabelig Efterskrift.} p. 83. ---
\textsc{Kroman}: \textit{Vor Naturerkendelse.} p. 270---274.

\textsuperscript{29}) (S. 33). \textsc{Ernst Mach}: \textit{Beiträge zur
Analyse der Empfindungen'.} p. 144.

\textsuperscript{30}) (S. 33). \textsc{Maxwell}: Scientific Papers. II. p. 360.
--- \textsc{H. Hertz}: \textit{Einleitung.} (Werke III. p. 1 f.). --- Smlgn. en
interessant Diskussion \textit{De la valeur des lois physiques} i Bulletin de
la société française de philosophie. Mai 1901.

\textsuperscript{31}) (S. 36). \textit{Kritik der reinen Vernunft.} p. 653.

\textsuperscript{31} b) (S. 38). \textit{Den nyere Filosofis Historie.} II. p.
114 f; 222.

\textsuperscript{31} c) (S. 38). Udtrykket „statisk Sandhedsbegreb“ blev brugt
--- med en noget forskellig Motivering --- af \textsc{Louis Weber} paa den
filosofiske Kongres i Paris 1900.

\textsuperscript{32}) (S. 39). \textit{Den nyere Filosofis Historie.} I. p. 160---167

\textsuperscript{33}) (S. 41). Se Note 30.

\textsuperscript{34}) (
% --- p. 88 ---
\opage{88}S. 41). \textsc{H. Hertz}: \textit{Ueber die Beziehungen zwischen
Licht und Elektrizität.} 1889. p. 29. --- \textsc{Boltzmans's} Foredrag paa det
tyske Naturforskermøde 1900 (oversat i The Monist. January 1901: The recent
development of method in theoretical Physics) og hans Afhandling i
Wienerakademiets Oversigter CV, 8 (oversat i The Monist. Octbr. 1901: On the
necessity of atomic theories in Physics). --- Smlgn. om det hele Spørgsmaal
ogsaa \textsc{C. Christiansen}: \textit{Den elektromagnetiske Lysteori.} (Det
danske Vid. Selsk.s Oversigter. 1889), og samme Forfatter i \textit{Fysisk
Tidsskrift.} I. p. 4---5.

\textsuperscript{35}) (S. 42). \textsc{Maxwell}: \textit{Scientific Papers.}
II. p. 26---33; 302---305 (hvor M. viser, hvad han skylder Faraday og Lord
Kelvin); 326; 777 f.

\textsuperscript{36}) (S. 42). Smlgn. mine \textit{Psykologiske Undersøgelser.}
p. 77---81.

\textsuperscript{37}) (S. 43). Dette Punkt lader \textsc{Ostwald} ikke komme
til sin Ret, naar han vil føre Alt i Naturen tilhage til „Energi“. Baade Masse,
Rumudfyldning, Vægt og kemiske Egenskaber skulle kunne føres tilbage til
Energibegrebet. Rigtigt er, at alle hine Begreber forudsætte Energibegrebet;
men er dette det eneste Fundamentalbegreb i Naturfilosofien? I saa Fald maatte
man kunne udlede de „geometriske“ Egenskaber af de dynamiske, og derpaa gør
Ostwald intet Forsøg. I sit Skrift \textit{Die Ueberwindung des
wissenschaftlichen Materialismus} (1895) definerer han endog Materie som „en
rumligt [sic] ordnet Gruppe af forskellige Energier“ (p. 28). Denne rumlige
Ordning, der her anerkendes som særligt Moment, skydes --- forekommer det mig
--- noget til Side i hans \textit{Naturphilosophie} (1902), hvor „det
energetiske Verdensbillede“ udvikles. Dog er hans Hovedsætning, hvorpaa han
bygger sit System, den, at „enhver Proces uden Undtagelse kan exakt og
udtømmende beskrives eller fremstilles ved, at man angiver, hvilke Energier der
undergaa tidslige eller rumlige [sic] Forandringer“ (p. 152). At Begreberne
Energi og Tid ikke kunne skilles, er selvfølgeligt, da Energi betyder Evne til
at overvinde en Modstand, og enhver Overvinden tager Tid. Ikke saa
selvfølgeligt er Sammenhængen mellem Begrebet Energi og de rumlige
Forandringer. Denne Sammenhæng skriver sig hos Ostwald simpelt hen fra, at han
har abstraheret Begrebet Energi fra rumlige Fænomener, fra Undersøgelser af
Forandringer paa bestemte Steder i Rummet. Kun ved at dette geometriske Moment
skydes til Side i hans Energibegreb, bliver det ham muligt at tro, at dette
Begreb uden videre er det samme for psykiske og for fysiske Fænomener, saa at
Problemet om „Sjæl og Legeme“ skulde være løst ved dets Opstilling. Se ovenfor
Note 24.

\textsuperscript{38}) (S. 43). \textsc{Hume}: \textit{Treatise on human
nature.} I, 3, 6: Your appeal to past experience decides nothing in the present
case, and at the utmost can only prove, that that very object, which produc'd
any other, was \textit{at that very instant} endow'd with such a power.

\textsuperscript{39}) (S. 44). \textit{Kritik der reinen Vernunft}: p. 181: Wir
werden also durch diese Grundsätze [Kausalitetssætningen og de beslægtede
Grundsætninger om „Substansens“ Bestaaen og om Vexelvirkningen mellem alt
Existerende] die Erscheinungen nur \textit{nach einer Analogie mit der
logischen und allgemeinen Einheit der Begriffe} zusammen zu setzen berechtigt
werden. --- p. 180: Eine Analogie der Erfahrung wird also nur eine Regel sein,
nach welcher aus Wahrnehmungen Einheit der Erfahrung entspringen soll, und als
Grundsatz von den Gegenständen (der Erscheinungen) nicht constitutiv, sondern
nur regulativ gelten.

\textsuperscript{40}) (S. 41). \textsc{Francis Bradley}: \textit{Appearance and
Reality.} London 1893. Chap. 4---7; 18. --- \textsc{Bernard Bosanquet}:
\textit{Logic or the morphology of knowledge.} Oxford 1888. Book I. Chap. 6.

\textsuperscript{41}) (
% --- p. 89 ---
\opage{89}S. 45). \textsc{Rich. Avenarius}: \textit{Kritik der reinen
Erfahrung.} II. p. 332---339. --- \textsc{Ernst Mach}: \textit{Zur Analyse der
Empfindungen'.} p. 168: „Die Thatsache der Nichtumkehrbarkeit der Zeit reducirt
sich darauf, dass die Werthänderungen der physikalischen Grössen in einem
\textit{bestimmten Sinne} stattfinden. Von den beiden analytischen
Möglichkeiten ist nur die eine wirklich. Ein metaphysisches Problem brauchen
wir hierin nicht zu sehen.“ --- Smlgn. om disse Theorier \textsc{Heinrich}
\textsc{Grünbaum}: \textit{Zur Kritik der modernen Causalanschanungen} (Archiv
für systematische Philosophie. 1899. p. 392---409).

\textsuperscript{42}) (S 46). %
% PRINTED AS IS: no point after S
 \textsc{Bosanquet} indrømmer da ogsaa dette:\textit{ At any given moment} we
have no choice but to say that the future is conditioned by the past,... effect
by cause. \textit{Logic} I. p. 272.

\textsuperscript{43}) (S. 47). Se \textit{Det psykologiske Grundlag for logiske
Domme.} p. 46---50.

\textsuperscript{44}) (S. 48). \textsc{Jevons}: \textit{Principles of Science.}
2. ed. London 1877. p. 43. (Allerede \textsc{De Morgan} har, som Jevons
bemærker, udtalt, at vi sædvanligvis tænke og argumentere indenfor en begrænset
Verden eller Sfære af Begreber, selv naar det ikke udtrykkeligt erklæres.) ---
Smlgn. min \textit{Formel Logik} 3 § 16 og § 22.

\textsuperscript{45}) (S. 48). \textsc{Francis Bradley} har klart set
Modsigelsen mellem Tiden og „det Absolute“, naar han erklærer: If time is not
unreal, I admit that our Absolute is a delusion (App. and Reality: p. 206).
(The Absolute has no seasons. ib. p. 500). --- Kritiken af og Reaktionen mod en
spekulativ Retning støtter sig derfor ogsaa ofte til Tidens Realitet som et
Hovedargument. Smlgn. \textsc{C. H. Weisse's} og \textsc{S. Kierkegaards}
Forhold til Hegel. %
% UNSURE: a faint point after „Hegel“ (not in the e-book)
(\textit{Den nyere Filosofis Historie.} II. p. 243; 261 f.) I sit skarpsindige
Skrift \textit{Tidsexistensens Apologi. Et stycke relationsteori} (Upsala 1888)
har \textsc{Pontus Wikner} forsøgt at rokke ved det Enten-Eller, der her
frembyder sig. Han vil vise, at den højeste Fuldkommenhed kun kan naas ved
successiv Udfoldelse af Kvaliteter, der som samtidige vilde udelukke hverandre.
Men den Begrænsning, der ligger i selve Successionen, ophæves derved ikke!

\textsuperscript{46}) (S. 50). Se min Afhandling \textit{Kontinuiteten i Kants
filosofiske Udviklingsgang.} p. 26---35.

\textsuperscript{47}) (S. 51). Smlgn. \textit{Det psykologiske Grundlag for
logiske Domme.} p. 39 f; 46 f, og (for religiøse Dommes Vedkommende)
\textit{Religionsfilosofi.} p. 67---71.

\textsuperscript{48}) (S. 51). Se min Afhandling \textit{Filosofien og Livet.}
(Tilskueren 1901). p. 34. --- Om det Imaginære i erkendelsesteoretisk Betydning
smlgn. \textsc{Wundt}: \textit{System der Philosophie.} p. 195---199.

\textsuperscript{49}) (S. 54). Smlgn. herom min \textit{Religionsfilosofi.} p.
57---65. --- I det erkendelsesteoretiske Afsnit af Religionsfilosofien har jeg
i det Hele allerede fremsat min Opfattelse af Tilværelsesproblemet og dets
Behandling.

\textsuperscript{50}) (S. 56). \textsc{Goethe}: \textit{Farbenlehre.} § 124;
175---177. Smlgn. \textit{Gespräche mit Eckermann.} 18. Febr. 1829; 21. Decbr.
1831. --- Goethe anvendte ikke dette Begreb blot paa Farvelærens Omraade, ogsaa
ved Magnetismen og i Botaniken. Se \textit{Sprüche in Prosa.} (Ueber
Naturwissenschaft) og \textit{Gespr. mit Eckermann.} 27. Januar 1830. ---
\textsc{Hegel} var (som det ses af Goethes \textit{Nachträge zur Farbenlehre})
meget begejstret over Begrebet Urfænomen, saaledes som Goethe havde opstillet det.

\textsuperscript{51}) (S. 58). Se mine \textit{Mindre Arbejder.} p. 180---194.
--- Om Schopenhauer og Lange se \textit{Den nyere Filosofis Historie.} II. p.
198 f; 501 f.

\textsuperscript{52}) (S. 64). Se \textit{Den nyere Filosofis Historie.} II. p.
165 f. --- Det er karakteristisk, at de engelske Tænkere, som i den nyeste Tid
have arbejdet med
% --- p. 90 ---
\opage{90}nogle af Hegels Grundtanker, indrømme, at det snarere er Analogi end
Bevisførelse, der ker er Tale om. \textsc{Francis Bradley} anser end ikke
Analogien for forsvarlig; den højeste Realitet kan, efter ham, hverken kaldes
Sjæl eller Legeme, skønt Sjælen mere end Legemet har den Forbindelse af Omfang
og Enhed, den „self-consisting“, som er sand Realitets Kendemærke.
(\textit{Appearance and Reality.} p. 307; 357). --- I sin Afhandling
\textit{Hegel's treatment of the categories of the idea} (Mind 1900. p. 149 ff)
udtaler\textsc{ Mc. Taggart,} at vi kun have et eneste Exempel paa Kategorien
„Aand“, der hos Hegel faar teologisk eller kosmologisk Betydning. --- Naar man
(som \textsc{A. E. Taylor} i Mind 1900 p. 245) skelner mellem to idealistiske
Skoler, en, der staar Leibniz og Lotze nærmest, og en, der staar Hegel nærmest,
saa maa jeg hævde, at det er den første af disse Skoler, der klarest har set
det virkelige filosofiske Grundlag for metafysisk Idealisme. (At det er en
Analogi med Menneskeaanden, der begge Steder ligger til Grund, indrømmer Taylor
egentligt selv, naar han siger, at begge Retninger ere enige om „the main
principle that it is in mind, and nowhere else, that we are face to face with
the central reality of the universe“. ibid.)

\textsuperscript{53}) (S. 66). Smlgn. om denne Type min
\textit{Religionsfilosofi.} p. 50; 121 f; 273 f.

\textsuperscript{54}) (S. 71). I min \textit{Etik} (2. Udg. p. 21) gik jeg ikke
nærmere ind paa Forholdet mellem Begreberne Værdi, Formaal og Norm. Se derimod
\textit{Det psykologiske Grundlag for logiske Domme.} p. 47 (smlgn. p. 65) og
min \textit{Religionsfilosofi.} p. 12.

\textsuperscript{55}) (S. 74). Smlgn. min \textit{Etik}. 2 p. 144---151.

\textsuperscript{56}) (S. 76). Smlgn. min \textit{Psykologi.} VI. B., 1---5.
--- \textit{Etik} 2 XIII., 4.

\textsuperscript{57}) (S. 77). \textsc{Francis Bradley}: \textit{Appearance and
Reality.} p. 414---418. --- \textsc{A. E. Taylor}: \textit{The Problem of
Conduct.} London 1901. p. 201.

\textsuperscript{58}) (S. 78). Smlgn. min \textit{Etik.} 2 p. 70---73.

\textsuperscript{59}) (S. 80). I min \textit{Religionsfilosofi} har jeg --- til
mere omfattende Belysning af Problemet --- foruden den psykologisk-historiske
Undersøgelse ogsaa anstillet en erkendelsesteoretisk og en etisk Undersøgelse.
Her i denne kortere Fremstilling holder jeg mig især til den
psykologisk-historiske Undersøgelse og antyder kun lejlighedsvis ganske kort de
erkendelsesteoretiske og etiske Synspunkter.

\textsuperscript{60}) (S. 82). \textsc{Platon}: \textit{Fajdros.} p. 249 C
(πρὸς οἷσπερ θεὸς ὢν θεῖός ἐστι). (Nøjagtigere Oversættelse: „det, som det er
Guds Guddommelighed at leve i“).
\end{small}

\end{document}
